% !TEX root = ../main.tex

\section{フックス型微分方程式}

\subsection{2階線型常微分方程式}
%#region
\begin{equation}\label{eq:2-ode}
\frac{d^{2}y}{dx^{2}} + p_{1}(x)\frac{dy}{dx} + p_{2}(x)y = 0   
\end{equation}
を考える。$x\in\mathbb{C}$とする。よって解$y$は解析的である。
また、$p_{1}(x),p_{2}(x)$は$x$の有理関数とする。よって\eqref{eq:2-ode}は広くても$\mathbb{P}^{1}(\mathbb{C})$上の微分方程式である。
%#endregion
%#region
\begin{definition}[$x=\infty$における微分方程式]
    \eqref{eq:2-ode}で$x=\frac{1}{u}$とする。$u=0$の近傍で考えたものを\textbf{$x=\infty$における微分方程式}という。
    （\cite[p.49]{book:okamoto} を参照）
\end{definition}
%#endregion
%#region
ある有理関数$q_{1},q_{2}$を用いて
\[
\frac{d^{2}y}{du^{2}} + q_{1}(u)\frac{dy}{du} + q_{2}(u)y = 0
\]
のように書いたとき、\eqref{eq:2-ode}に連鎖律 $\frac{d}{dx} = -u^2 \frac{d}{du}$ を用いたものとの係数比較により、
\[
q_{1}(u) = \frac{2}{u}-\frac{1}{u^{2}}p_{1}\left(\frac{1}{u}\right),\, q_2(u)=\frac{1}{u^{4}}p_{2}\left(\frac{1}{u}\right)
\]
という関係がわかる。
%#endregion
%#region
$p_{1}, p_{2}$ を $x=x_{0}$ で正則とする。このとき、コーシーの存在定理（Cauchy's existence theorem）から、任意の $(y_{0}, y'_{0}) \in \mathbb{C}^{2}$ に対して、初期条件
$$
y(x_{0}) = y_{0}, \quad \frac{dy}{dx}(x_{0}) = y'_{0}
$$
を満たす微分方程式 \eqref{eq:2-ode} の解 $y = \varphi(x; x_{0}, y_{0}, y'_{0})$ がただ一つ存在する。

ここで、$u = x - x_{0}$ とおく。$p_1, p_2$ および解 $\varphi$ は $x=x_0$ で正則であるから、以下のようにテイラー展開できる。
$$
\varphi(x; x_{0}, y_{0}, y'_{0}) = \sum_{n=0}^{\infty} c_{n} u^{n} \quad (c_{0} = y_{0}, \, c_{1} = y'_{0})
$$
$$
p_{1}(x) = \sum_{n=0}^{\infty} a_{n} u^{n}, \quad p_{2}(x) = \sum_{n=0}^{\infty} b_{n} u^{n}
$$
これらを \eqref{eq:2-ode} に代入すると、
$$
\sum_{n=2}^{\infty} c_{n} n(n-1) u^{n-2} + \left( \sum_{n=0}^{\infty} a_{n} u^{n} \right) \left( \sum_{n=1}^{\infty} c_{n} n u^{n-1} \right) + \left( \sum_{n=0}^{\infty} b_{n} u^{n} \right) \left( \sum_{n=0}^{\infty} c_{n} u^{n} \right) = 0
$$
となる。この式の $u$ の各べきの係数を比較することで、$c_n$ を漸化式として具体的に順次決定することができる。

\begin{definition}[解析関数の特異点の厳密な定義]\label{def:singular_point_strict}
    ある領域 $D$ で正則な関数（微分方程式の解の分枝）$f(x)$ について、その境界上の点 $\zeta \in \partial D$ が特異点であるとは、$\zeta$ を中心とするいかなる小さな開円板 $U$ をとっても、$f(x)$ を $D$ から $U$ へ解析接続できないこと、すなわち「共通部分 $D \cap U$ において $f(x)$ と一致する $U$ 上の正則関数が存在しないこと」をいう。
\end{definition}

\begin{Q}[【核心】「限界まで」という曖昧な表現を、「いかなる近傍への解析接続も不可能」と厳密化することで得られる論理的な利点は何か？]
\textbf{A.} 特異点であることの定義が「条件を満たす正則関数が存在しない」という明確な否定命題になるため、今回のように「$\xi$ を含む開円板上で正則なもの（解析接続）が構成できる」ことを示せば、直ちに「特異点の定義を満たさない」と結論づける論理（対偶や背理法）が数学的な隙なく回るようになるからである。
\end{Q}

\begin{Q}[【核心】なぜ特異点を単なる「最初の収束円の外側」ではなく、「解析接続の限界（境界）」として定義する必要があるのか？]
\textbf{A.} 単に一つの収束円の外側というだけでは、別の点を中心に展開し直す（解析接続する）ことでその先の点でも正則になり得るため、「いかなる局所的な展開を用いても絶対に到達・通過できない行き止まり」を特異点と定義して初めて、解の真の存在範囲と性質を議論できるからである。
\end{Q}

また、得られた級数 $\sum_{n=0}^{\infty} c_{n} u^{n}$ の収束半径を $r$ とすると、命題\ref{prop:singular_point_on_convergence_circle}により、収束円周 $|x-x_0|=r$ 上には解 $\varphi$ の特異点が少なくとも1つ存在する。

\begin{Q}[【補足】解の冪級数の収束半径 $r$ は、具体的にどのくらい大きくなることが保証されているか？]
\textbf{A.} 初期値 $x_0$ から、「係数関数 $p_1(x), p_2(x)$ の特異点」までの距離の最小値以上になることが保証されている。線形常微分方程式においては、解の特異点は係数の特異点以外には決して新たに生じない（特異点が固定されている）という強力な性質があるためである。
\end{Q}

\begin{Q}[【補足】「ある展開中心での冪級数の収束円上の特異点」は、大域的な関数にとっての「真の特異点」と常に一致するのか？]
\textbf{A.} 常に一致するとは限らない。収束円上の特異点はあくまで「その点を中心としたテイラー展開が破綻する限界点」に過ぎず、別の点を経由して解析接続すれば正則に振る舞うこともある。しかし、線形常微分方程式においては前述の通り特異点が係数によって固定されるため、収束限界点は「係数の特異点」あるいは「多価関数の分岐点」といった真の特異点構造に由来するものと一致する。
\end{Q}
%#endregion
%#region
\begin{proposition}\label{prop:linear_ode_singularities_fixed}
    $x=\xi$ を $y=\varphi(x;x_{0};y_{0},y_{0}')$ の特異点とすれば、$\xi$ は $p_{1}(x)$ または $p_{2}(x)$ の特異点である。
    （\cite[p.50 命題2.1]{book:okamoto} を参照）
\end{proposition}

$n$階でも成り立つので以下$n$階で示す。

\begin{proof}
    対偶を示す。すなわち「$x=\xi$ が係数 $p_{1}(x),\dots,p_{n}(x)$ の正則点であれば、解 $y=\varphi(x;x_{0};y_{0},\dots,y_{0}^{(n-1)})$ は $x=\xi$ において正則である（特異点を持たない）」ことを示す。

    $\xi$ に向かう解析接続の経路上で、$\xi$ に十分近い点 $x=x_{1}$ をとる。具体的には、$x_{1}$ から「係数 $p_1(x),\dots,p_n(x)$ の最も近い特異点」までの距離を $R$ としたとき、$|x_1 - \xi| < R$ となるように $x_1$ を選ぶ（$\xi$ は係数の正則点なので、$\xi$ に十分近づけば必ずこのような $x_1$ が取れる）。

    系\ref{cor:linear_maximal_disk}より、$x_1$ での解の冪級数展開は、少なくとも半径 $R$ の開円板上で正則となる。
    $|x_1 - \xi| < R$ であるから、この開円板は $x=\xi$ を完全に含んでいる。
    したがって、解 $y=\varphi(x;x_{0};y_{0},\dots,y_{0}^{(n-1)})$ は $x=\xi$ を含む開円板上まで解析接続可能となり、$x=\xi$ で正則となる。
    ゆえに定義\ref{def:singular_point_strict}より、$x=\xi$ は解の特異点とはなり得ない。
\end{proof}

\begin{Q}[【補足】なぜ線形微分方程式では、解が勝手に（係数にない場所で）特異点を持つことはないのか？]
\textbf{A.} コーシーの存在定理により、係数が正則である限り、解は局所的に必ず冪級数として一意に構成でき、その収束半径は次の「係数の特異点」にぶつかるまで解析接続で伸ばし続けられるからである。
\end{Q}

\begin{Q}[【重要】この命題が、後のパンルヴェ方程式の学習においてなぜ重要になるのか？]
\textbf{A.} 線形方程式では特異点の場所が初期値によらず係数のみで決まる（固定特異点）のに対し、パンルヴェ方程式のような「非線形」方程式では、初期値によって解の特異点の場所が動く（動く特異点）という決定的な違いを対比させるための基礎となるからである。
\end{Q}
%#endregion
%#region
\begin{definition}[線形微分方程式の特異点]\label{def:linear_equation_singularity}
    \textbf{線形微分方程式の特異点}とは、係数関数 $p_{1}(x), p_{2}(x)$ の特異点のことをいう。
    また、$x=\infty$ での特異点は、$x=\frac{1}{u}$ と変数変換したときの $u=0$ における係数の特異点として定義する。
\end{definition}

\begin{Q}[【核心】命題\ref{prop:linear_ode_singularities_fixed}の逆、すなわち「方程式の係数が特異点を持つならば、解も必ずそこで特異点を持つ」は成り立つか？]
\textbf{A.} 成り立たない。以下の例のように、方程式の係数としては特異点であっても、すべての解がそこで正則に振る舞う「見掛けの特異点（apparent singularity）」が存在し得るからである。
\end{Q}

\begin{example}[見掛けの特異点の例]\label{ex:apparent_singularity}
    次の1階線形微分方程式を考える。
    \begin{equation}\label{eq:apparent_ex}
    \frac{dy}{dx}=\frac{y}{x}
    \end{equation}
    この方程式は係数関数が $1/x$ であるため、$x=0$ に特異点をもつ。
    しかし、この微分方程式の一般解は任意の初期値（積分定数）$C \in \mathbb{C}$ に対して $y=Cx$ となり、$x=0$ において正則である。
    したがって、$x=0$ は方程式の特異点であるが、解の特異点とはならない。
\end{example}

\begin{definition}
    このように、線形微分方程式の特異点であって、任意の初期値に対して、解の正則点となる点を\textbf{みかけの特異点}という。
\end{definition}

%#endregion
%#region

$\Xi$を\eqref{eq:2-ode}の特異点の集合とし、$\mathbb{X}=\mathbb{P}^{1}(\mathbb{C})\setminus\Xi$とする。解の基本系を$\vec{\varphi}(x)$とする。

\begin{definition}
    \[
    W(\vec{\varphi})=
    \left[
    \begin{matrix}
        \varphi_{1}(x) & \varphi_{2}(x) \\
        \varphi_{1}'(x) & \varphi_{2}'(x) \\
    \end{matrix}
    \right]
    \]
    を$\vec{\varphi}(x)$の\textbf{ロンスキー行列}といい、
    \[
    w(\vec{\varphi})=\text{det}(W(\vec{\varphi}))
    \]
    を$\vec{\varphi}(x)$の\textbf{ロンスキアン}という
\end{definition}

\begin{proposition}\label{prop:wronskian-ode}
    \eqref{eq:2-ode}の解の基本系$\vec{\varphi}(x)$のロンスキアンは
    \[
    \frac{dw}{dx}+p_{1}(x)w=0
    \]
    なる$1$階微分方程式を満たす。
    （\cite[p.51 命題2.2]{book:okamoto} を参照）
\end{proposition}

この命題は$n$階でも成り立つ。

\begin{proof}
    定理\ref{thm:derivative-of-determinant}から、
    \[
    \frac{d}{dx}(w(\vec{\varphi})(x))=\sum_{i=1}^{n}w_{i}(\vec{\varphi})(x)
    \]
    である。ここで、$w_{i}(\vec{\varphi})(x)$は$w(\vec{\varphi})(x)$の$i$行目を微分した行列の行列式である。
    ここで$i=1,\cdots,n-1$に対して、$i$行目と$i+1$行目が一致するので、
    \[
    w_{i}(\vec{\varphi})(x)=0
    \]
    よって
    \[
    \frac{d}{dx}(w(\vec{\varphi})(x))=w_{n}(\vec{\varphi})(x)
    \]
    だが、$w_{n}(\vec{\varphi})(x)$の$n$行目について
    \[
    (\varphi_{1}^{(n)}(X),\cdots,\varphi_{n}^{(n)}(x))
    \]
    となっているが、満たす微分方程式によって
    \[
    (-p_{1}(x)\varphi_{1}^{(n-1)}(x)-\cdots-p_{n}(x)\varphi_{1}(x),\cdots,-p_{1}(x)\varphi_{n}^{(n-1)}(x)-\cdots-p_{n}(x)\varphi_{n}(x))
    \]
    とあらわせる。detの多重線形性と合わせると$n-2$階微分以下は0になるので、
    \[
    w_{n}(\vec{\varphi})(x)=-p_{1}(x)w(\vec{\varphi})(x)
    \]
    である。よって、
    \[
    \frac{d}{dx}(w(\vec{\varphi})(x))+p_{1}(x)w(\vec{\varphi})(x)=0
    \]
\end{proof}

\begin{proposition}
$\mathbb{X}$内の$x_{0}$を始点とする2つの閉じた道が互いにホモトープならば、
\[
\Gamma(\gamma)=\Gamma(\gamma ')
\]
が成り立つ
証明は定理\ref{thm:monodromy}を参照
\end{proposition}

\begin{Q}[一価性定理の仮定は満たしてるの？]
\textbf{A.} 係数が正則なら定理\ref{thm:linear_ode_existence_and_continuation}からすべての曲線にそって解析接続可能である。\\
\end{Q}

\begin{Q}[ホモトープって何？]
\textbf{A.} 曲線が連続変形できること。(定義\ref{def:homotope}を参照)\\
\end{Q}

\begin{Q}[解を解析接続したものも解になるの？]
\textbf{A.} なる。(定理\ref{thm:ode_solution_continuation}を参照)\\
\end{Q}
%TODO(基本群等について書いてもいい)

\begin{definition}[モノドロミー行列]
    $\Gamma(\gamma)$を\textbf{$\gamma$に対する$\vec{\varphi}(x)$の回路行列}、あるいは、\textbf{$\gamma$に関する$\vec{\varphi}(x)$のモノドロミー行列}という。
\end{definition}

$\Gamma(\gamma_{2}\cdot\gamma_{1})=\Gamma(\gamma_{2})\Gamma(\gamma_{1})$や、$\Gamma(\gamma^{-1})=\Gamma(\gamma)^{-1}$がわかるので、モノドロミー行列は通常の積に関して群構造をもち、その群を$\mathcal{G}(\vec{\varphi}(x))$とかく。\\
また基本群$\pi_1(x_{0};\mathbb{X})$から$\mathcal{G}(\vec{\varphi}(x))$への準同型
\[
\rho([\gamma])=\Gamma(\gamma)
\]
があり、これは$\pi_1(x_{0};\mathbb{X})$の$\mathbb{C}^{2}$上の表現となっている。

\begin{Q}[$\pi_1(x_{0};\mathbb{X})$の$\mathbb{C}^{2}$上の表現って？]
\textbf{A.} 定義\ref{dfn:representation}を参照。\\
\end{Q}

\begin{definition}
    $\mathcal{G}(\vec{\varphi}(x))$を、\eqref{eq:2-ode}の解の基本形$\vec{\varphi}(x)$の\textbf{モノドロミー群}という。
\end{definition}

ここでもう一つの解の基本系 $\vec{\psi}(x)$ をとっておくと、可逆行列 $C$ を用いて
\[
\vec{\psi}(x) = \vec{\varphi}(x)C
\]
となる。また、解析接続と線型操作は交換できるので、
\[
\vec{\psi}^{\gamma}(x) = \vec{\varphi}^{\gamma}(x)C
\]
よって、$\vec{\psi}(x)$ のモノドロミー行列は、
\begin{equation}\label{eq:monodromy_change_of_basis}
\vec{\psi}^{\gamma}(x) = \vec{\varphi}^{\gamma}(x)C= \vec{\varphi}(x)\Gamma(\gamma)C = \vec{\psi}(x)C^{-1}\Gamma(\gamma)C
\end{equation}
より、$C^{-1}\Gamma(\gamma)C$ となる。
よって、
\[
\mathcal{G}(\vec{\psi}(x))=C^{-1}\mathcal{G}(\vec{\varphi}(x))C
\]
より、表現の同値類、$\hat{\rho}:\pi(x_{0};\mathbb{X})\to GL(\mathbb{C})/\sim$を考えることができる。

\begin{definition}[モノドロミー]\label{def:monodromy}
     $\hat{\rho} : \pi_{1}(x_{0};\mathbb{X})\to GL(\mathbb{C})/\sim ; [\gamma] \mapsto [\Gamma(\gamma)]$ を\textbf{モノドロミー}という。
\end{definition}
%TODO (例を書いてもいい)
%#endregion

\subsection{確定特異点}
%#region
2階微分方程式\eqref{eq:2-ode}を考える。$x=\xi$を\eqref{eq:2-ode}の特異点とする。\eqref{eq:2-ode}の特異点ではない点$x_{0}$を取り、$x_{0}$を始点・終点として、$\xi$の周りを1周する閉曲線$\gamma$をとる。

\begin{Q}[そんな閉曲線$\gamma$は常にとれるの？]
\textbf{A.} 今は$p_{1}, p_{2}$を有理関数としているため特異点は離散的である。したがって、他の特異点を巻き込まないように$\xi$の周りを回る道をとることができる。
\end{Q}

\eqref{eq:2-ode}の解の基本系$\vec{\varphi}(x)$の$\gamma$に沿った解析接続を考える。$u=x-\xi$とする。このとき、解が
\[
\vec{\varphi}(x)=(\varphi_{1}(x),\varphi_{2}(x))=\left(u^{\alpha}\sum_{n=0}^{\infty}a_{n}u^{n},u^{\beta}\sum_{n=0}^{\infty}b_{n}u^{n}\right)
\]
とかけるとする。

\begin{Q}[冪級数部分の収束円盤内に$x_{0}$をとることはできる？]
\textbf{A.} できる。特異点は離散的なので、$\xi$から最も近い他の特異点までの距離より小さい収束半径を持つ円盤を考え、その内部に$x_{0}$を設定すればよい。
\end{Q}

このとき、$\gamma$に沿った解の解析接続$\vec{\varphi}^{\gamma}(x)$は、$u \to ue^{2\pi i}$ の変化（偏角が $2\pi$ 増える）を考慮して
\[
\vec{\varphi}^{\gamma}(x)=\left(u^{\alpha}e^{2\pi i\alpha}\sum_{n=0}^{\infty}a_{n}u^{n},u^{\beta}e^{2\pi i\beta}\sum_{n=0}^{\infty}b_{n}u^{n}\right)
\]
となり、行列表記を用いると
\[
\vec{\varphi}^{\gamma}(x)=\vec{\varphi}(x)
\begin{pmatrix}
    e^{2\pi i\alpha} & 0 \\
    0 & e^{2\pi i\beta} \\
\end{pmatrix}
=
\begin{pmatrix}
    e(\alpha) & 0 \\
    0 & e(\beta) \\
\end{pmatrix}
=:\vec{\varphi}(x)\Gamma(\gamma)
\]
となる。\\
また、任意の解の基本系$\vec{\phi}(x)$について、\eqref{eq:2-ode}は線型であるから、ある可逆行列$C$を用いて$\vec{\phi}(x)=\vec{\varphi}(x)C$とかける。このとき、モノドロミー行列は\eqref{eq:monodromy_change_of_basis}から$C^{-1}\Gamma(\gamma) C$とかける。

以下では、具体的な解の表示を仮定せずに一般論を進める。
\eqref{eq:2-ode}の解の基本系を$\vec{\varphi}(x)$とし、閉曲線$\gamma$に沿った解析接続のモノドロミー行列を$\Gamma(\gamma)$とする。その固有値を $e^{2\pi i \alpha}, e^{2\pi i \beta}$ とおく（ここでは記述の簡略化のため $e(\alpha)=e^{2\pi i \alpha}$ のように書くこととする）。

\begin{Q}[$\alpha, \beta$ の取り方には不定性がある？]
\textbf{A.} ある。任意の整数 $n \in \mathbb{Z}$ に対して $e(\alpha) = e(\alpha+n)$ となるため、$\alpha, \beta$ の値には整数分の不定性が生じる。しかし、ここではある特定の $\alpha, \beta$ を一つ固定したとして議論を進める。
\end{Q}

\begin{Q}[なぜ $\alpha-\beta \in \mathbb{Z}$ かどうかで場合分けするのか？]
\textbf{A.} モノドロミー行列 $\Gamma(\gamma)$ が対角化可能か、それともジョルダン細胞（非対角成分に1が残る形）になるかの分岐点だからである。特性指数の差が整数の場合、解に $\log$（対数項）が現れる可能性がある。
\end{Q}

$\alpha-\beta \in \mathbb{Z}$ の場合と $\alpha-\beta \not\in \mathbb{Z}$ の場合で分けて考える。

\begin{itemize}
    \item[\textbf{[1]}] $\alpha-\beta \notin \mathbb{Z}$ のとき\\
    このとき2つの固有値は異なるので $\Gamma(\gamma)$ は対角化可能である。ある正則行列$P$を用いて
    \[
    P^{-1}\Gamma(\gamma)P=
    \begin{pmatrix}
        e(\alpha) & 0 \\
        0 & e(\beta) \\
    \end{pmatrix}
    \]
    とできる。よって、\eqref{eq:monodromy_change_of_basis}により $\vec{\psi}(x)=\vec{\varphi}(x)P$ なる解の基本系に取り直すことで、モノドロミー行列は最初から対角行列であるとしてよい。
    ここで、行列関数 $\Phi(x)$ を
    \[
    \Phi(x)=
    \begin{pmatrix}
        (x-\xi)^{\alpha} & 0 \\
        0 & (x-\xi)^{\beta} \\
    \end{pmatrix}
    \]
    で定める。$\Phi(x)$ が $\gamma$ に沿って1周すると $\Phi^{\gamma}(x) = \Phi(x) \begin{pmatrix} e(\alpha) & 0 \\ 0 & e(\beta) \end{pmatrix}$ となる。
    ここで $\vec{z}(x)=\vec{\psi}(x)\Phi(x)^{-1}$ とおくと、
    \[
    \vec{z}^{\gamma}(x) = \vec{\psi}^{\gamma}(x)\left(\Phi^{\gamma}(x)\right)^{-1} 
    = \vec{\psi}(x)\begin{pmatrix} e(\alpha) & 0 \\ 0 & e(\beta) \end{pmatrix} \begin{pmatrix} e(-\alpha) & 0 \\ 0 & e(-\beta) \end{pmatrix}\Phi(x)^{-1} 
    = \vec{\psi}(x)\Phi(x)^{-1} = \vec{z}(x)
    \]
    となり、$\vec{z}(x)$ は $x=\xi$ の周りで1価であることがわかる。

    \item[\textbf{[2]}] $\alpha-\beta \in \mathbb{Z}$ のとき\\
    $\Gamma(\gamma)$ が対角化可能であれば、[1]のときと同様にして $\vec{z}(x)$ は1価となる。
    一方、$\Gamma(\gamma)$ が対角化不可能で、ある正則行列$P$を用いて
    \[
    P^{-1}\Gamma(\gamma)P=
    \begin{pmatrix}
        e(\alpha) & 1 \\
        0 & e(\alpha) \\
    \end{pmatrix}
    \]
    となるジョルダン細胞のケースを考える。[1]と同様に解の基本系 $\vec{\psi}(x)$ を取り直すことで、モノドロミー行列がこの形であるとしてよい。このとき、
    \[
    \Phi(x)=
    \begin{pmatrix}
        (x-\xi)^{\alpha} & \frac{1}{2\pi i e(\alpha)}(x-\xi)^{\alpha}\log(x-\xi) \\
        0 & (x-\xi)^{\alpha} \\
    \end{pmatrix}
    \]
    と置く。$\log(x-\xi)$ は1周すると $\log(x-\xi) + 2\pi i$ となるため、
    \[
    \Phi^{\gamma}(x)= \Phi(x)
    \begin{pmatrix}
        e(\alpha) & 1 \\
        0 & e(\alpha) \\
    \end{pmatrix}
    \]
    を満たすことが確認できる。したがって、[1]と同様に $\vec{z}(x)=\vec{\psi}(x)\Phi(x)^{-1}$ と置くことで、逆行列との積が打ち消し合い $\vec{z}^{\gamma}(x) = \vec{z}(x)$ となるため、$\vec{z}(x)$ が1価であることがわかる。
\end{itemize}

以上より、以下の極めて重要な命題が成り立つ。

\begin{proposition}\label{prop:single_valued_z}
    \eqref{eq:2-ode}の任意の解の基本系 $\vec{\varphi}(x)$ に対して、ある適当な2次の行列関数 $\Phi(x)$ を選ぶことで、
    \[
    \vec{z}(x)=\vec{\varphi}(x)\Phi(x)^{-1} \quad (\text{すなわち } \vec{\varphi}(x) = \vec{z}(x)\Phi(x))
    \]
    で定められる $\vec{z}(x)$ が、$x=\xi$ の周りで1価の解析関数になるようにできる。\\
    また、$\Gamma(\gamma)$が対角化可能であるという条件のもとで、上の$\vec{\psi}(x)$に対しては$\Phi(x)$を対角行列になるようにとることができる。
    （\cite[p.59 命題2.5]{book:okamoto} を参照）
\end{proposition}

命題\ref{prop:single_valued_z}の$\vec{z}(x)$は定理\ref{thm:laurent_expansion}より、$x=\xi$の周りで、
\begin{equation}\label{eq:power_series_z}
\vec{z}(x)=\sum_{n=-\infty}^{\infty}\vec{a}_{n}u^{n}
\end{equation}
とローラン級数に展開される。


\begin{definition}[微分方程式の確定特異点]
    \eqref{eq:power_series_z}において、ある整数$j_{0}$があって、$j\leq j_{0}$となるすべてのjに対し、
    \[
    \vec{a}_{j} = \vec{0}
    \]
    となる時$\xi$は\eqref{eq:2-ode}の\textbf{確定特異点}という。確定特異点でない時\textbf{不確定特異点}という。
    （\cite[p.59 定義2.4]{book:okamoto} を参照）
\end{definition}

よって微分方程式の特異点$x=\xi$が微分方程式の意味で確定特異点でない時、命題\ref{prop:single_valued_z}の$\vec{z}(x)$は$x=\xi$を真性特異点としてもつ。\\

微分方程式だけでなく解析関数の分枝にも確定特異点は定義されている。\\
$\theta_{1},\theta_{2}\in\mathbb{R}$と$r>0$に対して、
\[
S(\theta_{1},\theta_{2}) = \{u\,|\, 0<|x-\xi|<r, \theta_{1} < \arg u < \theta_{2} \}
\]
とおく。ここで$u=x-\xi$である。

\begin{definition}[解析的な解の確定特異点]
    $x=\xi$で特異点$\omega$をもつ解析関数の分枝$\varphi(x)$に対して、$S(\theta_{1},\theta_{2})$で$\varphi(x)$が定義されるように$r$を十分小さくとる。
    ある$A>0$がとれて$\theta_{1}<\theta_{2}$に対し、
    \[
    \lim_{\substack{u\to 0 \\ u\in S(\theta_{1},\theta_{2})}}|u|^{A}|\varphi(x)| = 0
    \]
    となるとき、$\varphi(x)$は$x=\xi$に\textbf{確定特異点}を持つという。
    （\cite[p.60 定義2.5]{book:okamoto} を参照）
\end{definition}

\begin{Q}[なぜ特異点をわざわざ $\omega$ と書き分けているの？]
\textbf{A.} $\varphi(x)$ が対数関数や累乗根のような多価関数であり、その「分枝（branch）」を厳密に区別するためである。
\begin{itemize}
    \item \textbf{厳密な理由（原著の立場）:} 複素幾何学（リーマン面の理論）の観点では、ベース空間の $x=\xi$ という1点の上にも、分枝の選び方によって異なる複数の特異点 $\omega_1, \omega_2, \dots$ が重なって存在し得る。
    \item \textbf{本ノートでの扱い:} 確定特異点の定義における極限評価 $\lim |u|^A|\varphi(x)|=0$ では、ベース空間での距離 $|u| = |x-\xi|$ のみが計算に関与する。「特定の分枝を選び、扇形領域内で近づく」という前提さえあれば解析的な論理は完結するため、本ノートでは無駄を省き、原則として $x=\xi$ と表記を統一する。
\end{itemize}
\end{Q}

\begin{example}\label{ex:regular_singular_points}
    $\alpha\in\mathbb{C}$に対して、以下の関数の特異点の性質を調べる。
    
    \textbf{(1)} $x^{\alpha}$ は $x=0$ に確定特異点を持つ。
    \begin{Q}[なぜ $x=0$ は確定特異点といえるのか？]
    \textbf{A.} 定義に従い $\lim_{x\to 0} |x|^A |x^\alpha| = 0$ となる $A>0$ が存在するか調べる。
    $x = r e^{i\theta}$ とおくと、
    \[
    |x^{\alpha}| = |e^{\alpha(\log r + i\theta)}| = e^{\text{Re}(\alpha)\log r - \text{Im}(\alpha)\theta} = e^{-\text{Im}(\alpha)\theta} r^{\text{Re}(\alpha)}
    \]
    となる。扇形領域内で $\theta$ は有界であるから、$e^{-\text{Im}(\alpha)\theta}$ はある定数 $B(\theta)$ で抑えられる。
    ここで $|x|^A |x^\alpha| \propto r^{A + \text{Re}(\alpha)}$ となるため、$A + \text{Re}(\alpha) > 0$、すなわち $A > -\text{Re}(\alpha)$ となるように十分大きな $A>0$ をとれば、極限は $0$ になる。
    \end{Q}

    \textbf{(2)} $x^{\alpha}\log x$ は $x=0$ に確定特異点を持つ。
    \begin{Q}[$\log x$ があっても確定特異点になるのはなぜ？]
    \textbf{A.} 対数関数の発散は、多項式（べき乗）の減衰よりも遅いためである。
    (1)と同様に考えると $|x^{\alpha}\log x| \propto r^{\text{Re}(\alpha)}|\log r + i\theta|$ となる。
    極限 $\lim_{r\to 0} r^{A + \text{Re}(\alpha)} |\log r + i\theta|$ を評価する際、任意の $\epsilon > 0$ に対して $\lim_{r\to 0} r^\epsilon \log r = 0$ であるから、(1)と同じく $A > -\text{Re}(\alpha)$ となるように $A$ をとれば、多項式部分が $\log$ を押し潰して極限は $0$ になる。
    \end{Q}

    \textbf{(3)} $x^{\alpha}e^{x}$ の特異点 $x=\infty$ は確定特異点ではない（不確定特異点である）。
    \begin{Q}[なぜ $x=\infty$ で条件を満たさないのか？]
    \textbf{A.} 指数関数の爆発スピードが、いかなる多項式の減衰スピードも上回ってしまうためである。
    $u = 1/x$ とおいて $u \to 0$ の極限を考える。例えば正の実軸に沿って $u \to +0$ と近づく場合、
    \[
    |u|^A |x^{\alpha}e^{x}| = |u|^{A - \text{Re}(\alpha)} e^{\frac{1}{u}}
    \]
    となる。$u \to +0$ のとき、指数関数 $e^{1/u}$ の発散は多項式 $|u|^{A-\text{Re}(\alpha)}$ の減衰よりも圧倒的に速いため、$A$ をどれだけ大きく（正に）とっても極限は $0$ にならない。したがって確定特異点ではない（真性特異点である）。
    \end{Q}
    
    （\cite[p.60 例2.2]{book:okamoto} を参照）
\end{example}

また、次の極めて重要な命題が成り立つ。

\begin{proposition}\label{prop:equivalence_regular_singular}
    線型微分方程式\eqref{eq:2-ode}の特異点$x=\xi$が確定特異点であるための必要十分条件は、\eqref{eq:2-ode}の任意の解の基本系$\vec{\phi}(x)$の各成分$\phi_{1}(x),\phi_{2}(x)$が$x=\xi$上に確定特異点を持つことである。
    （\cite[p.61 命題2.6]{book:okamoto} を参照）
\end{proposition}

\begin{Q}[代数的な定義（ローラン展開）と解析的な定義（増大度）が一致するのはなぜ？]
\textbf{A.} 複素関数論における「孤立特異点の分類定理（リーマンの定理）」が背後で強力に効いているためである。証明の核心は前節で示した定理\ref{thm:growth_and_pole}（極と増大度の同値性）にある。この定理を見通しよく使うために、まず以下の補題を用意する。
\end{Q}

\begin{lemma}[多項式増大度の和と積]\label{lem:polynomial_growth_closure}
関数 $f(x), g(x)$ が $x=\xi$ （$u = x-\xi = 0$）の近傍で「多項式増大度」を持つとする。すなわち、ある定数 $A, B \ge 0$ と $C_1, C_2 > 0$ が存在して、十分小さな穴あき近傍 $0 < |u| < r$ で
\[
|f(x)| \le C_1|u|^{-A}, \quad |g(x)| \le C_2|u|^{-B}
\]
を満たすとする。このとき、$f(x)+g(x)$ および $f(x)g(x)$ もまた $x=\xi$ の近傍で多項式増大度を持つ。
具体的には、ある定数 $C_3, C_4 > 0$ に対して以下が成り立つ。
\begin{enumerate}
    \item $|f(x)+g(x)| \le C_3|u|^{-\max(A, B)}$
    \item $|f(x)g(x)| \le C_4|u|^{-(A+B)}$
\end{enumerate}
\end{lemma}

\begin{proof}
和については三角不等式より $|f(x)+g(x)| \le |f(x)| + |g(x)| \le C_1|u|^{-A} + C_2|u|^{-B}$ であり、$|u|<1$ の近傍では次数の低い（マイナスに大きい）方が支配的になるため、定数を調整すれば $|u|^{-\max(A, B)}$ で抑えられる。積についてはそのまま $|f(x)g(x)| \le C_1 C_2 |u|^{-(A+B)}$ となる。
\end{proof}

\begin{proof}[命題\ref{prop:equivalence_regular_singular}の証明]
\begin{itemize}
    \item \textbf{【$\Rightarrow$ の証明（ODE確定 $\implies$ 解析的確定）】}\\
    ODEの確定特異点の定義より、命題\ref{prop:single_valued_z}で定まる1価関数 $z_{k}(x)$ のローラン展開は有限の負べきで打ち切られる（すなわち高々極である）。よって定理\ref{thm:growth_and_pole}より、$z_k(x)$ は $x=\xi$ の近傍で多項式増大度を持つ。\\
    また、基本解行列 $\Phi(x)$ の各成分 $\Phi_{kj}(x)$ は $u^{\alpha}$ や $u^{\alpha}\log u$ の線型結合であるため、例題\ref{ex:regular_singular_points}で見たようにこれらも多項式増大度を持つ。\\
    ここで、解の各成分 $\phi_i(x)$ は
    \[
    \phi_{i}(x) = z_{1}(x)\Phi_{1i}(x) + z_{2}(x)\Phi_{2i}(x)
    \]
    と表される。これは「多項式増大度を持つ関数の積と和」であるから、補題\ref{lem:polynomial_growth_closure}より $\phi_i(x)$ 自身も多項式増大度を持つ。すなわち、解析的な意味で確定特異点であることが示された。

    \item \textbf{【$\Leftarrow$ の証明（解析的確定 $\implies$ ODE確定）】}\\
    仮定より、解 $\phi_i(x)$ は $x=\xi$ の近傍で多項式増大度を持つ。\\
    これをODEの確定特異点の定義（$z_k(x)$ のローラン展開が有限の負べきで打ち切られること）に結びつけるには、定理\ref{thm:growth_and_pole}より、$z_k(x)$ が多項式増大度を持つことを示せば十分である。\\
    命題\ref{prop:single_valued_z}より $\vec{z}(x) = \vec{\phi}(x)\Phi(x)^{-1}$ であるから、各成分は
    \[
    z_{i}(x) = \phi_{1}(x)\Phi_{1i}^{-1}(x) + \phi_{2}(x)\Phi_{2i}^{-1}(x)
    \]
    と書ける。ここで、逆行列 $\Phi(x)^{-1}$ の各成分 $\Phi_{ji}^{-1}(x)$ は $u^{-\alpha}$ や $u^{-\alpha}\log u$ の線型結合で表される（$\log u$ の逆数などは出現しない）ため、これらもやはり多項式増大度を持つ。\\
    したがって、補題\ref{lem:polynomial_growth_closure}より、その積と和で表される $z_i(x)$ も多項式増大度を持つ。\\
    $z_i(x)$ は $x=\xi$ の周りで1価の解析関数であり、かつ多項式増大度を持つため、定理\ref{thm:growth_and_pole}によりそのローラン展開は有限の負べきで必ず打ち切られる。これはまさにODEの確定特異点の定義そのものである。
\end{itemize}
\end{proof}

\begin{Q}[【補足】$\Phi(x)^{-1}$ を計算するとき、$\log u$ の逆数（$1/\log u$）のような扱いにくい項は出てこないの？]
\textbf{A.} 全く出てこない。行列の逆行列計算の美しい性質により、$\log u$ は分子に留まる。
$\alpha - \beta \in \mathbb{Z}$ でジョルダン細胞になる場合（最も複雑なケース）、$\Phi(x)$ は
\[
\Phi(x) = 
\begin{pmatrix}
    u^\alpha & K u^\alpha \log u \\
    0 & u^\alpha
\end{pmatrix}
\]
という形であった（$K$ は定数）。この行列の行列式は $\det \Phi(x) = u^{2\alpha}$ である。逆行列 $\Phi(x)^{-1}$ を真面目に計算すると、
\[
\Phi(x)^{-1} = \frac{1}{u^{2\alpha}}
\begin{pmatrix}
    u^\alpha & -K u^\alpha \log u \\
    0 & u^\alpha
\end{pmatrix}
=
\begin{pmatrix}
    u^{-\alpha} & -K u^{-\alpha} \log u \\
    0 & u^{-\alpha}
\end{pmatrix}
\]
となる。見ての通り、$\Phi(x)^{-1}$ の成分に現れるのは $u^{-\alpha}$ と $u^{-\alpha} \log u$ のみであり、$1/\log u$ は出現しない。これらは例題\ref{ex:regular_singular_points}で確認した通り、すべて多項式オーダー $|u|^{-M}$ で安全に上から抑えることができる。
\end{Q}

\begin{Q}[【補足】「ローラン展開の負べきが有限個（極）」だと、なぜ $|u|^{-N}$ で抑えられると直感的に言えるの？]
\textbf{A.} 原点付近（$u \to 0$）では、最も次数の低い（マイナスに大きい）項が圧倒的に支配的になるからである。
$z_k(x)$ が $x=\xi$ ($u=0$) に $N$ 位の極を持つとすると、ローラン展開は
\[
z_k(x) = \frac{a_{-N}}{u^N} + \frac{a_{-N+1}}{u^{N-1}} + \cdots + a_0 + a_1 u + \cdots
\]
となる。両辺に $u^N$ をかけると、
\[
u^N z_k(x) = a_{-N} + a_{-N+1}u + \cdots 
\]
となり、これは $u \to 0$ で正則（有限の値 $a_{-N}$ に収束）である。したがって、$u$ が十分小さい領域では $|u^N z_k(x)| \le C$ (ある定数) と有界になり、両辺を割ることで $|z_k(x)| \le C|u|^{-N}$ という多項式オーダーの評価が直ちに得られる。無限に続く真性特異点（例：$e^{1/u}$）ではこの議論が破綻する。（※厳密な同値性の証明は定理\ref{thm:growth_and_pole}を参照）
\end{Q}
%#endregion
%#region
\begin{definition}[ガウスの微分方程式]
    複素パラメータ $a, b, c$ に対して、
    \begin{equation}\label{eq:gauss_hypergeometric}
    x(1-x)\frac{d^{2}y}{dx^{2}}+(c-(a+b+1)x)\frac{dy}{dx}-aby = 0
    \end{equation}
    と書かれる2階線形微分方程式を\textbf{ガウスの微分方程式（超幾何微分方程式）}という。
\end{definition}

\begin{lemma}[ガウスの超幾何関数]\label{lem:hypergeometric}
    \[
        F(a,b,c;x)=\sum_{j=0}^{\infty}\frac{(a)_{j}(b)_{j}}{(c)_{j}j!}x^{j} \qquad \left( \text{ただし } (a)_{j}=a(a+1)(a+2)\dots(a+j-1) \right)
    \]
    の収束半径は1である。これは\textbf{ガウスの超幾何関数}と呼ばれる。
\end{lemma}

\begin{proof}
    命題\ref{prop:dalembert_radius}より収束半径は
    \[
    \lim_{j\to\infty}\left| \frac{(a)_{j}(b)_{j}}{(c)_{j}j!} \frac{(c)_{j+1}(j+1)!}{(a)_{j+1}(b)_{j+1}} \right| = \lim_{j\to\infty}\left| \frac{(c+j)(j+1)}{(a+j)(b+j)} \right| = 1
    \]
    となる。
\end{proof}

\begin{Q}[【核心】ガウスの超幾何関数の収束半径が「1」になることは、ガウスの微分方程式の係数の形からどのように「予言」できるか？]
\textbf{A.} ガウスの微分方程式を標準形 $\frac{d^2y}{dx^2} + p_1(x)\frac{dy}{dx} + p_2(x)y = 0$ で表すと、分母に $x(1-x)$ が来るため、方程式の特異点は $x=0$ と $x=1$（および $\infty$）であることがわかる。超幾何関数 $F(a,b,c;x)$ は原点 $x=0$ まわりでの解であるため、そこから「最も近い次の特異点 $x=1$」までの距離がちょうど1であり、これが収束半径の限界と一致するからである。
\end{Q}
%#endregion
%#region
\begin{lemma}
    \eqref{eq:gauss_hypergeometric}は$x=0,1,\infty$で特異点をもつ。
\end{lemma}

\begin{proof}
    $x=0, 1$ が特異点であることは、方程式の両辺を $x(1-x)$ で割ることで係数の分母がゼロになることからすぐにわかる。

    無限遠点 $x=\infty$ については、$x=\frac{1}{u}$ を代入して $u=0$ での振る舞いを調べる。微分の連鎖律を用いて \eqref{eq:gauss_hypergeometric} を書き換えると、
    \begin{align*}
    \frac{1}{u}\left(1-\frac{1}{u}\right)\left(u^{4}\frac{d^{2}y}{du^{2}}+2u^{3}\frac{dy}{du}\right)+\left(c-(a+b+1)\frac{1}{u}\right)\left(-u^{2}\frac{dy}{du}\right)-aby &= 0 \\
    u^{2}(u-1)\frac{d^{2}y}{du^{2}}+2u(u-1)\frac{dy}{du}-u(cu-(a+b+1))\frac{dy}{du}-aby &= 0 \\
    u^{2}(u-1)\frac{d^{2}y}{du^{2}}+u((2-c)u+(a+b-1))\frac{dy}{du}-aby &= 0
    \end{align*}
    となる。両辺を $u^{2}(u-1)$ で割ると、各係数が $u=0$ を極（特異点）として持つことがわかる。つまり $x=\infty$ も特異点である。

    また、これら以外の点では係数は明らかに正則であるため、特異点の集合は $\{0,1,\infty\}$ である。
\end{proof}

\begin{Q}[【核心】なぜ複素平面だけでなく、無限遠点 $\infty$ を含めた「リーマン球面」全体で特異点の個数を数える必要があるのか？]
\textbf{A.} リーマン球面全体で見渡したとき、ガウスの微分方程式は「特異点をちょうど3つ $\{0, 1, \infty\}$ だけ持つ」という極めて対称的で純粋な構造を持っていることがわかり、実はこの「特異点が3つ」という幾何学的な条件だけで方程式の形が（変数変換を除いて）一意に決定されてしまうという、複素微分方程式の根源的な性質（リーマンのP方程式）に繋がるからである。
\end{Q}
%#endregion
%#region
\begin{lemma}\label{lem:hypergeometric_solution}
    $F(a,b,c;x)$ は \eqref{eq:gauss_hypergeometric} の解である。
\end{lemma}

\begin{proof}
    級数解 $y = \sum_{j=0}^{\infty}\frac{(a)_{j}(b)_{j}}{(c)_{j}j!}x^{j}$ の微分はそれぞれ、
    \[
    y' = \sum_{j=1}^{\infty}\frac{(a)_{j}(b)_{j}}{(c)_{j}(j-1)!}x^{j-1}, \quad
    y'' = \sum_{j=2}^{\infty}\frac{(a)_{j}(b)_{j}}{(c)_{j}(j-2)!}x^{j-2}
    \]
    となる。これらを \eqref{eq:gauss_hypergeometric} の左辺に代入して展開すると、
    \begin{align*}
    (\text{左辺}) &= x(1-x)\sum_{j=2}^{\infty}\frac{(a)_{j}(b)_{j}}{(c)_{j}(j-2)!}x^{j-2} + (c-(a+b+1)x)\sum_{j=1}^{\infty}\frac{(a)_{j}(b)_{j}}{(c)_{j}(j-1)!}x^{j-1} - ab\sum_{j=0}^{\infty}\frac{(a)_{j}(b)_{j}}{(c)_{j}j!}x^{j} \\
    &= \sum_{j=1}^{\infty}\frac{(a)_{j+1}(b)_{j+1}}{(c)_{j+1}(j-1)!}x^{j} - \sum_{j=2}^{\infty}\frac{(a)_{j}(b)_{j}}{(c)_{j}(j-2)!}x^{j} \\
    &\quad + c\sum_{j=0}^{\infty}\frac{(a)_{j+1}(b)_{j+1}}{(c)_{j+1}j!}x^{j} - (a+b+1)\sum_{j=1}^{\infty}\frac{(a)_{j}(b)_{j}}{(c)_{j}(j-1)!}x^{j} - ab\sum_{j=0}^{\infty}\frac{(a)_{j}(b)_{j}}{(c)_{j}j!}x^{j}
    \end{align*}
    となる（$x y''$ と $c y'$ の項ではインデックスを $j \to j+1$ とずらした）。
    定数項（$x^0$ の係数）は $c\frac{ab}{c} - ab = 0$ となる。
    $x^j \ (j \ge 1)$ の係数を集め、$\frac{(a)_{j}(b)_{j}}{(c)_{j}j!}$ をくくり出すと、
    \begin{align*}
    & \frac{(a)_{j+1}(b)_{j+1}}{(c)_{j+1}j!} (j + c) - \frac{(a)_{j}(b)_{j}}{(c)_{j}j!} \left( j(j-1) + (a+b+1)j + ab \right) \\
    &= \frac{(a)_{j}(b)_{j}}{(c)_{j}j!} \left[ \frac{(a+j)(b+j)}{c+j}(j+c) - (j^2 + (a+b)j + ab) \right] \\
    &= \frac{(a)_{j}(b)_{j}}{(c)_{j}j!} \left[ (a+j)(b+j) - (a+j)(b+j) \right] = 0
    \end{align*}
    となり、すべての次数の係数が $0$ となる。したがって $F(a,b,c;x)$ は方程式を満たす。
\end{proof}

\begin{Q}[【核心】ガウスの超幾何関数という複雑な級数は、どこから天下り的にやってきたのか？]
\textbf{A.} 天下りではなく、未定係数法による必然である。微分方程式の解を $y = \sum A_j x^j$ と仮定して代入し、各次数の係数がゼロになるように整理すると、$A_{j+1}$ と $A_j$ の間に $A_{j+1} = \frac{(a+j)(b+j)}{(c+j)(j+1)}A_j$ という2項間漸化式が必ず導かれる。これを $A_0 = 1$ から順次解き下した結果が、ポッホハマー記号を用いた $F(a,b,c;x)$ の係数に他ならない。
\end{Q}
%#endregion
%#region
\begin{definition}[複素領域における解の「分枝」の1次独立性とロンスキアン]\label{def:linear_independence_complex}
    複素領域における2階線形微分方程式について、特異点を含まないある単連結な領域（例えば小さな開円板やスリット領域）$D$ 上で固定された\textbf{解の分枝のペア $(\varphi_1(x), \varphi_2(x))$} が\textbf{1次独立（線形独立）}であるとは、任意の複素定数 $C_1, C_2$ に対して $D$ 上で
    \[ 
    C_1\varphi_1(x) + C_2\varphi_2(x) \equiv 0 \implies C_1 = 0 \text{ かつ } C_2 = 0 
    \]
    が成り立つことをいう。解析的には、領域 $D$ 上でロンスキアン $W(\varphi_1, \varphi_2)(x) \not\equiv 0$ であることと同値である。
    
    （※大域的な多価関数そのものが無条件に独立であるかを問うのではなく、あくまで「局所領域 $D$ 上で切り出された特定の分枝の組み合わせ」に対して独立性が定義・判定されることに注意する。）
\end{definition}

\begin{Q}[【核心】局所的な領域 $D$ での判定だけで、多価関数という複雑な対象の「全域での独立性」を保証できるのはなぜか？]
\textbf{A.} 複素解析の「一致の定理」と微分方程式の「アーベルの公式」が、局所と大域を鉄の論理で結んでいるからである。
ある1つの領域 $D$ でロンスキアンが非ゼロ（独立）ならば、特異点を回って別の分枝に移っても、それは正則なモノドロミー行列を掛ける「基底の変換」に過ぎず、独立性は全域で不変である。逆に、もし最初から従属なペアを選んでいれば、領域 $D$ での判定で確実に $0$ となり、偽陽性（誤判定）は絶対に起こらない。
\end{Q}

\begin{Q}[【核心】2つの多価関数を比較する際、なぜ「解析接続の経路を同期させる（シートを揃える）」ことが絶対ルールなのか？]
\textbf{A.} リーマン面という螺旋階段において、同じ階層（シート）で評価しなければ数学的に無意味だからである。
同じ複素座標 $x$ であっても、多価関数には「0周目の $x$」や「1周目の $x$」といった異なる実体が存在する。これらを混同して比較すると、本来独立なはずの解がたまたま同じ向きを向いて従属に見えてしまう「分枝の罠」に嵌まる。これを防ぐため、必ず共通の基準点から「全く同じ経路」で解析接続された値同士を戦わせる必要がある。
\end{Q}

\begin{Q}[【核心】特異点を回って別の分枝に移動したとき、解を結びつけていた定数 $C_1, C_2$ が変化してしまう心配はないのか？]
\textbf{A.} 全くない。$C_1, C_2$ は解析接続の影響を受けない「普遍的な複素定数」だからである。
2つの解を同じ経路で同期させて解析接続する限り、多価性によって関数自身の値が $\lambda$ 倍に変化しても、等式の両辺でその因子が約分されるため、線形関係を規定する定数は全分枝において完全に同一のまま保たれる。
\end{Q}
%#endregion
%#region
\begin{example}[$x=0$ まわりの局所解]\label{ex:gauss_solutions_zero}
    $c\notin\mathbb{Z}$ ならば、ガウスの微分方程式 \eqref{eq:gauss_hypergeometric} は $x=0$ の近傍において、
    \[
    \varphi_{1}(x)=F(a,b,c;x),\quad \varphi_{2}(x)=x^{1-c}F(a-c+1,b-c+1,2-c;x)
    \]
    という1次独立な2つの解をもつ。
\end{example}

\begin{proof}
    定理と命題を用いて、以下の3つのステップに分けて証明する。

    \begin{description}
        \item[【ステップ1】$\varphi_2(x)$ が解であること]\mbox{}\\
        $\varphi_{1}(x)$ が解であることは既に補題\ref{lem:hypergeometric_solution}で示した。$\varphi_{2}(x)$ について、$y=x^{1-c}z$ とおいて \eqref{eq:gauss_hypergeometric} に代入し、微分の連鎖律を適用すると、
        \begin{gather*}
        x(1-x)\left\{ c(c-1)x^{-1-c}z + 2(1-c)x^{-c}\frac{dz}{dx} + x^{1-c}\frac{d^{2}z}{dx^{2}} \right\} \\
        + (c-(a+b+1)x)\left\{ (1-c)x^{-c}z + x^{1-c}\frac{dz}{dx} \right\} - abx^{1-c}z = 0
        \end{gather*}
        となる。超幾何微分方程式の標準形に合わせるため、この両辺に $x^{c-1}$ を掛けて整理すると、
        \begin{gather*}
        x(1-x)\frac{d^{2}z}{dx^{2}} + \left\{ 2(1-c)(1-x) + (c-(a+b+1)x) \right\}\frac{dz}{dx} \\
        + \left\{ c(c-1)x^{-1}(1-x) + (c-(a+b+1)x)(1-c)x^{-1} - ab \right\}z = 0 \\
        \intertext{$z', z$ の係数をそれぞれ展開して因数分解すると、}
        x(1-x)\frac{d^{2}z}{dx^{2}} + \left\{ (2-c) - ((a-c+1)+(b-c+1)+1)x \right\}\frac{dz}{dx} - (a-c+1)(b-c+1)z = 0
        \end{gather*}
        が得られる。これはパラメータを $(a-c+1), (b-c+1), (2-c)$ としたガウスの超幾何微分方程式そのものである。
        したがって、補題\ref{lem:hypergeometric_solution}よりその解として $z=F(a-c+1,b-c+1,2-c;x)$ がとれる。よって、$y=x^{1-c}F(a-c+1,b-c+1,2-c;x)$ は \eqref{eq:gauss_hypergeometric} の解である。

        
        \item[【ステップ2】2つの解が1次独立であること]\mbox{}\\
        まず、特異点 $x=0$ を含まない単連結な領域 $D$ 上で $\log x$ の分枝を1つ固定し、$\varphi_2(x)$ を $D$ 上の一価正則な関数として確定させる。
        このとき、$\varphi_1, \varphi_2$ が $D$ 上で1次独立でないと仮定すると、ある定数 $(C_{1},C_{2})\in\mathbb{C}^{2}\setminus\{(0,0)\}$ に対して、
        \[
        C_{1}\varphi_{1}(x)+C_{2}\varphi_{2}(x)\equiv 0
        \]
        が成り立つ。命題\ref{prop:independence_branch_invariance}より、この等式全体を $x=0$ の周りに反時計回りに1周する共通の経路に沿って解析接続したとき、元の領域 $D$ に戻ってきて得られる新しい分枝同士の間でも、同じ定数 $(C_1, C_2)$ による線形関係が成立する。

        1周解析接続した結果、$\varphi_1(x)$ は一価関数であるため元の分枝と完全に一致するが、$\varphi_2(x)$ は $x^{1-c} = \exp((1-c)\log x)$ の因子を持つため、元の分枝から $e^{2\pi i(1-c)}$ 倍された値となる。したがって $D$ 上で、
        \[
        C_{1}\varphi_{1}(x)+C_{2}e^{2\pi i(1-c)}\varphi_{2}(x)\equiv 0
        \]
        となる。この新旧2つの等式の差をとると $C_2(e^{2\pi i(1-c)}-1)\varphi_2(x) \equiv 0$ となるが、$c \notin \mathbb{Z}$ より $1-c$ は整数ではないため $e^{2\pi i(1-c)} \neq 1$ である。よって $C_2=0$ となり、最初の式から $C_1=0$ も導かれるため仮定に矛盾する。ゆえに、選び出した分枝の組によらず $\varphi_1, \varphi_2$ は1次独立である。
        
        \item[【ステップ3】解がこれら2つで全てであること]\mbox{}\\
        もし、上に挙げた2つ以外にこれらと1次独立な解が存在すると仮定すると、解空間の次元が3次元以上となってしまう。しかしこれは、定理\ref{thm:solution_space_dimension}で保証された「解空間が2次元ベクトル空間をなす」という事実に矛盾する。よって、局所解はこの2つのみで完全に尽くされている。
    \end{description}
\end{proof}

\begin{Q}[【核心】切り込み（スリット）を越えて解析接続を行うと、分枝の定義域を超えてしまい論理が破綻しないか？]
\textbf{A.} 破綻しない。なぜなら、切り込みの上で方程式を評価するのではなく、始点の「安全な領域 $D$」上の等式を出発点とし、1周回って「再び領域 $D$ に戻ってきた新しい分枝」と元の分枝を比較しているからである。一致の定理により定数の関係は保たれるため、同じ $D$ 上での比較が正当化される。
\end{Q}

\begin{Q}[【核心】$\varphi_{2}(x)$ は $x=0$ で多価性や発散を伴うが、微分方程式の「解」として問題ないか？]
\textbf{A.} 全く問題ない。$x=0$ は方程式自身の「特異点」であるため、解が特異性を持つことは自然である。むしろ特異点の周りでは、このような $x^\rho \times (\text{正則関数})$ の形（フロベニウスの解）を想定して基本解を探すのが定石である。
\end{Q}

\begin{Q}[【核心】$x=0$ の真上ではコーシーの定理が使えないのに、なぜ解空間が「2次元」であり、この2つの解で「全て」だと断言できるのか？]
\textbf{A.} $x=0$ を除外した「穴あき近傍（$0 < |x| < r$）」で考えているからである。正則点で確立された「解空間は2次元」という代数的な事実は穴あき近傍でも変わらないため、そこで1次独立な2つの解を見つけさえすれば、いかなる解もその線形結合で完全に表現できる。
\end{Q}

\begin{Q}[【核心】特異点を回って解が「多価関数」になる現象は、「解空間は2次元」という枠組みと矛盾しないか？]
\textbf{A.} 矛盾しない。方程式の係数が一価であるため、特異点を1周して得た「新しい分枝」も再び解となる。解空間が局所的に2次元である以上、いかに多価に分岐しようとも、その新しい分枝は必ず元の2つの基底の線形結合で表され、多価性は「2次元空間内の線形変換（モノドロミー行列）」として吸収される。
\end{Q}

\begin{Q}[【実践】$c \notin \mathbb{Z}$ のとき、2つの解の1次独立性が計算せずとも直ちにわかるのはなぜか？]
\textbf{A.} 特異点 $x \to 0$ における漸近挙動が、一方は $\varphi_1(x) \to 1$、もう一方は $\varphi_2(x) \sim x^{1-c}$ と決定的に異なるからである。先頭のオーダーが異なる関数同士が定数倍で結びつくことは絶対に不可能であるため、これだけで1次独立性が確定する。
\end{Q}
%#endregion
%#region
これまでの議論から、$\{\varphi_{1}(x), \varphi_{2}(x)\}$ は方程式 \eqref{eq:gauss_hypergeometric} の基本系をなし、それぞれ $x=0$ に確定特異点をもつことがわかっている。

\begin{Q}[命題\ref{prop:equivalence_regular_singular}を適用するには「\textbf{任意の}解の基本系」が $x=0$ に確定特異点をもつことを示す必要がある。なぜ特定の基本系 $\{\varphi_{1}(x), \varphi_{2}(x)\}$ の性質から「任意」へと議論を拡張できるのか？]
\textbf{A.} 微分方程式の線型性と、定義\ref{def:linear_independence_complex}における確定特異点（増大度）の条件が、線型結合によって保たれる性質だからである。

具体的には以下の論理による：
\begin{enumerate}
    \item $\{\varphi_{1}(x), \varphi_{2}(x)\}$ が基本系であるため、\eqref{eq:gauss_hypergeometric} の任意の解は $c_{1}\varphi_{1}(x) + c_{2}\varphi_{2}(x)$ という定数係数の線型結合で一意に表される。
    \item 確定特異点をもつ（すなわち高々多項式オーダーの増大度をもつ）関数の線型結合もまた、同じ特異点において確定特異点をもつ。
    \item ゆえに、\eqref{eq:gauss_hypergeometric} の「任意の解の基本系」を構成する各成分も、すべて $x=0$ に確定特異点をもつことが保証される。
\end{enumerate}
\end{Q}

上記の理由により、\eqref{eq:gauss_hypergeometric} の任意の解の基本系は $x=0$ を確定特異点にもつ。よって、命題\ref{prop:equivalence_regular_singular}より、$x=0$ は微分方程式 \eqref{eq:gauss_hypergeometric} の確定特異点である。また、$x=1,x=\infty$も確定特異点であることが知られている。
%#endregion
%#region
\begin{definition}
    \eqref{eq:2-ode}の特異点$x=\xi$がすべて確定特異点である時、\eqref{eq:2-ode}を\textbf{フックス型微分方程式}という。
\end{definition}
%#endregion
%#region
\begin{example}
    ガウスの超幾何微分方程式 \eqref{eq:gauss_hypergeometric} はフックス型微分方程式である。
\end{example}

\begin{proof}
    今までの議論から、\eqref{eq:gauss_hypergeometric} の特異点は $x=0,1,\infty$ のみであり、それらはすべて確定特異点であることがわかっているので、定義からフックス型微分方程式である。
\end{proof}

\begin{Q}[ガウスの超幾何微分方程式のように「リーマン球面上に確定特異点をちょうど3つだけ持つ」フックス型方程式は、パンルヴェ方程式などの非線形理論へ向かう上でどのような特別な立ち位置にあるか？]
\textbf{A.} 「変形」の余地を持たない、最も基本的で**「剛（rigid）」**な方程式としての立ち位置です。

リーマン球面上の任意の3点は、一次分数変換（モビウス変換）によって常に $0, 1, \infty$ に移すことができます。そのため、確定特異点が3つのフックス型方程式は、各特異点における特性指数（局所的な振る舞い）を指定するだけで**方程式全体が一意に決定されてしまいます**（リーマンのP関数）。
ここから特異点を4つ（例えば $0, 1, \infty, t$）に増やし、その位置 $t$ を動かしても方程式の全体的な性質（モノドロミー）が変わらないように要請する（モノドロミー保存変形）ことで、初めてパンルヴェ方程式のような非線形微分方程式が舞台に上がってきます。
\end{Q}
%#endregion
%#region
\begin{proposition}
    \eqref{eq:2-ode}の特異点$x=\xi$が確定特異点になる必要十分条件は$(x-\xi)^{i}p_{i}(x)$が$x=\xi$で正則になることである。
\end{proposition}

\begin{proof}
    命題\ref{prop:fuchs_theorem_n}の$n=2$の場合である。
\end{proof}
%#endregion
%#region
上の命題で冪級数の構成部分を再考する。
特異点 $x=\xi$ の周りでの局所座標を $u = x - \xi$ とおく。
\eqref{eq:2-ode}の両辺に $u^{2}$ をかけると、$\frac{d}{dx} = \frac{d}{du}$ であるから、
\[
u^{2}\frac{d^{2}y}{du^{2}} + (u p_{1}(x)) u\frac{dy}{du} + (u^{2} p_{2}(x)) y = 0
\]
となる。

\begin{Q}[【核心】なぜわざわざ $u^2$ を掛け、$u p_1(x)$ と $u^2 p_2(x)$ が正則であるという状況を考えるのか？]
\textbf{A.} この条件こそが、$x=\xi$ が方程式の「確定特異点（regular singular point）」であることの定義だからである。フックス（Fuchs）の定理により、特異点であってもこの条件を満たせば、解の少なくとも1つは $u^s \sum c_n u^n$ という冪級数（フロベニウス級数）の形で正則な解を持つことが保証されるため、この前提が不可欠となる。
\end{Q}

上の仮定より、$u p_{1}(x), u^{2} p_{2}(x)$ は $u=0$ において正則なので、それぞれ $u$ の関数として $q(u), r(u)$ と置き直すと、方程式は次のように書ける。
\begin{equation}\label{eq:ode_u}
u^{2}\frac{d^{2}y}{du^{2}} + q(u) u\frac{dy}{du} + r(u) y = 0
\end{equation}
ここで、解 $\varphi(u)$ が
\[
\varphi(u) = u^{s}\sum_{n=0}^{\infty}c_{n}u^{n} \quad (c_{0}\neq 0)
\]
と書けるとする。$q(u), r(u)$ もテイラー展開して
\[
q(u) = \sum_{n=0}^{\infty}q_{n}u^{n}, \quad r(u) = \sum_{n=0}^{\infty}r_{n}u^{n}
\]
とし、これらを\eqref{eq:ode_u}に代入して $u^{s+n}$ の係数を整理すると、
\[
\sum_{n=0}^{\infty} u^{s+n} \left[ (s+n)(s+n-1)c_{n} + \sum_{k=0}^{n} (s+k)c_{k}q_{n-k} + \sum_{k=0}^{n} c_{k}r_{n-k} \right] = 0
\]
を得る。$n=0$（すなわち最低次 $u^s$ の係数）を見ると、
\[
(s(s-1) + sq_{0} + r_{0})c_{0} = 0
\]
となる。仮定より $c_{0} \neq 0$ であるため、次の方程式が得られる。
\begin{equation}\label{eq:indicial}
f(s) := s(s-1) + sq_{0} + r_{0} = 0
\end{equation}

\begin{definition}
\eqref{eq:indicial}を、$x=\xi$ における元の微分方程式\eqref{eq:2-ode}の\textbf{決定方程式}（または特性方程式）という。また、この2次方程式の2つの根を\textbf{決定指数}（特性指数）という。
\end{definition}

決定指数のうち、実部が大きい方（あるいは任意の1つ）を $s_{1}$ とし、もう1つを $s_{2}$ とする。
一般の $u^{s+n}$ の係数が $0$ になる条件から、以下の漸化式が得られる。
\begin{equation}\label{eq:recurrence}
f(s+n)c_{n} + \sum_{k=0}^{n-1} \left( (s+k)q_{n-k} + r_{n-k} \right) c_{k} = 0
\end{equation}
右側の和を既知の項 $R_n(c_{0}, c_{1}, \cdots, c_{n-1})$ とおけば、$f(s+n)c_{n} = -R_n$ となる。
したがって、$f(s+n) \neq 0$ であれば、$c_{n}$ は $c_{0}$ から順に一意に定まる。

決定指数の差 $s_{1} - s_{2}$ が整数でない場合、任意の自然数 $n$ に対して $f(s_{1}+n) \neq 0$ かつ $f(s_{2}+n) \neq 0$ となる。よって $s=s_{1}$ と $s=s_{2}$ のそれぞれから解 $\varphi_{1}(x), \varphi_{2}(x)$ が構成できる。これらは最低次の項（$u^{s_1}$ と $u^{s_2}$）の冪が異なるため、互いに1次独立な解となる。

\begin{Q}[【核心】決定指数の差が整数のとき、なぜ解に対数項 $\log u$ が加わることがあるのか？]
\textbf{A.} $s_{1} - s_{2} = J \in \mathbb{Z}_{\ge 0}$ と書ける場合、小さい方の根 $s=s_{2}$ を用いて漸化式\eqref{eq:recurrence}を解こうとすると、$n=J$ のステップで $f(s_2+J) = f(s_1) = 0$ となる。
このとき、漸化式は $0 \cdot c_{J} = -R_J$ となり、もし $R_J \neq 0$ であれば、これを満たす係数 $c_{J}$ は存在せず、解の構成が破綻してしまうからである。この破綻を吸収し独立な解を2つ見つけるためには、フロベニウス級数の形を拡張し、必然的に $\log u$ を含む項を導入しなければならない。
\end{Q}
%#endregion
%#region
\eqref{eq:2-ode}において、無限遠点 $x=\infty$ での振る舞いを調べるため、局所座標 $x = \frac{1}{u}$ を導入する。
連鎖律より $\frac{dy}{dx} = -u^{2}\frac{dy}{du}$、$\frac{d^{2}y}{dx^{2}} = 2u^{3}\frac{dy}{du} + u^{4}\frac{d^{2}y}{du^{2}}$ となるため、これを\eqref{eq:2-ode}に代入して $u^{4}$ で割ると、次の方程式を得る。
\begin{equation}\label{eq:ode_infinity}
\frac{d^{2}y}{du^{2}} + \left(\frac{2}{u} - \frac{1}{u^{2}}p_{1}\left(\frac{1}{u}\right)\right)\frac{dy}{du} + \frac{1}{u^{4}}p_{2}\left(\frac{1}{u}\right)y = 0
\end{equation}

\begin{Q}[【核心】なぜ $x=\infty$ の特異性を調べるのに、わざわざ $x=1/u$ という変換を行って $u=0$ での振る舞いを見るのか？]
\textbf{A.} 複素解析において、無限遠点 $\infty$ は特別な場所ではなく、リーマン球面 $\hat{\mathbb{C}} = \mathbb{C} \cup \{\infty\}$ 上の1点に過ぎないからである。しかし、我々が持っているフロベニウスの方法などの「局所的な解析手法」は原点（有限の点）の周りで構築されているため、チャート変換 $x=1/u$ を用いて $\infty$ を原点 $u=0$ に引き戻すことで、有限の特異点と全く同じ土俵で解析できるようになる。
\end{Q}

\eqref{eq:ode_infinity}に対して、有限の確定特異点の条件（フックスの定理の前提）を適用することで、以下の命題が成り立つ。

\begin{proposition}\label{prop:regular_infinity}
$x$ の有理関数を係数とする\eqref{eq:2-ode}に対して、無限遠点 $x=\infty$ が確定特異点になるための必要十分条件は、$\frac{1}{u}p_{1}\left(\frac{1}{u}\right)$ および $\frac{1}{u^{2}}p_{2}\left(\frac{1}{u}\right)$ が $u=0$ で正則になることである。
（※これは言い換えれば、$x \to \infty$ において $p_{1}(x) = O(x^{-1})$ かつ $p_{2}(x) = O(x^{-2})$ となることと同値である。）
\end{proposition}
%#endregion
%#region
次に、特異点において決定指数（特性指数）の差が整数となる例外的なケースを考察する。
通常、決定指数の差が整数の場合は対数項 $\log u$ が解に出現するが、漸化式における障害となる項 $R_J$ が偶然 $0$ になる場合、対数項を持たない正則な解が存在する。

\begin{definition}[非対数的特異点 / 見かけの特異点]\label{def:non_logarithmic}
確定特異点 $x=\xi$ における2つの決定指数 $s_{1}, s_{2}$ （$s_{1} \ge s_{2}$）の差が整数 $s_{1} - s_{2} = J \in \mathbb{Z}_{\ge 0}$ であり、かつ小さい方の根 $s=s_{2}$ に対する漸化式の第 $J$ ステップにおいて障害項が消滅する（すなわち $R_{J}(c_{0}, \cdots, c_{J-1}) = 0$ が成り立つ）とき、この特異点を\textbf{非対数的特異点}（または\textbf{見かけの特異点}）と呼ぶ。
\end{definition}

非対数的特異点においては、$n=J$ での漸化式が $0 \cdot c_{J} = 0$ となるため、$c_{J}$ は矛盾を引き起こすことなく、任意の複素定数 $h \in \mathbb{C}$ として自由に選んで計算を進めることができる。

\begin{Q}[【核心】小さい方の根 $s_1$ から出発して $c_J = h$ と置いた解は、なぜ大きい方の根 $s_2$ の解 $\varphi_{2}(x)$ を $h$ 倍して足し合わせた形になるのか？]
\textbf{A.} $c_J = h$ としたとき、$n > J$ における漸化式は $c_0, \cdots, c_{J-1}, h$ に依存する項 $R_n$ を持ち、次のように書ける。
\[
f(s_1+n)c_n + R_n(c_0, \cdots, c_{J-1}, h) = 0
\]
ここで、$R_n$ は各係数に対して線形であるため、次のように分離できる。
\[
R_n(c_0, \cdots, c_{J-1}, h) = R_n(c_0, \cdots, c_{J-1}, 0) + h R_n(0, \cdots, 0, 1)
\]
この分離により、以降の係数 $c_n$ も「$h=0$ としたときの係数」と「$h$ に比例する係数」の和に完全に分かれる。
後半の $h$ に比例する部分は、初項を $c_J=1$ として $u^{s_2+J} = u^{s_1}$ から始まる級数であり、これはまさに大きい根 $s_2$ に対応する解 $\varphi_{2}(x)$ そのものである。
したがって、任意の定数 $h$ に対する全体の解は、
\[
\varphi(x) = \varphi_{1}(x) + h\varphi_{2}(x)
\]
（ただし $\varphi_{1}(x)$ は $h=0$ とした場合の解）という形に自然に分離されるのである。
\end{Q}
決定指数の差が整数となる場合について、方程式を正規化してさらに深く考察する。
特異点 $x=\xi$ における小さい方の決定指数を $s_{1}$、大きい方の決定指数を $s_{2}$ とし、その差を $J = s_{2} - s_{1} \in \mathbb{Z}_{\ge 0}$ とする。
未知関数 $y$ を
\[
y = (x-\xi)^{s_{1}} w
\]
と変換して\eqref{eq:2-ode}を $w$ についての微分方程式に書き換えると、変換後の微分方程式の $x=\xi$ における決定指数は $s_{1}-s_{1}=0$ および $s_{2}-s_{1}=J$ となる。
この整数 $J \ge 0$ の値によって、変換後の特異点の性質は以下のように分類される。

\begin{itemize}
    \item \textbf{$J=0$ のとき:} 決定指数は $0, 0$ となり重根を持つため、必ず対数項が発生し、非対数的特異点にはなり得ない。
    \item \textbf{$J=1$ のとき:} 決定指数は $0, 1$ となる。これが非対数的特異点である場合、2つの独立な解はともに $x=\xi$ で正則となる。フックスの定理の逆から、このとき $x=\xi$ は変換後の微分方程式にとって特異点ではなく、単なる正則点（ordinary point）となる。
    \item \textbf{$J \ge 2$ のとき:} 変換後の微分方程式は $x=\xi$ で依然として確定特異点を持つ。これが非対数的特異点である場合、決定指数が $0, J$ であるため、変換後の任意の解は $x=\xi$ で正則なテイラー展開を持つ。
\end{itemize}

\begin{Q}[【核心】$J=1$ と $J \ge 2$ では、どちらも「2つの解が正則」になるのに、なぜ $J=1$ は「正則点」で、$J \ge 2$ は「見かけの特異点」になるのか？]
\textbf{A.} その違いは、解から微分方程式の係数を逆算したときの「ロンスキアン（Wronskian）」の値に現れる。
微分方程式 $w'' + P(x)w' + Q(x)w = 0$ の係数 $P, Q$ は、2つの解 $w_{1}, w_{2}$ のロンスキアン $W = w_{1}w_{2}' - w_{1}'w_{2}$ を分母に持つ。
決定指数が $0$ と $J$ のとき、$w_{1}$ は定数項から始まり（$w_{1}(\xi) \neq 0$）、$w_{2}$ は $(x-\xi)^{J}$ から始まる。
\begin{itemize}
    \item \textbf{$J=1$ の場合：} $w_{2}'(\xi) \neq 0$ となるため、$x=\xi$ において $W(\xi) \neq 0$ となる。分母が $0$ にならないため、係数 $P(x), Q(x)$ も正則となり、そもそも微分方程式にとって特異点ではない（正則点である）。
    \item \textbf{$J \ge 2$ の場合：} $w_{2}(\xi) = 0$ だけでなく $w_{2}'(\xi) = 0$ も成り立つため、$x=\xi$ において $W(\xi) = 0$ となってしまう。分母が $0$ になるため係数 $P(x), Q(x)$ は極を持ち、方程式の特異性が残る（見かけの特異点となる）。
\end{itemize}
\end{Q}

\begin{Q}[【核心】$W(\xi)=0$ になるとしても、$P(x) = -\frac{w_{1}w_{2}'' - w_{1}''w_{2}}{W}$ の分子も同時に $0$ になって、極がキャンセルされる（除去可能な特異点になる）可能性はないのか？]
\textbf{A.} 結論から言えば、分子も $0$ になるが、零点の位数が異なるためキャンセルしきれず、必ず極が残る。
局所座標を $u = x-\xi$ とする。2つの解のテイラー展開の最低次を見ると、
\begin{align*}
w_{1}(u) &= c_{0} + c_{1}u + c_{2}u^{2} + \cdots \quad (c_{0} \neq 0) \\
w_{2}(u) &= u^{J}(d_{0} + d_{1}u + \cdots) \quad (d_{0} \neq 0)
\end{align*}
と書ける。分母のロンスキアン $W = w_{1}w_{2}' - w_{1}'w_{2}$ の最低次の項を計算すると、
\[
W(u) \approx c_{0} \cdot (J d_{0} u^{J-1}) - c_{1} \cdot (d_{0} u^{J}) \approx J c_{0} d_{0} u^{J-1}
\]
となり、$W$ は位数 $J-1$ の零点を持つ。
一方、$P(x)$ の分子 $N_{P} = w_{1}w_{2}'' - w_{1}''w_{2}$ の最低次の項を計算すると、
\[
N_{P}(u) \approx c_{0} \cdot J(J-1)d_{0}u^{J-2} - 2c_{2} \cdot d_{0}u^{J} \approx J(J-1)c_{0}d_{0}u^{J-2}
\]
となり、分子は位数 $J-2$ の零点を持つ。
したがって、極の位数は引き算されて、
\[
P(u) = -\frac{N_{P}(u)}{W(u)} \approx -\frac{J(J-1)c_{0}d_{0}u^{J-2}}{J c_{0} d_{0} u^{J-1}} = -\frac{J-1}{u}
\]
となる。$J \ge 2$ のとき $J-1 \ge 1$ であるため、キャンセルしきれずに必ず $u=0$ に $1$ 位の極（留数 $1-J$）が残るのである。

（※ちなみに、このとき $u P(u) \to 1-J$ となることから、決定方程式 $t(t-1) + p_{0}t + q_{0} = 0$ に $p_{0} = 1-J, q_{0} = 0$ を代入すると $t(t-J) = 0$ となり、決定指数が $0$ と $J$ になる事実とも見事に辻褄が合う。）
\end{Q}

\begin{example}
微分方程式
\begin{equation}\label{eq:euler_example}
x^{2}\frac{d^{2}y}{dx^{2}} + (1-2s)x\frac{dy}{dx} + (s^{2}-1)y = 0
\end{equation}
において、$x=0$ は非対数的特異点である。
\end{example}

\begin{proof}
方程式\eqref{eq:euler_example}の両辺を $x^2$ で割ると、$p_{1}(x) = \frac{1-2s}{x}, p_{2}(x) = \frac{s^{2}-1}{x^{2}}$ となる。
$x p_{1}(x) = 1-2s$ および $x^{2} p_{2}(x) = s^{2}-1$ はともに定数（正則）であるため、$x=0$ は確定特異点である。
また、これらは定数関数であるため、展開係数は $q_{0} = 1-2s, r_{0} = s^{2}-1$ であり、$n \ge 1$ に対しては $q_{n} = 0, r_{n} = 0$ となる。

決定方程式は
\[
f(t) = t(t-1) + (1-2s)t + (s^{2}-1) = t^{2} - 2st + s^{2} - 1 = (t-s-1)(t-s+1) = 0
\]
より、決定指数は $t = s+1, s-1$ である。
小さい方を $s_{1} = s-1$、大きい方を $s_{2} = s+1$ とすると、その差は $J = s_{2} - s_{1} = 2$ である。

ここで、小さい方の根 $t=s_{1}$ から出発する漸化式を考える。一般の漸化式において、障害となる項は
\[
R_{n} = \sum_{k=0}^{n-1} \left( (t+k)q_{n-k} + r_{n-k} \right) c_{k}
\]
と書けるが、今 $n \ge 1$ に対して $q_{n} = r_{n} = 0$ であるため、和の中身は常に $0$ となる。
したがって、任意の $n \ge 1$ に対して $R_{n} \equiv 0$ である。
特に $n = J = 2$ においても $R_{2}(c_{0}, c_{1}) = 0$ が自動的に成り立つため、$x=0$ は非対数的特異点である。
\end{proof}

\begin{Q}[【核心】この例題において、$R_n \equiv 0$ となるのは偶然か？]
\textbf{A.} 偶然ではない。この方程式は、各項の次数が揃っている「オイラーの微分方程式（Euler's equidimensional equation）」である。オイラー方程式は $y = x^{\lambda}$ という単項式の代入だけで厳密解が求まるという性質を持つため、そもそも無限級数（フロベニウス級数）で解を構成する必要がない。級数の高次の項が不要であることの代数的な現れが、「$n \ge 1$ で $q_n=r_n=0$ となり、$R_n \equiv 0$ となる」という構造そのものなのである。
\end{Q}
%#endregion

\subsection{不確定特異点}

\subsection{高階フックス型線型常微分方程式}
%#region
$n$階のフックス型微分方程式
\begin{equation}\label{eq:linear-ode-n}
    \frac{d^{n}y}{dx^{n}}+\sum_{j=1}^{n} p_{j}(x)\frac{d^{n-j}y}{dx^{n-j}} = 0   
\end{equation}
について考える。

\begin{definition}
    \eqref{eq:linear-ode-n}の特異点$x=\xi$が\textbf{確定特異点}であるとは、全ての解$\varphi(x)$が$x=\xi$で正則であるか、確定特異点を持つことである
\end{definition}

\begin{proposition}[フックスの定理]\label{prop:fuchs_theorem_n}
    $n$階線型常微分方程式\eqref{eq:linear-ode-n}
    の特異点 $x=\xi$ が確定特異点になる（すなわち任意の解が $x=\xi$ において確定増大度をもつ）ための必要十分条件は、各 $i=1,\dots,n$ について $(x-\xi)^{i}p_{i}(x)$ が $x=\xi$ で正則になることである。
    これは$n$階でも成り立つ。
    （\cite[定理13.4 p.128]{book:takano}）
\end{proposition}

\begin{Q}[$n$階のフックスの定理（命題\ref{prop:fuchs_theorem_n}）の証明は、どのような数学的アイデアに基づいているか？]
\textbf{A.} 十分性（係数の正則性 $\implies$ 確定特異点）は「フロベニウスの方法」による解の直接構成で示し、必要性（確定特異点 $\implies$ 係数の正則性）は「ロンスキアンの増大度評価」を用いて示します。

証明の骨格は以下の通り：
\begin{itemize}
    \item \textbf{十分性の証明（解の構成）：}\\
    $(x-\xi)^{i}p_{i}(x)$ が正則であると仮定する。方程式 \eqref{eq:linear-ode-n} の両辺に $(x-\xi)^n$ をかけると、一般化されたオイラーの微分方程式（$x^n y^{(n)} + \cdots = 0$ の形）の摂動系とみなせる。このとき、解をフロベニウス級数 $y = (x-\xi)^{\rho} \sum_{k=0}^{\infty} c_k (x-\xi)^k$ の形で仮定し、決定方程式（indicial equation）を解くことで、$n$個の線型独立な解（対数項を含む場合もある）を具体的に構成できる。これらはすべて高々多項式オーダーの増大度を持つため、確定特異点となる。

    \item \textbf{必要性の証明（係数の逆算）：}\\
    任意の解の基本系 $y_1, \dots, y_n$ がすべて確定増大度をもつと仮定する。方程式の係数 $p_i(x)$ は、基本系から作られるロンスキアン $W(x) = \det(y_j^{(k-1)})$ を用いてクラメルの公式から逆算できる（例えば $p_1(x) = -W'(x)/W(x)$）。解が確定増大度を持つことから各導関数の増大度も評価でき、結果として得られる $p_i(x)$ の極の位数が「高々 $i$ であること（すなわち $(x-\xi)^i p_i(x)$ が正則であること）」が導かれる。
\end{itemize}
\end{Q}

\begin{Q}[フックスの定理の証明において、十分性（解の構成）に比べて、必要性（ロンスキアンを用いた逆算）の証明が圧倒的に「面倒」なのはなぜか？]
\textbf{A.} 最大の理由は、複素領域特有の「モノドロミー（解の多価性）」の処理と、分母に現れる行列式の「零点の評価」が極めて繊細な作業を要求するから。

それぞれの証明のアプローチの違いから、難易度に以下のような差が生じる：
\begin{itemize}
    \item \textbf{十分性（構成的で楽）：}\\
    係数の形（極の位数）が既に与えられているため、フロベニウス級数（べき級数＋対数項）を代入し、決定方程式と漸化式を機械的に解くだけで済む。計算は長いものの、論理はアルゴリズム的で平坦である。
    
    \item \textbf{必要性（技巧的で面倒）：}\\
    「すべての解が確定増大度を持つ」という解析的な仮定だけから、係数 $p_i(x)$ の極の位数を評価しなければならない。係数はロンスキアンを用いた分数の形で表されますが、特異点近傍で分母が $0$ に近づく際の挙動を正確に抑え込む必要がある。さらに、解が特異点の周りを回ったときの多価性を処理するため、「モノドロミー行列のジョルダン標準形」を用いて、評価しやすい特別な基本系（標準基本系）をわざわざ選び直すという、高度に技巧的なステップを踏むことになる。
\end{itemize}
\end{Q}

一般の $n$ 階線形微分方程式
\begin{equation}\label{eq:n-ode}
\frac{d^{n}y}{dx^{n}} + \sum_{j=1}^{n} q_{j}(x)\frac{d^{n-j}y}{dx^{n-j}} = 0
\end{equation}
について、$x=\infty$ での振る舞いを調べるため、局所座標 $x = \frac{1}{u}$ を導入する。

\begin{proposition}\label{prop:regular_infinity_n}
方程式\eqref{eq:n-ode}に対して、特異点 $x=\infty$ が確定特異点になるための必要十分条件は、各 $j = 1, 2, \dots, n$ について、$x^{j}q_{j}(x)$ が $x=\infty$ において正則になることである。
（すなわち、$x \to \infty$ において $q_{j}(x) = O(x^{-j})$ となることと同値である。）
\end{proposition}

\begin{Q}[【核心】$x=\frac{1}{u}$ と変換したとき、連鎖律によって微分方程式は低階の項が複雑に混ざり合った形になるはずである。それにもかかわらず、なぜ $x^{j}q_{j}(x)$ が正則という、各項が独立したシンプルな条件に帰着するのか？]
\textbf{A.} 変換によって生じる $u$ の冪乗の「最高階部分の寄与」に注目すれば明らかになる。
連鎖律 $\frac{d}{dx} = -u^{2}\frac{d}{du}$ を繰り返し適用すると、$k$ 階微分の最高階部分は
\[
\frac{d^{k}y}{dx^{k}} = (-1)^{k} u^{2k} \frac{d^{k}y}{du^{k}} + (\text{低階の微分項})
\]
となる。これを\eqref{eq:n-ode}に代入すると、全体の最高階項は $(-1)^{n} u^{2n} \frac{d^{n}y}{du^{n}}$ となる。
新しい微分方程式を標準形（最高階の係数を $1$）にするため、全体を $(-1)^{n} u^{2n}$ で割る。このとき、第 $j$ 項（元の $\frac{d^{n-j}y}{dx^{n-j}}$ に対応する部分）の最高階の係数には、
\[
\frac{(-1)^{n-j} u^{2(n-j)}}{(-1)^{n} u^{2n}} \cdot q_{j}\left(\frac{1}{u}\right) = (-1)^{-j} u^{-2j} q_{j}\left(\frac{1}{u}\right)
\]
がかかることになる。
$u=0$ が確定特異点になるためには（フックスの定理より）、この $\frac{d^{n-j}y}{du^{n-j}}$ の係数が $u=0$ で高々 $j$ 位の極（つまり $u^{-j}$ の特異性）に収まらなければならない。
すなわち、$u^{-2j} q_{j}(\frac{1}{u}) = O(u^{-j})$ が要求される。両辺に $u^{2j}$ を掛ければ $q_{j}(\frac{1}{u}) = O(u^{j})$ となり、これを元の $x$ に戻せば $q_{j}(x) = O(x^{-j})$ 、つまり $x^{j}q_{j}(x)$ が $x=\infty$ で正則であることが導かれる。
連鎖律による低階の複雑な混ざり合いは、常にこの最も厳しい条件（最高階からの寄与）に吸収されるため、結果として各係数ごとの美しい条件式だけが残るのである。
\end{Q}
%#endregion
%#region
$n$ 階線形微分方程式\eqref{eq:linear-ode-n}において、特異点 $x=\xi$ の周りの局所座標を $u = x-\xi$ とする。
方程式の両辺に $u^{n}$ をかけると、
\[
u^{n}\frac{d^{n}y}{du^{n}} + \sum_{j=1}^{n} u^{n-j} \left( u^{j}p_{j}(u+\xi) \right) \frac{d^{n-j}y}{du^{n-j}} = 0
\]
となる。ここで、$x=\xi$ が確定特異点であるという仮定から、$u^{j}p_{j}(u+\xi)$ は $u=0$ で正則となるため、これを $q_{j}(u)$ と置き直す。

解をフロベニウス級数 $\varphi(u) = u^{s} \sum_{k=0}^{\infty} c_{k} u^{k} \quad (c_{0} \neq 0)$ と仮定し、上式に代入して最低次 $u^{s}$ の係数を比較する。
微分演算子について $u^{k} \frac{d^{k}}{du^{k}} (u^{s}) = s(s-1)\cdots(s-k+1)u^{s}$ となる性質を用いると、$u^{s}$ の係数が消えるための条件として次の方程式が得られる。

\begin{equation}\label{eq:indicial-n}
f(s) = s(s-1)\cdots(s-n+1) + \sum_{j=1}^{n-1} q_{j}(0) s(s-1)\cdots(s-n+j+1) + q_{n}(0) = 0
\end{equation}
（ただし、$q_{j}(0)$ は $(x-\xi)^{j}p_{j}(x)$ の $x=\xi$ における極限値である。）

\begin{definition}
\eqref{eq:linear-ode-n}に対し、多項式 $f(s)$ を特異点 $x=\xi$ における\textbf{決定多項式}（または特性多項式）という。さらに、$f(s)=0$ を $x=\xi$ における\textbf{決定方程式}（または特性方程式）と呼び、その根 $s$ を\textbf{決定指数}（または特性指数）という。
\end{definition}

\begin{Q}[【核心】決定方程式 $f(s)=0$ の各項が、$s^{n}$ のような単純な冪乗ではなく、$s(s-1)\cdots(s-n+1)$ のような「下降階乗冪（Falling factorial）」の形をとるのはなぜか？]
\textbf{A.} これは、特異点近傍での主要な振る舞いが「指数関数 $e^{sx}$」ではなく「冪関数 $u^{s}$」によって支配されているからである。
定数係数の線形微分方程式では解を $e^{sx}$ と仮定するため、微分するたびに $s$ が前に出て $s^{n}$ が現れる。しかし確定特異点においては、方程式が本質的にオイラーの微分方程式（各項が $x^{k} y^{(k)}$ の形を持つ）に帰着される。冪関数 $u^{s}$ を $k$ 回微分して $u^{k}$ を掛けると、$s^{k}$ ではなく $s(s-1)\cdots(s-k+1)$ という係数が飛び出してくる。この微分の代数的な構造が、決定方程式に下降階乗冪を強制しているのである。
\end{Q}

決定方程式の根（決定指数）の一つを $s$ とする。
解が $u=0$ の周りで
\[
\varphi(u) = u^{s}\sum_{k=0}^{\infty} c_{k}u^{k} \quad (c_{0}\neq 0)
\]
のように展開されると仮定し、方程式に代入して $u^{s+k}$ の係数を整理すると、一般に次のような漸化式が得られる。
\begin{equation}\label{eq:recurrence-n}
f(s+k)c_{k} + R_{k}(c_{0}, \cdots, c_{k-1}) = 0 \quad (k \ge 1)
\end{equation}
したがって、$f(s+k) = 0$ となる正の整数 $k$ が存在しない限り、$c_{k}$ は $c_{0}$ から順番に矛盾なく一意に決定できる。

次に、この特異点 $x=\xi$ における $n$ 個の決定指数 $s_{1}, s_{2}, \dots, s_{n}$ の和について考察する。
方程式\eqref{eq:indicial-n}の決定多項式 $f(s)$ に対して「解と係数の関係」を適用することで、以下の美しい関係が導かれる。
\begin{equation}\label{eq:sum-of-exponents}
\sum_{j=1}^{n} s_{j} = \frac{1}{2}n(n-1) - q_{1}(0)
\end{equation}

\begin{Q}[【核心】決定多項式 $f(s)$ から、なぜこの特性指数の和の式が導出されるのか？]
\textbf{A.} $n$ 次方程式 $f(s)=0$ において、根の和 $\sum s_{j}$ は「$s^{n-1}$ の係数にマイナスをつけたもの」となる（解と係数の関係）。
決定多項式\eqref{eq:indicial-n}の各項の展開を考えると、
\begin{align*}
f(s) &= s(s-1)\cdots(s-n+1) + q_{1}(0)s(s-1)\cdots(s-n+2) + \cdots \\
&= \left( s^{n} - (0+1+2+\cdots+n-1)s^{n-1} + \cdots \right) + q_{1}(0)\left( s^{n-1} + \cdots \right) + \cdots \\
&= s^{n} + \left( -\frac{1}{2}n(n-1) + q_{1}(0) \right)s^{n-1} + (\text{低次の項})
\end{align*}
となる。したがって、$s^{n-1}$ の係数の符号を反転させることで、ただちに $\sum s_{j} = \frac{1}{2}n(n-1) - q_{1}(0)$ が得られる。
\end{Q}

\begin{Q}[【核心】なぜわざわざ、特異点における特性指数の「和」を計算したのか？これにはどのような大域的な意味があるのか？]
\textbf{A.} この局所的な和こそが、リーマン球面 $\hat{\mathbb{C}}$ 全体における極めて強力な拘束条件「フックスの関係式（Fuchs' relation）」を導くためのパーツだからである。
各特異点での $q_{1}(0)$ は、微分方程式の第2項の係数 $p_{1}(x)$ のその特異点における留数（Residue）に相当する。複素解析の留数定理により、リーマン球面上のすべての特異点（$\infty$ を含む）の留数を足し合わせると必ず $0$ になる。
つまり、すべての特異点について\eqref{eq:sum-of-exponents}を足し合わせると $q_{1}(0)$ の寄与が完全にキャンセルされ、「全特異点の特性指数の総和は、空間のトポロジー（特異点の数と階数 $n$）だけで決まる一定の整数になる」という大域的な定理が成立する。パンルヴェ方程式などのモノドロミー保存変形を考える際、この関係式は系の自由度を落とす絶対的な土台として機能する。
\end{Q}
%#endregion
%#region
ここで
\[
f(u;s)=s(s-1)\cdots(s-n+1) + \sum_{j=1}^{n-1} q_{j}(u) s(s-1)\cdots(s-n+j+1) + q_{n}(u)
\]
とおく。
\begin{proposition}
    \eqref{eq:linear-ode-n}は
    \[
    f(u;D)y=0\quad D=u\frac{d}{du}
    \]
    と表せる。
\end{proposition}

\begin{proof}
補題\ref{lem:falling_factorial}から従う。
\end{proof}

\begin{lemma}\label{lem:falling_factorial}
    \[
    u^{k}\frac{d^{k}}{du^{k}}=D(D-1)(D-2)\cdots(D-k+1)
    \]
\end{lemma}

\begin{proof}
$k=1$のとき、$u\frac{d}{du} = D$ であるため、成立する。
$k>1$のとき、帰納法を用いる。$u^{k}\frac{d^{k}}{du^{k}} = u \cdot u^{k-1} \frac{d^{k-1}}{du^{k-1}} \cdot \frac{d}{du}$ と書ける。帰納法の仮定から、$u^{k-1} \frac{d^{k-1}}{du^{k-1}} = D(D-1)\cdots(D-k+2)$ であるため、
\[
u^{k}\frac{d^{k}}{du^{k}} = u \cdot D(D-1)\cdots(D-k+2) \cdot \frac{d}{du}
\]
となる。ここで、$(D-a)\frac{d}{du}=\frac{d}{du}(D-a-1)$ という関係を用いると、
\[
u^{k}\frac{d^{k}}{du^{k}} = u \cdot D(D-1)\cdots(D-k+2) \cdot \frac{d}{du} = u\frac{d}{du}(D-1)\cdots(D-k+2)(D-k+1) = D(D-1)(D-2)\cdots(D-k+1)
\]
となり、主張が成立する。
\end{proof}

\begin{definition}[フロベニウスの方法]
このようにして、冪級数解を求める方法を\textbf{フロベニウスの方法}（Frobenius method）と呼ぶ。
\end{definition}

無限遠点 $x=\infty$ のまわりでの微分方程式の表現について、オイラー作用素（Euler operator）を用いてより代数的に洗練された表示を考える。
方程式\eqref{eq:linear-ode-n}に $x^{n}$ を掛け、$q_{j}\left(\frac{1}{x}\right) = x^{j}p_{j}(x)$ とおくと、
\begin{equation}\label{eq:ode-euler-x}
x^{n}\frac{d^{n}y}{dx^{n}} + \sum_{j=1}^{n} q_{j}\left(\frac{1}{x}\right) x^{n-j}\frac{d^{n-j}y}{dx^{n-j}} = 0
\end{equation}
となる。ここで補題\ref{lem:falling_factorial}より、微分演算子 $D_{x} = x\frac{d}{dx}$ を用いると $x^{k}\frac{d^{k}}{dx^{k}} = D_{x}(D_{x}-1)\cdots(D_{x}-k+1)$ と表せるため、上式は次のような演算子の形で書ける。
\[
\left[ D_{x}(D_{x}-1)\cdots(D_{x}-n+1) + \sum_{j=1}^{n}q_{j}\left(\frac{1}{x}\right) D_{x}(D_{x}-1)\cdots(D_{x}-n+j+1) \right] y = 0
\]

ここで局所座標 $u = \frac{1}{x}$ を導入する。連鎖律より $\frac{d}{dx} = -u^{2}\frac{d}{du}$ であるから、オイラー作用素は
\[
D_{x} = x\frac{d}{dx} = \frac{1}{u}\left(-u^{2}\frac{d}{du}\right) = -u\frac{d}{du} = -D_{u}
\]
と極めてシンプルに変換される。ここで、2変数多項式 $f(u; t)$ を
\[
f(u; t) = t(t-1)\cdots(t-n+1) + \sum_{j=1}^{n} q_{j}(u) t(t-1)\cdots(t-n+j+1)
\]
と定義すれば、以下の命題が成り立つ。

\begin{proposition}
\eqref{eq:linear-ode-n}は、$x=\infty$ （すなわち $u=0$）のまわりでは、オイラー作用素 $D_{u} = u\frac{d}{du}$ を用いて
\[
f(u; -D_{u})y = 0
\]
と簡潔に表現できる。
\end{proposition}

\begin{Q}[【核心】$x=\frac{1}{u}$ という変換で生じるはずの「連鎖律による複雑な低階項の混ざり合い」は、どこへ消えたのか？]
\textbf{A.} オイラー作用素 $D = x\frac{d}{dx}$ の持つ「スケール不変性」に吸収された。
通常の微分演算子 $\frac{d}{dx}$ は座標変換に対して複雑に振る舞うが、オイラー作用素は $x \to x^{a}$ という変換に対して $D_{x} \to \frac{1}{a} D_{u}$ と振る舞う。特に $x = u^{-1}$ の反転変換に対しては単に符号が反転するだけ（$D_{x} = -D_{u}$）であり、微分の階数が混ざることは一切ない。方程式全体をあらかじめ $D_{x}$ の多項式として組み立てておくことで、座標反転の代数的な複雑さを「$-1$ を代入するだけ」という操作に還元できたのである。
\end{Q}

作用素 $f(u; -D_{u})$ を展開し、フックスの定理を適用できる標準形
\[
u^{n}\frac{d^{n}y}{du^{n}} + \sum_{k=1}^{n} r_{k}(u) u^{n-k}\frac{d^{n-k}y}{du^{n-k}} = 0
\]
に書き直すことを考える。
ここで重要なのは、作用素 $(-D_{u})(-D_{u}-1)\cdots(-D_{u}-m+1)$ の最高階の微分が $(-1)^{m} u^{m} \frac{d^{m}}{du^{m}}$ となり、それ以下の階数の微分は定数係数の線形結合になるという事実である。

\begin{proposition}\label{prop:infinity_fuchs_q}
\eqref{eq:linear-ode-n}に対して、特異点 $x=\infty$ が確定特異点になるための必要十分条件は、各 $j = 1, 2, \dots, n$ について $q_{j}(u)$ が $u=0$ において正則になることである。
またこのとき、$x=\infty$ における特性指数は、決定方程式 $f(0; -s) = 0$ の根 $s$ として得られる。
\end{proposition}

\begin{Q}[【核心】作用素を展開した後の係数 $r_{k}(u)$ の正則性と、元の $q_{j}(u)$ の正則性が完全に「同値」になるのはなぜか？]
\textbf{A.} 両者の関係が、「最高階の微分から順番に決まる、三角行列の構造（帰納的構造）」を持っているからである。
作用素 $f(u; -D_{u})$ における第 $j$ 項 $q_{j}(u) (-D_{u})\cdots(-D_{u}-n+j+1)$ は、最大で $n-j$ 階の微分までしか持たない。
方程式全体を最高階の係数 $(-1)^{n}$ で割って標準形にしたとき、$(n-k)$ 階微分の係数 $r_{k}(u)$ を構成するのは以下の3つのみである。
\begin{enumerate}
    \item 第 $k$ 項から出る最高階の寄与： $(-1)^{-k} q_{k}(u)$
    \item 第 $j < k$ 項から出る低階の寄与： $q_{1}(u), \dots, q_{k-1}(u)$ の定数倍
    \item 先頭の作用素（$q$ が付かない項）から出る低階の寄与： 単なる定数
\end{enumerate}
つまり、$r_{k}(u) = (-1)^{-k} q_{k}(u) + \sum_{j=1}^{k-1} c_{k,j} q_{j}(u) + c_{k,0}$ という形になる。
\begin{itemize}
    \item $k=1$ のとき： $r_{1}(u) = -q_{1}(u) + c_{1,0}$ より、$r_{1}$ が正則 $\iff$ $q_{1}$ が正則。
    \item $k=2$ のとき： $r_{2}(u) = q_{2}(u) + c_{2,1}q_{1}(u) + c_{2,0}$ となり、$q_{1}$ はすでに正則であるため、$r_{2}$ が正則 $\iff$ $q_{2}$ が正則。
\end{itemize}
このように、上から帰納的に連鎖していくため、すべての $r_{k}(u)$ が正則であること（確定特異点であること）と、すべての $q_{j}(u)$ が正則であることは、代数的に完全に同値となるのである。
\end{Q}

\begin{Q}[【核心】決定方程式が $f(0; s)=0$ ではなく $f(0; -s)=0$ となるのはなぜか。また、この根 $s$ は「$x=\infty$ における特性指数」とどう対応するのか？]
\textbf{A.} 演算子が $D_{x} = -D_{u}$ と反転した結果、固有値も反転するからである。
解を $u$ のフロベニウス級数 $y \sim u^{s}$ と仮定したとき、$D_{u} (u^{s}) = s u^{s}$ であるため、作用素 $f(u; -D_{u})$ を適用した最低次項は $f(0; -s) = 0$ となる。
一方、$x=\infty$ における特性指数は通常、$x$ の級数 $y \sim x^{\rho}$ として定義される。$x^{\rho} = (1/u)^{\rho} = u^{-\rho}$ であるため、$s = -\rho$ の関係がある。
つまり、特性指数 $\rho$ を直接求める決定方程式は $f(0; \rho) = 0$ となり、有限の特異点と全く同じ形の多項式 $f$ に $\rho$ を代入して解けばよい、という美しい結論が得られる。
\end{Q}

\begin{example}
ガウスの超幾何微分方程式\eqref{eq:gauss_hypergeometric}において、局所座標 $u=\frac{1}{x}$ とオイラー作用素 $D_{u} = u\frac{d}{du}$ を用いると、無限遠点 $x=\infty$ の周りでの方程式は
\[
f(u; -D_{u})y = 0, \quad f(u; -D_{u}) = D_{u}(D_{u}+1) - \frac{cu-(a+b+1)}{u-1}D_{u} - \frac{ab}{u-1}
\]
と表される。このとき、$x=\infty$ での特性指数は $s=a, b$ となる。
\end{example}

\begin{proof}
方程式\eqref{eq:gauss_hypergeometric}を、オイラー作用素 $D_{x} = x\frac{d}{dx}$ （$x^{2}\frac{d^{2}}{dx^{2}} = D_{x}(D_{x}-1)$ および $x\frac{d}{dx} = D_{x}$）を用いて書き換えるため、全体を $x$ で割ると、
\[
(1-x) \frac{1}{x} D_{x}(D_{x}-1)y + \left( \frac{c}{x} - (a+b+1) \right) D_{x}y - aby = 0
\]
となる。したがって、演算子 $f(u; D_{x})$ は（$u=1/x$ として）
\[
f(u; D_{x}) = (u-1)D_{x}(D_{x}-1) + (cu-(a+b+1))D_{x} - ab
\]
と書ける。ここに $D_{x} = -D_{u}$ を代入すると、
\begin{align*}
f(u; -D_{u}) &= (u-1)(-D_{u})(-D_{u}-1) + (cu-(a+b+1))(-D_{u}) - ab \\
&= (u-1)D_{u}(D_{u}+1) - (cu-(a+b+1))D_{u} - ab
\end{align*}
全体を最高階の係数 $(u-1)$ で割って標準形にすると、
\[
D_{u}(D_{u}+1) - \frac{cu-(a+b+1)}{u-1}D_{u} - \frac{ab}{u-1}
\]
となる。
決定方程式は $u=0$ とし、作用素 $D_{u}$ を $s$ に置き換える（$f(0; -s) = 0$）ことで得られる。
\begin{align*}
s(s+1) - \frac{-(a+b+1)}{-1}s - \frac{ab}{-1} &= s(s+1) - (a+b+1)s + ab \\
&= s^{2} - (a+b)s + ab \\
&= (s-a)(s-b) = 0
\end{align*}
したがって、$x=\infty$ における特性指数は $s=a, b$ である。
\end{proof}

\begin{Q}[【核心】ガウスの超幾何関数 $F(a,b,c; x)$ のパラメータ $a, b, c$ は、一見すると不規則に並んでいるように見えるが、これらは微分方程式の特異点とどのような関係にあるのか？]
\textbf{A.} パラメータ $a, b, c$ は、まさに「各特異点における特性指数」そのものを指定するための変数である。
ガウスの微分方程式は $x=0, 1, \infty$ の3点のみを確定特異点に持つ。それぞれの特性指数を計算すると、
\begin{itemize}
    \item $x=0$ での特性指数： $0, 1-c$
    \item $x=1$ での特性指数： $0, c-a-b$
    \item $x=\infty$ での特性指数： $a, b$
\end{itemize}
となる。つまり、$\infty$ での特性指数がそのまま関数の基本パラメータ $a, b$ として採用されているのである。
（※ベルンハルト・リーマンは、この「3つの特異点とそれぞれの特性指数」を指定するだけで、微分方程式全体が大域的に一意に定まってしまうことを証明した。超幾何関数とは、パラメータの代数的な操作の産物ではなく、「リーマン球面上に3つの穴を開け、そこでの局所的な振る舞い（特性指数）を指定することで定義される幾何学的な関数」なのである。）
\end{Q}
したがって$x=\infty$における特性指数$s_{1}^{(\infty)},\cdots,s_{n}^{(\infty)}$に対して、解と係数の関係より
\begin{equation}\label{eq:sum-of-infinity-exponents}
\sum_{j=1}^{n} s_{j}^{(\infty)} = -\frac{1}{2}n(n-1) + q_{1}(0)
\end{equation}
が成り立つ。
%#endregion
%#region
ところで、有限の確定特異点 $x=\xi_{j}$ （$j=1, \dots, m$）の条件から、微分方程式の第2項の係数 $p_{1}(x)$ は有理関数であり、各特異点での留数を $a_{j}$ とすると、部分分数分解により
\[
p_{1}(x) = \sum_{j=1}^{m} \frac{a_{j}}{x-\xi_{j}} + (x\text{の多項式})
\]
と表せる。

\begin{Q}[【核心】有理関数 $p_{1}(x)$ に含まれるはずの「$x$の多項式」の部分は、なぜ完全に $0$ になると断言できるのか？]
\textbf{A.} 無限遠点 $x=\infty$ も確定特異点であるという条件が、多項式部分を禁止するからである。
命題\ref{prop:infinity_fuchs_q}より、$x=\infty$ が確定特異点であるためには $x p_{1}(x)$ が $x=\infty$ で正則、すなわち $x \to \infty$ において $p_{1}(x) = O(x^{-1})$ とならなければならない。もし $p_{1}(x)$ に定数項や $x$ の1次以上の項が含まれていれば、$x \to \infty$ で発散または $O(1)$ となり、この条件に反してしまう。したがって、多項式部分は恒等的に $0$ でなければならない。（※これはリュービルの定理の帰結でもある）。
\end{Q}

したがって、$p_{1}(x)$ は厳密に
\[
p_{1}(x) = \sum_{j=1}^{m} \frac{a_{j}}{x-\xi_{j}}
\]
となる。
各有限特異点 $x=\xi_{l}$ の局所座標 $u = x-\xi_{l}$ において、$q_{1}^{(l)}(0) = \lim_{x \to \xi_{l}} (x-\xi_{l})p_{1}(x) = a_{l}$ である。よって、$x=\xi_{l}$ における $n$ 個の特性指数 $s_{1}^{(l)}, \dots, s_{n}^{(l)}$ に対して、\eqref{eq:sum-of-exponents}より
\[
\sum_{j=1}^{n} s_{j}^{(l)} = \frac{1}{2}n(n-1) - a_{l}
\]
が成り立つ。
一方、$x=\infty$ における特性指数の和を求める。局所座標 $u=1/x$ を用いると、
\[
q_{1}^{(\infty)}(0) = \left. \frac{1}{u}p_{1}\left(\frac{1}{u}\right) \right|_{u=0} = \left. \sum_{j=1}^{m} a_{j}\frac{1}{1-u\xi_{j}} \right|_{u=0} = \sum_{j=1}^{m} a_{j}
\]
となる。

\begin{Q}[【核心】有限特異点での特性指数の和は $\frac{1}{2}n(n-1) - a_{l}$ であったのに対し、無限遠点での和はなぜ符号が反転して $-\frac{1}{2}n(n-1) + \sum a_{j}$ となるのか？]
\textbf{A.} 無限遠点における決定方程式が $f(0; -s) = 0$ であり、$s$ ではなく $-s$ が代入されることで「上昇階乗冪」が現れるからである。
有限特異点の決定方程式には下降階乗冪 $s(s-1)\cdots(s-n+1) = s^{n} - \frac{1}{2}n(n-1)s^{n-1} + \cdots$ が現れた。
一方、無限遠点では $(-s)(-s-1)\cdots(-s-n+1) = (-1)^{n} s(s+1)\cdots(s+n-1)$ となり、これを展開すると $(-1)^{n} (s^{n} + \frac{1}{2}n(n-1)s^{n-1} + \cdots)$ という上昇階乗冪になる。
全体を $(-1)^{n}$ で割って最高次の係数を $1$ にすると、$s^{n-1}$ の係数が $-\frac{1}{2}n(n-1)$ から $+\frac{1}{2}n(n-1)$ に反転していることがわかる。したがって、解と係数の関係を用いた際にも和の符号が反転し、$\sum s^{(\infty)} = -\frac{1}{2}n(n-1) + q_{1}^{(\infty)}(0)$ となるのである。
\end{Q}

これより、$x=\infty$ における特性指数の和は
\[
\sum_{j=1}^{n}s_{j}^{(\infty)} = -\frac{1}{2}n(n-1) + \sum_{j=1}^{m} a_{j}
\]
となる。
有限特異点と無限遠点のすべての特性指数の和をとると、$a_{j}$ の項が見事に相殺され、以下の極めて重要な大域的関係式が成立する。

\begin{proposition}\label{prop:fuchs_relation}
リーマン球面上の特異点がすべて確定特異点であるようなフックス型微分方程式\eqref{eq:linear-ode-n}の特性指数には、
\begin{align*}
\sum_{l=1}^{m}\sum_{j=1}^{n} s_{j}^{(l)} + \sum_{j=1}^{n} s_{j}^{(\infty)} &= \sum_{l=1}^{m} \left( \frac{1}{2}n(n-1) - a_{l} \right) + \left( -\frac{1}{2}n(n-1) + \sum_{j=1}^{m} a_{j} \right) \\
&= \frac{1}{2}n(n-1)(m-1)
\end{align*}
という関係式が成立する。この方程式のパラメータを大域的に縛る拘束条件を\textbf{フックスの関係式（Fuchs' relation）}という。
\end{proposition}
%#endregion
%#region
微分方程式を完全に決定するためには、その係数に含まれる独立なパラメータの数を数え上げればよい。
$n$ 階フックス型微分方程式において、確定特異点の位置を $\Xi=\{\xi_{1}, \xi_{2}, \dots, \xi_{m}, \infty\}$ として考える。

まず、各特異点における局所的な振る舞いを決定する「特性指数」の自由度について考える。
特異点は全部で $m+1$ 個あり、各点に $n$ 個の特性指数が存在するため、特性指数の総数は $(m+1)n$ 個である。しかし、これらが満たすべき大域的な制約は命題\ref{prop:fuchs_relation}（フックスの関係式）の1つのみであるため、特性指数のうち自由に選ぶことができるパラメータの数は
\begin{equation}
E_{0}(n,m) = (m+1)n - 1
\end{equation}
個となる。

次に、微分方程式の係数 $p_{i}(x)$ 全体に含まれるパラメータの総数を数え上げる。
命題\ref{prop:fuchs_theorem_n}から、有限の特異点 $x=\xi_{l}$ における極の最大位数は $i$ であるため、$p_{i}(x)$ は $mi$ 個の未知定数 $a_{k}^{(l)}$ （$k=1,\dots,i$, $l=1,\dots,m$）を用いて、部分分数分解により
\[
p_{i}(x) = \sum_{l=1}^{m}\sum_{k=1}^{i} \frac{a_{k}^{(l)}}{(x-\xi_{l})^{k}} + (x\text{の多項式})
\]
と表される。
さらに命題\ref{prop:infinity_fuchs_q}より、無限遠点 $x=\infty$ が確定特異点であるためには $x^{i}p_{i}(x)$ が $\infty$ で正則でなければならない。これにより、$p_{i}(x)$ は $x \to \infty$ で $O(x^{-i})$ となることが要求されるため、多項式部分は $0$ となり、
\begin{equation}\label{eq:pi_fraction}
p_{i}(x) = \sum_{l=1}^{m}\sum_{k=1}^{i} \frac{a_{k}^{(l)}}{(x-\xi_{l})^{k}}
\end{equation}
に制限される。さらに、$p_{i}(x) = O(x^{-i})$ という条件から、$a_{k}^{(l)}$ たちの間に一定の拘束条件が生じる。

\begin{Q}[【核心】$p_{i}(x) = O(x^{-i})$ という無限遠点での条件は、具体的にいくつの独立なパラメータを奪う（束縛する）のか？]
\textbf{A.} 厳密に $i-1$ 個の独立な自由度を奪う。
局所座標 $u=1/x$ を用いて\eqref{eq:pi_fraction}を $u=0$ のまわりでテイラー展開する。
\[
\frac{1}{(x-\xi_{l})^{k}} = \frac{u^{k}}{(1-\xi_{l}u)^{k}} = u^{k} \left( 1 + k\xi_{l}u + \frac{k(k+1)}{2}\xi_{l}^{2}u^{2} + \cdots \right)
\]
これを $p_{i}(x)$ に代入して $u$ のべき級数として整理したとき、$u^{i}$ 以上のオーダー（$O(u^{i})$）にならなければならない。つまり、展開した際の $u^{1}, u^{2}, \dots, u^{i-1}$ の係数がすべて $0$ になる必要がある。
これは未知定数 $a_{k}^{(l)}$ に対する $i-1$ 個の線形連立方程式となる。特異点の位置 $\xi_{l}$ はすべて相異なるため、この連立方程式の係数行列（ヴァンデルモンド行列の微分型）はフルランクとなる。したがって、これら $i-1$ 個の条件式は完全に独立であり、冗長なものは一つもない。
\end{Q}

これにより、$p_{i}(x)$ に含まれる独立なパラメータの数は、元の $mi$ 個から独立な制約 $i-1$ 個を引いて
\[
mi - (i-1) = m i - i + 1
\]
個となる。
これを $i=1, 2, \dots, n$ について合計すれば、方程式全体に含まれるパラメータの総数が得られる。
\begin{align*}
\sum_{i=1}^{n} (mi - i + 1) &= (m-1)\sum_{i=1}^{n} i + \sum_{i=1}^{n} 1 \\
&= (m-1)\frac{1}{2}n(n+1) + n \\
&= \frac{1}{2}n \big( m(n+1) - (n-1) \big)
\end{align*}

以上の議論から、フックス型微分方程式のモジュライ空間の次元（パラメータ数）に関する以下の重要な命題を得る。

\begin{proposition}\label{prop:fuchs_parameters}
\eqref{eq:linear-ode-n}が $\Xi=\{\xi_{1}, \xi_{2}, \dots, \xi_{m}, \infty\}$ に確定特異点をもつフックス型微分方程式である時、この微分方程式の係数全体は、
\[
E(n,m) = \frac{1}{2}n \big( m(n+1) - (n-1) \big)
\]
個の独立なパラメータを含む。
\end{proposition}

命題\ref{prop:fuchs_parameters}より、方程式の係数に含まれる全パラメータ数は $E(n,m)$ である。一方、特性指数として自由に選べるパラメータ数は $E_{0}(n,m)$ であった。
これらの差を計算すると、
\begin{align*}
E(n,m) - E_{0}(n,m) &= \frac{1}{2}n \big( m(n+1) - (n-1) \big) - \big( (m+1)n - 1 \big) \\
&= \frac{1}{2}(n-1)(mn - n - 2)
\end{align*}
となる。$m \ge 2, n \ge 2$ のとき
\[
E(n,m) - E_{0}(n,m) \ge 0
\]
であるため、各特異点での特性指数をすべて指定したとしても、一般にはそれだけでは微分方程式を完全に決定するには至らないことがわかる。

\begin{definition}
このように、微分方程式の係数に含まれるパラメータのうち、特性指数の値に関係しない（特性指数を固定してもなお決定されない）独立なパラメータのことを\textbf{アクセサリー・パラメータ}（accessory parameter）という。その総数は $N_{acc} = \frac{1}{2}(n-1)(mn - n - 2)$ で与えられる。
\end{definition}

\begin{example}\label{ex:acc_parameter_zero}
2階フックス型微分方程式で、確定特異点が3つ（有限に $m=2$ 個、および無限遠点 $\infty$）の場合について考える。
このとき、$E(2,2) = 5$、$E_{0}(2,2) = 5$ であるため、$N_{acc} = 5 - 5 = 0$ となる。
したがって、この場合にはアクセサリー・パラメータは存在しない。
\end{example}

\begin{Q}[【核心】特異点が3つの場合（$m=2$）にアクセサリー・パラメータが $0$ になることは、解析学的に何を意味しているのか？]
\textbf{A.} 「特異点の位置」と「各特異点での特性指数（局所的な振る舞い）」を指定するだけで、微分方程式全体が大域的にただ1つに決定されてしまう（剛性を持つ）ことを意味している。
特異点3つの2階方程式といえば「ガウスの超幾何微分方程式」である。超幾何関数が持つ数々の美しい変換則や積分表示は、この「アクセサリー・パラメータを持たない（特性指数だけで方程式が一意にロックされる）」という強烈な代数的性質の恩恵に他ならない。
\end{Q}

\begin{Q}[【核心】特異点を1つ増やして4つ（$m=3$）にした場合、何が起きるのか？また、それがパンルヴェ方程式とどう繋がるのか？]
\textbf{A.} $n=2, m=3$ のとき、$N_{acc} = \frac{1}{2}(1)(6 - 2 - 2) = 1$ となり、アクセサリー・パラメータが初めて1つ出現する。これが「ホイン（Heun）の微分方程式」である。
特性指数をすべて固定しても、未知のパラメータが1つ残るため、解のモノドロミー群はこのパラメータに依存して変化してしまう。
岡本和夫の理論（モノドロミー保存変形）は、特異点の位置 $t$ を動かしたときに「モノドロミー群が変化しないように」このアクセサリー・パラメータをどう時間発展させればよいか、という非線形力学系を構築することである。その際、岡本はこの1次元のパラメータを直接扱うのではなく、あえて「見かけの特異点 $q$」と「その留数に由来する運動量 $p$」という新しい変数を導入することで、この変形方程式を美しいハミルトン系（パンルヴェVI型方程式）として見事に定式化したのである。
\end{Q}
%#endregion
%#region
ここで、\eqref{eq:linear-ode-n}の各特異点における特性指数を一覧として書き並べ、方程式の大域的な構造を視覚的に捉える記法を導入する。

\begin{definition}[リーマン図式]
    確定特異点 $x = \xi_1, \xi_2, \dots, \xi_m, \infty$ を持つフックス型微分方程式\eqref{eq:linear-ode-n}に対し、各特異点での特性指数を列挙した表
    \[
    \left\{
    \begin{matrix}
        x = \xi_1 & x = \xi_2 & \cdots & x = \xi_m & x = \infty \\
        s_1^{(1)} & s_1^{(2)} & \cdots & s_1^{(m)} & s_1^{(\infty)} \\
        s_2^{(1)} & s_2^{(2)} & \cdots & s_2^{(m)} & s_2^{(\infty)} \\
        \vdots    & \vdots    &        & \vdots    & \vdots \\
        s_n^{(1)} & s_n^{(2)} & \cdots & s_n^{(m)} & s_n^{(\infty)}
    \end{matrix}
    \right\}
    \]
    を、その微分方程式の\textbf{リーマン図式}（Riemann scheme）という。（※記号 $P$ を冠してリーマンのP関数と呼ばれることもある。）
\end{definition}

\begin{example}
    ガウスの超幾何微分方程式\eqref{eq:gauss_hypergeometric}のリーマン図式は、以下のようになる。
    \[
    \left\{
    \begin{matrix}
        x = 0 & x = 1 & x = \infty \\
        0 & 0 & a \\
        1-c & c-a-b & b \\
    \end{matrix}
    \right\}
    \]
\end{example}

\begin{proof}
    ガウスの超幾何微分方程式\eqref{eq:gauss_hypergeometric}は、$x=0, 1, \infty$ の3点に確定特異点を持つフックス型微分方程式である。
    それぞれの特異点における特性指数は、これまでの計算により
    \begin{itemize}
        \item $x=0$ での特性指数： $0, 1-c$
        \item $x=\infty$ での特性指数： $a, b$
    \end{itemize}
    であることが既にわかっている。残る $x=1$ での特性指数を求める。
    局所座標を $u = x-1$ とすると $x = u+1$ であり、$x(1-x) = (u+1)(-u) = -u(u+1)$ となる。方程式全体を $-u(u+1)$ で割って標準形にすると、
    \[
    \frac{d^{2}y}{du^{2}} - \frac{c - (a+b+1)(u+1)}{u(u+1)} \frac{dy}{du} + \frac{ab}{u(u+1)} y = 0
    \]
    となる。フックスの定理の条件に従い、$q_{1}(u), q_{2}(u)$ を計算して $u=0$ を代入する。
    \begin{align*}
        q_{1}(0) &= \left. u \left( - \frac{c - (a+b+1)(u+1)}{u(u+1)} \right) \right|_{u=0} = \left. - \frac{c - (a+b+1)(u+1)}{u+1} \right|_{u=0} = -c + a + b + 1 \\
        q_{2}(0) &= \left. u^{2} \left( \frac{ab}{u(u+1)} \right) \right|_{u=0} = \left. u \frac{ab}{u+1} \right|_{u=0} = 0
    \end{align*}
    したがって、$x=1$（すなわち $u=0$）における決定方程式は $s(s-1) + q_{1}(0)s + q_{2}(0) = 0$ より、
    \[
    s(s-1) + (a+b+1-c)s = s^{2} + (a+b-c)s = s(s - (c-a-b)) = 0
    \]
    となり、特性指数は $0, c-a-b$ であることが示された。
\end{proof}

\begin{Q}[【核心】この超幾何微分方程式のリーマン図式に対して、先ほど証明した「フックスの関係式」は本当に成り立っているか？]
\textbf{A.} 完璧に成り立っている。
特異点の数は $m+1 = 3$（有限が $m=2$）、階数は $n=2$ である。
フックスの関係式（命題\ref{prop:fuchs_relation}）の右辺は、$\frac{1}{2}n(n-1)(m-1) = \frac{1}{2}(2)(1)(1) = 1$ となる。
一方、左辺（リーマン図式に現れる全6個の特性指数の総和）を計算すると、
\[
0 + (1-c) + 0 + (c-a-b) + a + b = 1
\]
となり、パラメータ $a, b, c$ の値にいかなる依存もせず、見事に両辺が $1$ で一致する。空間のトポロジーが特性指数を大域的に縛っていることの美しい実例である。
\end{Q}

\begin{Q}[【核心】なぜわざわざ、特性指数をこのような「表（リーマン図式）」の形式で記述するのか？]
\textbf{A.} 前段で示した通り、$x=0, 1, \infty$ の3点のみを確定特異点に持つ2階方程式（$m=2, n=2$）は「アクセサリー・パラメータを持たない」からである。
アクセサリー・パラメータが $0$ 個であるということは、「特異点の位置」と「各特異点の特性指数」さえ決めれば、元の微分方程式が一意に完全に復元できる（剛性を持つ）ことを意味する。つまり、このリーマン図式は単なるパラメータのメモ書きではなく、これ自体が「ガウスの超幾何関数」という対象を完全に過不足なく定義する「関数記号（リーマンのP関数）」そのものとして機能するのである。
\end{Q}
%#endregion
%#region
ここで、フックス型微分方程式\eqref{eq:linear-ode-n}の特性指数を操作する基本的な変換について述べる。
任意の定数 $\alpha_{1}, \dots, \alpha_{m}$ を用い、未知関数 $y$ を
\begin{equation}\label{eq:gauge_transform}
y = \left( \prod_{l=1}^{m} (x-\xi_{l})^{\alpha_{l}} \right) z
\end{equation}
によって新しい未知関数 $z$ に変数変換しても、微分方程式はフックス型のままである。
\begin{Q}[【核心】$y = \left( \prod_{l=1}^{m} (x-\xi_{l})^{\alpha_{l}} \right) z$ という変換を行っても、微分方程式が「フックス型（確定特異点のみを持つ）」であるという性質は本当に壊れないのか？]
\textbf{A.} 本当である。その理由は、変換因子 $\Phi(x) = \prod_{l=1}^{m} (x-\xi_{l})^{\alpha_{l}}$ の「対数微分」が、各特異点において高々1位の極しか生み出さないからである。
簡単なため2階のフックス型方程式 $y'' + p_{1}y' + p_{2}y = 0$ で確認する。
$y = \Phi z$ とおいて微分すると、
\begin{align*}
y' &= \Phi z' + \Phi' z \\
y'' &= \Phi z'' + 2\Phi' z' + \Phi'' z
\end{align*}
となる。これらを元の方程式に代入して $\Phi$ で割ると、$z$ についての微分方程式は
\[
z'' + \left( p_{1} + 2\frac{\Phi'}{\Phi} \right) z' + \left( p_{2} + p_{1}\frac{\Phi'}{\Phi} + \frac{\Phi''}{\Phi} \right) z = 0
\]
となる。新しい係数を $\tilde{p}_{1}(x), \tilde{p}_{2}(x)$ とおく。
ここで、$\Phi(x)$ の対数微分 $\frac{\Phi'}{\Phi}$ を計算すると、
\[
\frac{\Phi'(x)}{\Phi(x)} = \frac{d}{dx} \log \Phi(x) = \frac{d}{dx} \sum_{l=1}^{m} \alpha_{l} \log(x-\xi_{l}) = \sum_{l=1}^{m} \frac{\alpha_{l}}{x-\xi_{l}}
\]
となり、これは各 $x=\xi_{l}$ において「ちょうど1位の極」を持つ有理関数となる。
さらに、$\frac{\Phi''}{\Phi} = \left(\frac{\Phi'}{\Phi}\right)' + \left(\frac{\Phi'}{\Phi}\right)^{2}$ であるため、1位の極の微分（2位の極）と平方（2位の極）の和となり、これも「高々2位の極」しか持たない。
元の $p_{1}$ が1位の極、$p_{2}$ が2位の極を持つため、足し合わせた新しい係数 $\tilde{p}_{1}$ は依然として1位の極を、$\tilde{p}_{2}$ は依然として2位の極を持つにとどまる。
無限遠点に関しても同様のオーダー計算が成り立つため、このゲージ変換はフックス型の条件（極の位数制限）を一切破壊しないのである。
\end{Q}
このとき、変換後の微分方程式（$z$ の方程式）のリーマン図式は、元の図式から以下のように特性指数がシフトする。
\begin{itemize}
    \item 有限の特異点 $x=\xi_{l}$ における特性指数は、$s_{j}^{(l)} - \alpha_{l}$ となる。
    \item 無限遠点 $x=\infty$ における特性指数は、$s_{j}^{(\infty)} + \sum_{l=1}^{m} \alpha_{l}$ となる。
\end{itemize}

\begin{Q}[【核心】有限特異点では特性指数が $-\alpha_{l}$ と引き算されるのに、なぜ無限遠点の特性指数だけ $+\sum \alpha_{l}$ と足し算になるのか？]
\textbf{A.} 無限遠点の局所座標 $u=1/x$ で変換式を展開すれば明らかになる。
変換の因子を $x \to \infty$ （すなわち $u \to 0$）で展開すると、
\[
\prod_{l=1}^{m} (x-\xi_{l})^{\alpha_{l}} = \prod_{l=1}^{m} \left( \frac{1}{u} - \xi_{l} \right)^{\alpha_{l}} = u^{-\sum \alpha_{l}} \prod_{l=1}^{m} (1 - \xi_{l}u)^{\alpha_{l}} \sim u^{-\sum \alpha_{l}}
\]
となる。
解 $y$ が無限遠点で $u^{s^{(\infty)}}$ のオーダーで振る舞うとすれば、$u^{s^{(\infty)}} \sim u^{-\sum \alpha_{l}} z$ となるため、新しい未知関数 $z$ は $u^{s^{(\infty)} + \sum \alpha_{l}}$ のオーダーで振る舞うことになる。
（※これによって、有限特異点でのシフトの合計 $-n\sum \alpha_{l}$ と無限遠点でのシフトの合計 $+n\sum \alpha_{l}$ が見事に相殺し、フックスの関係式が保たれるという極めて美しい代数幾何学的構造になっている。）
\end{Q}

この特性指数のシフトを利用することで、方程式を標準形に帰着させることができる。
各有限特異点において $\alpha_{l} = s_{1}^{(l)}$ と選べば、変換後の有限特異点における特性指数のうち1つを
\[
s_{1}^{(1)} - \alpha_{1} = s_{1}^{(2)} - \alpha_{2} = \cdots = s_{1}^{(m)} - \alpha_{m} = 0
\]
と、すべて $0$ に揃えることができる。

特に $n=2, m=2$ （確定特異点が全体で3つ）のときは、アクセサリー・パラメータを持たない（$N_{acc}=0$）。
さらに、複素平面上の任意の相異なる3点は、一次分数変換（メビウス変換）によって必ず $x=0, 1, \infty$ の3点に写すことができる。
したがって、特異点を $0, 1, \infty$ に移動させた後、上記の\eqref{eq:gauge_transform}の変換を行って $x=0, 1$ の特性指数の1つを $0$ に揃えれば、リーマン図式は以下の非常にシンプルな形に帰着される。
\[
    \left\{
    \begin{matrix}
        x = 0 & x = 1 & x = \infty \\
        0 & 0 & s_{1}^{(\infty)} + s_{1}^{(1)} + s_{1}^{(2)} \\
        s_{2}^{(1)} - s_{1}^{(1)} & s_{2}^{(2)} - s_{1}^{(2)} & s_{2}^{(\infty)} + s_{1}^{(1)} + s_{1}^{(2)} \\
    \end{matrix}
    \right\}
\]
この図式の形は、パラメータを適当に $a, b, c$ と置き直せば、まさにガウスの超幾何微分方程式のリーマン図式と完全に一致する。

\begin{Q}[【核心】これらの一連の変換によって、「ガウスの超幾何微分方程式」の数学的な位置づけはどう変わるか？]
\textbf{A.} 「数ある特殊関数の1つ」から、「3つの確定特異点を持つすべての2階線形微分方程式の、普遍的な標準形（Universal normal form）」へと昇華される。
アクセサリー・パラメータが $0$ であるため、リーマン図式が一致すれば、背後にある微分方程式も完全に一致する（剛性）。つまり、物理学や幾何学のあらゆる問題で「3つの確定特異点を持つ方程式」が現れた瞬間、それはメビウス変換と冪関数の掛け算という「自明な操作」を除いて、すべてガウスの超幾何関数に他ならないことが保証されるのである。（※これはベルンハルト・リーマンがその驚異的な洞察力で到達した結論である。）
\end{Q}
%#endregion

\subsection{みかけの特異点}
%#region
\eqref{eq:linear-ode-n}の解の基本系 $\vec{\varphi}(x)$ をとる。確定特異点の集合 $\Xi$ を通らない閉曲線 $\gamma$ について解析接続を考えると、回路行列 $\Gamma(\gamma)$ を用いて
\[
\vec{\varphi}^{\gamma}(x) = \vec{\varphi}(x)\Gamma(\gamma)
\]
とかける。ここでもう一つの解の基本系 $\vec{\psi}(x)$ をとっておくと、可逆行列 $C$ を用いて
\[
\vec{\varphi}(x) = \vec{\psi}(x)C
\]
となる。また、解析接続と線型操作は交換できるので、
\[
\vec{\varphi}^{\gamma}(x) = \vec{\psi}^{\gamma}(x)C
\]
よって、$\vec{\psi}(x)$ の回路行列は、
\[
\vec{\psi}^{\gamma}(x) = \vec{\varphi}^{\gamma}(x)C^{-1} = \vec{\varphi}(x)\Gamma(\gamma)C^{-1} = \vec{\psi}(x)C\Gamma(\gamma)C^{-1}
\]
より、$C\Gamma(\gamma)C^{-1}$ となる。

\begin{definition}[リーマンデータ]\label{def:riemann_data}
    フックス型微分方程式 \eqref{eq:linear-ode-n} に対して、定義域 $\mathbb{P}^{1}(\mathbb{C})$、確定特異点の集合 $\Xi$、およびモノドロミー $\hat{\rho}$ の3つ組
    \[
    (\mathbb{P}^{1}(\mathbb{C}), \Xi, \hat{\rho})
    \]
    のことを\textbf{リーマンデータ}という。
\end{definition}

\begin{Q}[モノドロミーって何？]
\textbf{A.} $\cdot$ 閉曲線に対して回路行列を対応させる写像のようなもの。\\
\phantom{\textbf{A.}} $\cdot$ これは基本群 $\pi_1(X)$ から行列群への表現とみなすことができ、微分方程式の解の「多価性」を調べるための中心的な道具となる。(定義\ref{def:monodromy}を参照)\\
\end{Q}
%#endregion
%#region
ここで、モノドロミー群を決定するために必要なパラメータの数を考えたい。
基点 $x_{0}$ をうまく取ることによって、全ての有限の特異点が異なる偏角を持つようにでき、特に $\arg(\xi_{1}-x_{0}) < \arg(\xi_{2}-x_{0}) < \cdots < \arg(\xi_{m}-x_{0})$ となるように並べることができる。
ここで、それぞれの特異点 $\xi_{l}$ を反時計回りに一周だけして基点に戻ってくる閉道を $\gamma_{l}$ とし、$x=\infty$ を反時計回り（無限遠点から見て反時計回り）に一周して戻ってくる閉道を $\gamma_{\infty}$ とすると、空間の位相的な構造から
\[
[\gamma_{\infty}] = ([\gamma_{1}] \cdot [\gamma_{2}] \cdots [\gamma_{m}])^{-1} = [\gamma_{m}]^{-1} \cdots [\gamma_{2}]^{-1} \cdot [\gamma_{1}]^{-1}
\]
というホモトピー同値の関係が成り立つ。
よって、各閉道に沿った解析接続から定まるモノドロミー行列を $\Gamma(\gamma)$ とすれば、
\[
\Gamma(\gamma_{\infty}) = \Gamma(\gamma_{m})^{-1} \cdots \Gamma(\gamma_{2})^{-1} \Gamma(\gamma_{1})^{-1}
\]
となり、無限遠点でのモノドロミー行列は、有限特異点のモノドロミー行列の組から完全に決定される。
したがって、モノドロミー群の生成元全体を決定するには、$m$ 個の $n \times n$ 行列 $\Gamma(\gamma_{1}), \dots, \Gamma(\gamma_{m})$、すなわち $mn^{2}$ 個のパラメータを決めれば良いことがわかる。

さらに、モノドロミー表現は解の基本系の取り替え（$GL(n, \mathbb{C})$ による全体の相似変換）という冗長な自由度を持つ。
相似変換に用いる正則行列が持つ $n^{2}$ 個の自由度のうち、変換に寄与しない定数倍（群の中心 $Z(GL(n, \mathbb{C}))$）の $1$ 次元分を除いた「実質的な相似変換の自由度」は $n^{2}-1$ である。この分を差し引いて、次の命題を得る。

\begin{proposition}\label{prop:monodromy_parameters}
    フックス型微分方程式\eqref{eq:linear-ode-n}のモノドロミー表現のモジュライ空間は、
    \[
    M(n,m):=mn^{2} - (n^{2}-1) = (m-1)n^{2} + 1
    \]
    個の独立なパラメータを含む。
\end{proposition}

\begin{Q}[【核心】微分方程式自体のパラメータ数 $E(n,m)$ と、モノドロミー表現のパラメータ数 $(m-1)n^{2} + 1$ を比較すると、常に微分方程式の方がパラメータが多くなる（自由度が高い）。これは何を意味しているのか？]
\textbf{A.} 「同じモノドロミー群を持つフックス型微分方程式が、無数に存在する」ということである。
実際に差を計算してみると、特異点の数 $m$ や階数 $n$ が大きくなるにつれて、方程式の自由度 $E(n,m)$ の方が圧倒的に大きくなる。（※唯一の例外が、両者が一致するガウスの超幾何方程式 $E(2,2)=5$ である）。
「与えられたモノドロミー表現を持つフックス型方程式を構成せよ」という問題をヒルベルトの第21問題（リーマン・ヒルベルト問題）と呼ぶ。パラメータ数にギャップがあるため、モノドロミーを指定するだけでは方程式は一意に定まらず、余剰な自由度（アクセサリー・パラメータ）が生じてしまう。
そこで、この余剰な自由度をあえて「見かけの特異点」として幾何学的に配置し、本質的な特異点 $t$ を動かしたときに「モノドロミー群が一定に保たれるように」見かけの特異点を時間発展させる、という壮大なアイデアが生まれた。これが岡本和夫らによるモノドロミー保存変形の理論であり、パンルヴェ方程式のハミルトニアン構造の起源なのである。
\end{Q}

ここまで見てきたように、一般にパラメータの数について
\[
E_{0}(n,m) \le E(n,m) \le M(n,m)
\]
が成り立つ。実際、$m \ge 2, n \ge 2$ の場合、右側の不等式について評価すると、
\begin{align}
M(n,m)-E(n,m) &= (m-1)n^{2} + 1-\frac{1}{2}n \big( m(n+1) - (n-1) \big) \notag \\
&= \frac{1}{2}n \big\{ (m-1)(n-1)-2 \big\} + 1 \ge 0
\end{align}
となる。ここで、整数 $m \ge 2, n \ge 2$ のもとで上式がちょうど $0$ となる（等号が成立する）条件を探ると、$n=2, m=2$ のケースのみであることがわかる。このとき、
\[
E_{0}(2,2)=E(2,2)=M(2,2)=5
\]
となり、この極めて特殊な状況が、ガウスの超幾何微分方程式によって実現されている。

\begin{Q}[【核心】$E(n,m)=M(n,m)$ が成り立つ（$n=2, m=2$）ということは、微分方程式論において何を意味しているのか？なぜそれがガウスの超幾何微分方程式に結びつくのか？]
\textbf{A.} $E$ は「方程式を決定するために必要なパラメータ数」、$M$ は「見かけの特異点などを伴うアクセサリー・パラメータの数」を意味する。$M - E = 0$ となることは、\textbf{「余分なアクセサリー・パラメータが一切存在せず、特異点の位置とRiemann scheme（特性指数）だけで方程式がグローバルにただ一つに決定される」}という強固な性質（剛性, rigidity）を示している。

$m=2$ は有限確定特異点が2つ（通常 $x=0, 1$）、これに $\infty$ を加えた計3つの確定特異点を持つ状況を指す。階数 $n=2$ でこの条件を満たすものは、まさにガウスの超幾何微分方程式に他ならない。計算上の等号成立条件が、歴史的な名方程式の存在条件と見事に一致しているのである。
\end{Q}
%#endregion
%#region
また、$E(n,m)-E_{0}(n,m)$ は方程式を決定するための「アクセサリー・パラメータ」の個数を表しているが、さらに
\[
K(n,m)=M(n,m)-E(n,m)
\]
が何を意味しているのかを考察する。

ここで、「モノドロミー表現（大域的な性質）を変えることなく、微分方程式の確定特異点の数（局所的な性質）を増やすことは可能か？」という問いを立てる。
式 \eqref{eq:linear-ode-n} のリーマンデータが
\[
(\mathbb{P}^{1}(\mathbb{C}),\Xi,\hat{\rho}), \quad \Xi=\{\xi_{1},\cdots,\xi_{m},\xi_{m+1}\}
\]
で与えられているとし、ここに新しい特異点の集合 $\Lambda=\{\lambda_{1},\cdots,\lambda_{g}\}$ が追加されたとする。
これらの追加された特異点がモノドロミーを変化させないためには、$\Lambda$ の各点の周りでの解が1価（分岐を持たない）でなければならない。すなわち、各 $x=\lambda_{i}$ において以下の2条件を満たす必要がある。
\begin{itemize}
    \item[(a)] 特性指数が全て整数である。（任意の解に無理数や虚数の複素べきが現れない）
    \item[(b)] 非対数的特異点である。（フロベニウス法で解を構成する際、任意の解に対数関数 $\log(x-\lambda_{i})$ が現れない）
\end{itemize}

\begin{definition}[みかけの特異点]
    上の(a), (b)の性質を同時に満たす確定特異点 $x=\lambda_{i}$ を、\textbf{みかけの特異点 (apparent singularity)} という。
\end{definition}

では、リーマンデータとして $(\mathbb{P}^{1}(\mathbb{C}),\Xi\cup\Lambda,\hat{\rho})$ を持つ微分方程式 \eqref{eq:linear-ode-n} のパラメータ数を数え上げる。

\begin{Q}[【核心】みかけの特異点を $g$ 個追加したとき、微分方程式のパラメータ数はどう変化するのか？また、その物理的・幾何学的な意味は何か？]
\textbf{A.} 結論から言えば、パラメータの数は元の $E(n,m)$ からちょうど $g$ 個だけ増え、$E(n,m)+g$ となる。

みかけの特異点が1つ（$x=\lambda$）増えたときのパラメータの増減をトレースすると以下のようになる。
\begin{itemize}
    \item \textbf{特異点追加によるベースの増加}: 新たな確定特異点が1つ増えるため、微分方程式の係数に与える自由度は本来 $\frac{1}{2}n(n+1)$ 個増える（$E(n,m+1)-E(n,m)$ に相当）。
    \item \textbf{条件(1)による減少}: 特性指数 $n$ 個が全て「整数」に固定される。連続な複素パラメータ（自由度）が離散的な値に縛られるため、自由度は $n$ 個減少する。
    \item \textbf{条件(2)による減少}: 対数項を消去するための適合条件（フロベニウス法における各特性指数の差からの共鳴条件）が課される。これは独立な代数方程式を $\frac{1}{2}n(n-1)$ 個要求するため、自由度がさらに $\frac{1}{2}n(n-1)$ 個減少する。
    \item \textbf{特異点の位置という新たな自由度}: 本来の特異点 $\Xi$ は固定されているが、みかけの特異点 $\lambda$ の「位置」自体は微分方程式の係数（有理関数）の極として自由に動かせるため、自由度が $+1$ される。
\end{itemize}
これらを足し合わせると、
\[
\frac{1}{2}n(n+1) - n - \frac{1}{2}n(n-1) + 1 = 1
\]
となり、みかけの特異点1つにつきパラメータはちょうど1つ（特異点の位置 $\lambda$ の分だけ）増えることがわかる。
\end{Q}

\begin{proposition}
    微分方程式 \eqref{eq:linear-ode-n} が、本来の確定特異点 $\Xi$ の他に、$g$ 個のみかけの特異点 $\Lambda$ を持つとすると、その微分方程式を決定するパラメータの総数は
    \[
    E(n,m)+g
    \]
    で与えられる。
\end{proposition}

\begin{Q}[【核心】予備知識なしに微分方程式の係数 $p_i(x)$ の形（有理関数）だけを見たとき、本来の確定特異点 $\xi$ とみかけの特異点 $\lambda$ は区別できるのか？]
\textbf{A.} 一見しただけでは全く区別できないが、「局所的な解析」を行えば完全に区別できる。

微分方程式の係数 $p_i(x)$ の分母の極として現れる点という視点では、$\xi$ も $\lambda$ も全く同じ「確定特異点」として振る舞う。しかし、ある特異点 $x=a$ のまわりでフロベニウス法を用いて級数解を構成しようとしたとき、以下の2つの厳しいテストをクリアした場合にのみ、その点 $a$ は「みかけの特異点である」と事後的に判定される。
\begin{enumerate}
    \item 決定方程式（indicial equation）を解くと、特性指数が「全て整数」になる。
    \item 特性指数の差が整数であることに起因する「対数項（$\log$）が発生するリスク」を、係数同士の絶妙な関係式（適合条件）によって完全に回避している。
\end{enumerate}
すなわち、みかけの特異点とは「係数の分母がゼロになるため見た目は特異点だが、解の分岐（モノドロミー）を一切引き起こさないように、係数パラメータが奇跡的なバランスで調整された点」なのである。
\end{Q}

\begin{Q}[【核心】みかけの特異点 $\lambda$ を追加した際のパラメータ増減を計算する際、なぜ最初は「特異点を固定して定義した」はずの公式 $E(n,m+1)$ を使って計算を出発しているのか？定義と矛盾していないか？]
\textbf{A.} 全く矛盾していない。これは「一旦、位置を固定した普通の確定特異点として追加してから、みかけの特異点へとクラスチェンジさせる」という段階的な思考実験（計算テクニック）を行っているからである。

具体的には、以下の3ステップの論理でパラメータを数え上げている。
\begin{enumerate}
    \item \textbf{枠組みの拡張（代数的な自由度の追加）}:
    まず、点 $x=\lambda$ を「位置が固定されたただの確定特異点」とみなして係数の分母に追加する。この時点では単なる極の追加であるため、微分方程式の係数が持つ代数的な自由度は $E(n,m+1) - E(n,m) = \frac{1}{2}n(n+1)$ 個増える。
    
    \item \textbf{条件の付加（自由度の剥奪）}:
    次に、この点 $x=\lambda$ が「みかけの特異点」になるための厳しい条件（特性指数が整数、対数項なし）を後から課す。これにより、ステップ1で追加された自由度のうち、ちょうど $\frac{1}{2}n(n+1)$ 個のパラメータが「条件を満たすための調整役」として完全に束縛される。結果として、係数由来の自由度の増加は「相殺されてゼロ」になる。
    
    \item \textbf{パラダイムシフト（位置のパラメータ化）}:
    最後に、条件を満たしたことで解の分岐（モノドロミー）に一切影響を与えなくなった点 $x=\lambda$ に対して、「位置を固定しなければならない」という足かせを外す。係数パラメータの自由度が増えなかった代わりに、「$\lambda$ という座標そのもの」が新たな動的パラメータ（$+1$）として復活する。
\end{enumerate}
すなわち、$E(n,m+1)$ は、計算の土俵に新しい極を乗せるための「一時的な踏み台」として機能しているのである。
\end{Q}

\begin{Q}[【核心】「みかけの特異点 $\lambda$ の位置がパラメータになる」という言葉の真意は何か？単に1つの微分方程式の中で、特異点を物理的に動かしているという意味か？]
\textbf{A.} そうではない。本質的な意味は、\textbf{「同じモノドロミー（大域的性質）を共有する微分方程式が、$\lambda$ の取り得る値の分だけ無数に存在し、ひとつの『族（family）』をなしている」}ということである。

本来の確定特異点 $\Xi$ とモノドロミー表現 $\hat{\rho}$ を完全に固定したとする。もしみかけの特異点が存在しなければ、方程式はただ1つに決定される（剛性）。しかし、みかけの特異点 $\lambda$ が存在すると、その位置 $\lambda$ を別の値 $\lambda'$ に変えるごとに、「係数の形は異なるが、モノドロミーは全く同じ別の微分方程式」を次々と作ることができる。

つまり、「$\lambda$ を動かす」というのは、\textbf{「同じモノドロミーを持つ方程式の空間（モジュライ空間）の中を連続的に移動し、別の方程式へ乗り換えていくこと」}を意味する。$\lambda$ は、その空間内の特定の方程式を指し示す「座標（インデックス）」として機能しているのである。このように、モノドロミーを不変に保ったまま微分方程式を連続的に変形させる概念を\textbf{等モノドロミー変形 (isomonodromic deformation)} と呼ぶ。
\end{Q}

ここで、改めて冒頭で定義した
\[
K(n,m)=M(n,m)-E(n,m)
\]
という量が何を意味しているのかを総括する。

\begin{Q}[【核心】モノドロミーのパラメータ数 $M$ と、微分方程式のパラメータ数 $E$ の差である $K(n,m)$ は、本質的に何を表している数なのか？]
\textbf{A.} 結論から言えば、$K(n,m)$ は「任意のモノドロミーを実現するために、\textbf{最低限追加しなければならない『みかけの特異点』の個数}」を表している。

微分方程式の理論における究極の目標の一つは、「与えられたモノドロミー表現 $\hat{\rho}$ から、それを持つフックス型微分方程式を逆算して構成すること」である（リーマン・ヒルベルト問題）。
しかし、一般に $M(n,m) > E(n,m)$ である場合、微分方程式側の自由度（パラメータ）が圧倒的に足りず、望みのモノドロミーを実現する方程式が存在しない。

そこで、モノドロミーを一切変えずに微分方程式側の自由度だけを増やす「助っ人」として、みかけの特異点 $\Lambda=\{\lambda_{1},\cdots,\lambda_{g}\}$ を導入する。みかけの特異点を $g$ 個追加すると、微分方程式のパラメータ数は $E(n,m) + g$ となる。
モノドロミーの自由度と微分方程式の自由度が完全に一致し、対応がつく（モジュライ空間の次元が等しくなる）ための条件は、
\[
E(n,m) + g = M(n,m)
\]
すなわち、
\[
g = M(n,m) - E(n,m) = K(n,m)
\]
となる。
つまり、一般のフックス型微分方程式の分類・構成において、我々は最初に本来の確定特異点 $\Xi$ を決めた後、「次元の不足分 $K(n,m)$ を補うために、同数のみかけの特異点 $\Lambda$ を導入する」という手続きを必然的に踏むことになるのである。
\end{Q}

\begin{Q}[【発展】次元の不足を補うためにみかけの特異点を追加するならば、最低限の $K(n,m)$ 個よりも多く（$g > K(n,m)$ 個）追加しても、与えられたモノドロミーを実現することは可能か？また、その場合に何が起こるのか？]
\textbf{A.} 実現することは可能であるが、微分方程式側に「冗長な自由度」が発生し、モノドロミーと微分方程式の1対1の美しい対応が崩れてしまう。

これを「微分方程式の係数の空間（次元 $E+g$）」から「モノドロミー表現の空間（次元 $M$）」への写像として捉えると、みかけの特異点の個数 $g$ によって以下のような状況の違いが生じる。

\begin{itemize}
    \item \textbf{$g < K$ の場合（自由度不足）}: 
    $E+g < M$ となり、写像の像がターゲット全体を覆えない。つまり、実現不可能なモノドロミーが多数存在する。
    
    \item \textbf{$g = K$ の場合（ジャスト・フィット）}:
    $E+g = M$ となり、両空間の次元が一致する。一般のモノドロミーに対して、それを実現する微分方程式が（本質的に）ただ1つ、あるいは有限個に定まる。無駄のない「剛性」が保たれた理想的な状態である。
    
    \item \textbf{$g > K$ の場合（自由度過剰）}:
    $E+g > M$ となるため、任意のモノドロミーを実現することは可能である（全射的）。しかし、次元が $g-K$ だけ余るため、写像は単射にはならない。結果として、「全く同じモノドロミー表現を持つのに、係数（みかけの特異点の配置など）が連続的に異なる微分方程式の族（次元 $g-K$ のファイバー）」が無数に発生してしまう。
\end{itemize}

のちにパンルヴェ方程式の理論（岡本和夫のハミルトン理論）において、みかけの特異点の位置 $q$ と、その点での微分方程式の係数に関するある量 $p$ が、ハミルトン力学系における「正準変数（座標と運動量）」として扱われる。力学系として非退化な（決定論的に未来が1つに決まる）相空間を構成するためには、この過剰な冗長性を排除し、$g=K$ という「必要最低限のモジュライ空間」を考えることが不可欠となるのである。
\end{Q}
%#endregion
%#region
また、みかけの特異点 $x=\lambda$ における特性指数が全て整数 $J_{1}, \cdots, J_{n}$ であるとき、未知関数を
\begin{equation}\label{eq:gauge_transform_local}
z = (x-\lambda)^{s} y
\end{equation}
と変換（ゲージ変換）すると、新しい関数 $z$ が満たす微分方程式の $x=\lambda$ における特性指数は、一律に $s$ だけシフトして
\[
s+J_{1}, \cdots, s+J_{n}
\]
となる。元の解 $y$ は対数項を持たずローラン展開可能（非対数的特異点）であったため、それに $(x-\lambda)^{s}$ を掛けただけの $z$ も当然対数項を持たない。
このとき、$z$ の微分方程式について点 $x=\lambda$ の周りを1周する閉経路 $\gamma$ に沿った解析接続を考えると、モノドロミー行列はスカラー行列
\[
\Gamma(\gamma) = e^{2\pi is} I_{n}
\]
となる。逆に、ある特異点でのモノドロミー行列がこのスカラー行列の形で書けるならば、逆変換 $y = (x-\lambda)^{-s} z$ を施すことで特性指数を全て整数にでき、かつ対数項を持たない（非対数的である）ため、それはみかけの特異点に帰着される。

\begin{definition}[みかけの特異点の拡張]
解が多価関数となる場合（すなわち $s \notin \mathbb{Z}$）であっても、適当な複素数 $s$ によって \eqref{eq:gauge_transform_local} の変換を行い、特性指数を全て整数にでき、かつ非対数的となるような特異点のことを、広い意味で \textbf{みかけの特異点} と呼ぶ。
\end{definition}

さて、$x=\lambda$ がみかけの特異点であるならば、変換 \eqref{eq:gauge_transform_local} で適切に $s$ を選ぶことで、すべての特性指数を「非負の整数」に揃えることができる。このとき、微分方程式の任意の解は $x=\lambda$ の近傍で正則（テイラー展開可能）となる。

\begin{Q}[【核心】任意の解が正則（特異性を持たない）になるように変換できたとしても、微分方程式の係数にとっては $x=\lambda$ が依然として極（特異点）になり得るのはなぜか？分子と分母が打ち消し合って極が消えることはないのか？]
\textbf{A.} 解が正則であっても、解のロンスキアン $W(x)$ が $x=\lambda$ でゼロになれば、必ず方程式の係数に極が発生する。

$n$ 階線形微分方程式 $y^{(n)} + p_1(x) y^{(n-1)} + \dots + p_n(x) y = 0$ において、係数 $p_1(x)$ と基本解のロンスキアン $W(x)$ の間には、アーベルの公式（リウヴィルの公式）により
\[
p_1(x) = -\frac{W'(x)}{W(x)}
\]
という関係がある。
解がすべて正則であるとき、$W(x)$ も正則関数となる。もし $W(\lambda) = 0$ ならば、$x=\lambda$ は $W(x)$ の $m$ 位のゼロ点（$m \ge 1$）となり、$W(x) = c(x-\lambda)^m + \cdots$ と展開できる。これをアーベルの公式に代入すると、
\[
p_1(x) = -\frac{cm(x-\lambda)^{m-1} + \cdots}{c(x-\lambda)^m + \cdots} = -\frac{m}{x-\lambda} + (\text{正則な項})
\]
となり、$p_1(x)$ は必ず $x=\lambda$ において留数 $-m$ の1位の極を持つ。したがって、分子との打ち消し合いで極が消滅することは原理的にあり得ず、微分方程式にとっては確実に特異点となる。

これが特異点ではなく「完全な正則点（係数も極を持たない点）」に還元されるための必要十分条件は、$W(\lambda) \neq 0$ となることである。
みかけの特異点においてこれを満たすのは、特性指数がちょうど
\[
0, 1, 2, \cdots, n-1
\]
という「抜けのない連続した整数」になる場合のみである。このとき、適切に基底を取り直すことで、各解のテイラー展開の先頭項を $1, (x-\lambda), \dots, (x-\lambda)^{n-1}$ とできる。これらの $x=\lambda$ における微分値を並べたロンスキアン行列は、対角成分に $0!, 1!, \dots, (n-1)!$ が並ぶ非退化な三角行列となる。したがって、行列式 $W(\lambda)$ は対角成分の積となり $W(\lambda) \neq 0$ が保証され、係数の極が完全に消滅する。
\end{Q}
%#endregion

\subsection{連立微分方程式系と単独高階微分方程式}

%#region
これまでの単独の高階微分方程式から視点を変え、以下の連立1階線形微分方程式系について考えたい。
\begin{equation}\label{eq:system_ode}
    \frac{d\vec{y}}{dx} = A(x)\vec{y}
\end{equation}
ここで、$A(x)$ は有理関数を成分とする $n \times n$ 行列であり、$\vec{y}$ は $n$ 次元の未知ベクトル関数である。

\begin{definition}
    方程式 \eqref{eq:system_ode} の任意の解ベクトル $\vec{\varphi}(x)$ の各成分 $\varphi_{i}(x)$ が、$x=\xi$ において高々多項式オーダーの増大度しか持たない（すなわち確定特異点である）とき、方程式 \eqref{eq:system_ode} は $x=\xi$ に \textbf{確定特異点} をもつという。
\end{definition}

\begin{definition}
    方程式 \eqref{eq:system_ode} の特異点が全て確定特異点であるとき、\eqref{eq:system_ode} を \textbf{フックス型微分方程式系 (Fuchsian system)} という。
\end{definition}

無限遠点 $x=\infty$ も特異点として許容し、確定特異点であるとする。また、解はベクトル値関数であるが、解析接続によるモノドロミー行列は単独方程式のときと同様に定義できるため、リーマンデータ $(\mathbb{P}^{1}(\mathbb{C}), \Xi, \hat{\rho})$ を考えることができる。

いま、有限の特異点に対しては $u=x-\xi$、無限遠点に対しては $u=\frac{1}{x}$ という局所座標を導入し、方程式 \eqref{eq:system_ode} を
\begin{equation}\label{eq:system_ode_local}
    u\frac{d\vec{y}}{du} = B(u)\vec{y}
\end{equation}
と書き換える。このとき、以下の十分条件が成り立つ。

\begin{proposition}\label{prop:fuchs_condition_system}
    \eqref{eq:system_ode_local} において、行列 $B(u)$ が $u=0$ の近傍で正則（各成分が極を持たない）ならば、$u=0$ （すなわち元の $x=\xi$ または $\infty$）は \eqref{eq:system_ode} の確定特異点である。
\end{proposition}

\begin{proof}
    $u=0$ を頂点とする任意の角領域（セクター）内で原点に向かう経路を考え、$u = r e^{i\theta}$ （$r \to +0$, $\theta$ は固定）と極座標表示する。
    このとき、微分演算子は $u \frac{d}{du} = r \frac{d}{dr}$ となるため、方程式 \eqref{eq:system_ode_local} は
    \[
    r \frac{d\vec{y}}{dr} = B(re^{i\theta})\vec{y}
    \]
    と書ける。
    仮定より $B(u)$ は $u=0$ の近傍で正則であるため、その成分は連続関数であり、原点を含むある閉円板 $|u| \le r_0$ において有界となる。すなわち、ある正の定数 $M$ が存在して、行列の作用素ノルムについて $\|B(u)\| \le M$ が成り立つ。
    
    ベクトル $\vec{y}(r)$ のノルム $\|\vec{y}\|$ について、微分不等式を評価する。
    \begin{align*}
        r \frac{d}{dr} \|\vec{y}\| &\le \left\| r \frac{d\vec{y}}{dr} \right\| \\
        &= \| B(re^{i\theta})\vec{y} \| \\
        &\le \| B(re^{i\theta}) \| \|\vec{y}\| \\
        &\le M \|\vec{y}\|
    \end{align*}
    両辺を $r \|\vec{y}\|$ ($>0$) で割ると、
    \[
    \frac{1}{\|\vec{y}\|} \frac{d\|\vec{y}\|}{dr} \le \frac{M}{r}
    \]
    これをある点 $r_1$ から $r$ ($0 < r < r_1 \le r_0$) まで積分する。経路の向きに注意して $- \int_r^{r_1}$ として評価すると、
    \[
    \log \|\vec{y}(r_1)\| - \log \|\vec{y}(r)\| \le M (\log r_1 - \log r)
    \]
    移項して整理すると、
    \[
    \log \|\vec{y}(r)\| \ge \log \|\vec{y}(r_1)\| - M \log(r_1 / r) \implies \|\vec{y}(r)\| \ge C r^M
    \]
    となる。（※注：特異点へ向かう極限は $r \to 0$ であり、マイナス乗の評価が必要なため、原点から遠ざかる方向へ不等式を逆に解く。）
    正確には、$\frac{d}{dr}(\log \|\vec{y}\|) \ge - \frac{M}{r}$ （絶対値の微分の性質から下限も評価可能）より、
    \[
    \|\vec{y}(r)\| \le C' r^{-M} = C' |u|^{-M}
    \]
    という上からの評価が得られる（$C'$ は積分定数に依存する正の定数）。
    これにより、解ベクトル $\vec{y}(u)$ の各成分は、$u \to 0$ において高々 $u$ の多項式（マイナス乗）のオーダーでしか増大しないことが示された。これは確定特異点の定義そのものであり、証明が完了する。
\end{proof}

\begin{Q}[【核心】上記の証明において、解の増大度が $\|\vec{y}(u)\| \le C' |u|^{-M}$ で抑えられたことが、なぜ「確定特異点であること」の直接的な証明になっているのか？]
\textbf{A.} 確定特異点（regular singular point）の現代的な定義が、「特異点へ任意のセクター内から近づくとき、解の増大度が高々多項式オーダー（$|u|^{-N}$ の形）に制限されること」だからである。

もし $u=0$ が不確定特異点（真性特異点のような激しい振る舞いをする点）であれば、解の成分には $e^{1/u}$ のような指数関数的な発散項が含まれる。このような関数は $u \to 0$ で $|u|^{-N}$ といういかなる多項式の逆数をもってしても上から抑え込むことができない。
証明の中で $B(u)$ の正則性（有界性 $\|B\| \le M$）を用いたことで、指数の肩に乗る発散のスピードが対数積分 $\int \frac{M}{r} dr = M \log r$ に制限され、結果として解の増大度が多項式オーダーに落ち着くことが完全に保証されたのである。
\end{Q}

この命題は、「$A(x)$ が $x=\xi$ で1位の極を持つ（単純極である）ならば、確定特異点である」ということを意味している。

\begin{Q}[【核心】上記の命題の逆、すなわち「$x=\xi$ が確定特異点ならば、必ず $B(u)$ は正則になる（$A(x)$ は高々1位の極しか持たない）」は成立するか？成り立たないとすれば、どのような反例があるか？]
\textbf{A.} 逆は全く成り立たない。驚くべきことに、単独 $n$ 階微分方程式を連立系に書き直した「コンパニオン行列」自体が、完璧な反例となっている。

単独のフックス型微分方程式 \eqref{eq:linear-ode-n} を考える。
\[
\vec{y}=
\begin{pmatrix}
    y_{1} \\ y_{2} \\ \vdots \\ y_{n}
\end{pmatrix}
=
\begin{pmatrix}
    y \\ \frac{dy}{dx} \\ \vdots \\ \frac{d^{n-1}y}{dx^{n-1}}
\end{pmatrix}
\]
とおくと、方程式は \eqref{eq:system_ode} の形になり、その係数行列 $A(x)$ は
\[
A(x)=
\begin{pmatrix}
    0 & 1 & 0 &\cdots & 0 \\
    0 & 0 & 1 & \cdots & 0\\
    \vdots&\vdots&\vdots& \ddots &\vdots\\
    0 & 0 & 0 & \cdots & 1 \\
    -p_{n}(x) & -p_{n-1}(x) & -p_{n-2}(x) & \cdots & -p_{1}(x)
\end{pmatrix}
\]
となる。元の単独方程式において $x=\xi$ が確定特異点であるとき、フックスの定理（命題\ref{prop:fuchs_theorem_n}）より $p_{k}(x)$ は $k$ 位の極を持つ。したがって、行列 $A(x)$ の一番左下の成分 $-p_{n}(x)$ は「$n$ 位の極」を持つことになる。
解ベクトル $\vec{y}$ の成分は全て確定特異点の性質を満たしているにもかかわらず、行列 $A(x)$ は1位の極どころか $n$ 位の極を持っており、$u A(u)$ は全く正則にならないのである。
\end{Q}

\begin{Q}[【核心】単独 $n$ 階方程式と連立1階系では、理論を展開する上でどのような利点と欠点の違いがあるか？]
\textbf{A.} 最大の違いは「表現の一意性（ゲージ変換の自由度）」にある。

\textbf{単独方程式の利点}は、方程式の「表示が一意」に定まることである。解の空間が決まれば方程式の係数 $p_k(x)$ はロンスキアンを通じて一意に決まり、フックスの定理によって特異点における「極の位数」が厳密に統制されるため、外見から特異点の性質を判定しやすい。

一方、\textbf{連立系の利点}は、行列式や固有値といった線形代数の強力なツールを直接適用できることである。しかし欠点として、$\vec{z} = T(x)\vec{y}$ （$T(x)$ は有理関数行列）というゲージ変換（基底の取り替え）を行うと、方程式は $\vec{z}' = (T'T^{-1} + TAT^{-1})\vec{z}$ となり、解の確定特異点としての性質（モノドロミーなど）を全く変えないまま、\textbf{係数行列の極の位数をいくらでも人為的に上げ下げできてしまう}。
この「表示の非一意性」こそが、連立系において見かけの極の位数に騙されず、本質的な特異点の性質を見抜くための理論（Poincaré rankの簡約化など）を必要とする理由である。
\end{Q}
%#endregion
%#region
$x$ の有理関数を成分とする可逆な（行列式が恒等的に0でない）$n$ 次正方行列 $T(x)$ をとって、未知関数を
\begin{equation}\label{eq:gauge_transform_x}
    \vec{y} = T(x)\vec{z}
\end{equation}
と変換（ゲージ変換）すると、方程式 \eqref{eq:system_ode} は微分の連鎖律により
\begin{equation}
    \frac{d\vec{z}}{dx} = C(x)\vec{z}, \quad \left( C(x) = T(x)^{-1}A(x)T(x) - T(x)^{-1}\frac{dT(x)}{dx} \right)
\end{equation}
と書ける。このとき、「$x=\xi$ が確定特異点ならば、局所的に $T(x)$ をうまく取ることで、$C(x)$ が $x=\xi$ において1位の極（単純極）を持つようにできる」という事実が知られている。前述の命題と合わせると、局所的には確定特異点と1位の極は同値な概念として扱える。

これを踏まえて、$A_{l} \ (l=1,\cdots,m)$ を $x$ に依らない定数行列として、特異点における留数を明示した以下の微分方程式系を考える。
\begin{equation}\label{eq:schlesinger}
    \frac{d\vec{y}}{dx} = \sum_{l=1}^{m}\frac{A_{l}}{x-\xi_{l}}\vec{y}
\end{equation}
これは $\Xi=\{\xi_{1},\cdots,\xi_{m},\infty\}$ を特異点にもち、すべての有限特異点で係数行列が1位の極を持つため、十分条件（命題\ref{prop:fuchs_condition_system}）を満たす確定特異点となる。無限遠点 $x=\infty$ についても局所座標を用いて計算すると確定特異点であることがわかるため、これは大域的にもフックス型微分方程式系である。

\begin{definition}[シュレージンガー型微分方程式]
    \eqref{eq:schlesinger} のように、係数行列が「1位の極の和」として表されるフックス型微分方程式系を、\textbf{シュレージンガー型微分方程式 (Schlesinger system)} という。
\end{definition}

\begin{Q}[【核心】「局所的」にはゲージ変換で確定特異点を1位の極に変形できるが、では「大域的」に、任意のフックス型微分方程式系を1つのゲージ変換でシュレージンガー型 \eqref{eq:schlesinger} に変形することは可能なのか？]
\textbf{A.} 一般には不可能である。（ここにリーマン・ヒルベルト問題の罠がある）

ある1つの特異点 $x=\xi$ の近傍だけであれば、有理型関数を用いた局所的なゲージ変換によって、確定特異点を必ず1位の極に還元できる。
しかし、1つの全域的な有理関数行列 $T(x)$ を用いて、\textbf{「すべての特異点を同時に1位の極に落としつつ、空間に余計な特異点を一切増やさない」}という都合の良い大域的な変換ができるとは限らない。

もしこれが常に可能であれば、任意のモノドロミー表現から必ずシュレージンガー型方程式が作れることになり、リーマン・ヒルベルト問題は無条件で肯定的に解けることになってしまう。長年そのように信じられていたが、1989年にボリブルフ (Bolibruch) が「シュレージンガー型の系では絶対に実現できないモノドロミー表現が存在する」という反例を構成した。
つまり、一般のフックス型系を無理やりシュレージンガー型の美しい形に書き換えようとすると、大域的な次元の辻褄を合わせるために、どうしても新たな「みかけの特異点」を系に発生させざるを得なくなるのである。
\end{Q}

シュレージンガー型微分方程式 \eqref{eq:schlesinger} に対して、$x$ に依存しない定数行列を用いた大域的なゲージ変換 $T(x) := G \in GL(n, \mathbb{C})$ を考えると、各特異点の留数行列は
\[
A_{l} \to G^{-1}A_{l}G
\]
と変換される。この定数ゲージ変換の自由度を用いることで、どれか1つの留数行列を標準形（ジョルダン標準形や対角化など）に固定することができる。
ここで、無限遠点 $x=\infty$ における留数行列 $A_{\infty}$ について考える。コンパクトなリーマン面（ここではリーマン球面 $\mathbb{P}^{1}(\mathbb{C})$）上の有理型微分形式の留数の総和は必ず $0$ になるという留数定理（あるいは $u=1/x$ による局所座標での直接計算）より、
\[
A_{\infty} = -\sum_{l=1}^{m}A_{l}
\]
が成り立つ。我々は慣例として、この $A_{\infty}$ をゲージ変換 $G$ によって標準形に固定することにする。

では、このシュレージンガー型方程式に含まれる「本質的な独立パラメータの数」を数え上げる。

\begin{Q}[【核心】シュレージンガー型微分方程式 \eqref{eq:schlesinger} が持つ本質的なパラメータの総数はいくつか？また、その計算の根拠は何か？]
\textbf{A.} パラメータの総数はちょうど $(m-1)n^{2}+1$ 個である。

計算の根拠は以下の通りである。
\begin{itemize}
    \item \textbf{元のパラメータ数}: 方程式を決定するのは $m$ 個の $n \times n$ 行列 $A_{1}, \cdots, A_{m}$ であるため、最初は $mn^{2}$ 個の自由度がある。（$A_{\infty}$ はこれらの和で決まるため独立ではない）
    \item \textbf{ゲージ変換による削減}: 定数行列 $G$ による共役変換 $G^{-1}(\cdot)G$ で標準形に固定できるため、$G$ の持つ自由度だけパラメータが減る。$G$ は $n \times n$ の可逆行列なので基本的には $n^{2}$ 個の自由度を持つ。
    \item \textbf{スカラー行列の除外}: しかし、$G = cI$（$c$ は非ゼロの定数、$I$ は単位行列）という「スカラー行列」を用いた場合、$c^{-1}I A_l cI = A_l$ となり、行列を全く変化させない。つまり、変換として有効に働く自由度はスカラー行列の分（1次元）を除いた $n^{2}-1$ 個である（射影一般線形群 $PGL(n, \mathbb{C})$ の次元）。
\end{itemize}
したがって、本質的なパラメータ数は $mn^{2}$ から有効なゲージ自由度 $n^{2}-1$ を引いて、
\[
mn^{2} - (n^{2} - 1) = (m-1)n^{2} + 1
\]
となる。
\end{Q}

\begin{Q}[【核心】このパラメータ数 $(m-1)n^{2}+1$ は、これまでの議論とどう結びつくのか？]
\textbf{A.} 驚くべきことに、この数は $\Xi=\{\xi_{1},\cdots,\xi_{m},\infty\}$ を特異点にもつ $n$ 階フックス型微分方程式の「モノドロミー表現のモジュライ空間の次元 $M(n,m)$」と完全に一致する。

単独のフックス型微分方程式では、パラメータ数 $E(n,m)$ が $M(n,m)$ に全く届かず、次元の不足を補うために「みかけの特異点」を多数導入する必要があった。
しかし、連立系である「シュレージンガー型微分方程式」を考えた瞬間、みかけの特異点を一切追加することなく、最初からモノドロミーと同じ自由度 $M(n,m)$ を獲得できている。この見事な次元の一致こそが、シュレージンガー系がリーマン・ヒルベルト問題（与えられたモノドロミーを実現する方程式の構成）を解くための「最も自然で強力な枠組み」として期待され、研究されてきた最大の理由である。
\end{Q}
%#endregion
%#region
\begin{definition}[可約な系]
    \eqref{eq:system_ode} の係数行列 $A(x)$ が、ある整数 $r$ （$1 \le r \le n-1$）に対して、$r$ 次正方行列 $A_{11}(x)$ と $n-r$ 次正方行列 $A_{22}(x)$ を用いて
    \[
    A(x)=
    \begin{pmatrix}
    A_{11}(x) & 0\\
    A_{21}(x) & A_{22}(x)
    \end{pmatrix}
    \]
    というブロック下三角行列の形で表せるとき、この微分方程式系 \eqref{eq:system_ode} は \textbf{可約 (reducible)} であるという。
\end{definition}

\begin{Q}[【核心】係数行列が可約（ブロック下三角）であるとき、微分方程式を解く上でどのような具体的なメリットがあるのか？]
\textbf{A.} $n$ 次元の巨大な連立方程式を、「$n-1$ 次元以下のより小さな部分系のカスケード（数珠つなぎ）解法」に完全に帰着できることである。

未知ベクトルを $\vec{y} = \begin{pmatrix} \vec{y}_{1} \\ \vec{y}_{2} \end{pmatrix}$ と分割して代入すると、系は以下の2つのステップに分解される。
\begin{enumerate}
    \item $\frac{d\vec{y}_{1}}{dx} = A_{11}(x)\vec{y}_{1}$
    \item $\frac{d\vec{y}_{2}}{dx} = A_{22}(x)\vec{y}_{2} + A_{21}(x)\vec{y}_{1}$
\end{enumerate}
まず、ステップ1の $r$ 次元斉次方程式を解いて $\vec{y}_{1}$ を完全に決定する。次に、その結果をステップ2の第2項（非斉次項）に代入し、定数変化法などを用いて $n-r$ 次元の非斉次方程式を解く。
いずれのステップも最大で $\max(r, n-r) \le n-1$ 次元の系を解けばよくなるため、方程式の難易度が劇的に下がるのである。
\end{Q}

\begin{Q}[【核心】単独 $n$ 階微分方程式を連立系に直した「コンパニオン行列（同伴行列）」は、そのままの形では絶対に可約にならないのはなぜか？]
\textbf{A.} コンパニオン行列の「一つ上の対角成分（超対角成分）」がすべて $1$ で埋まっているからである。

コンパニオン行列をどのように $r$ 行と $n-r$ 行でブロック分割しようとしても、右上ブロック（$A_{12}(x)$ に相当する部分）の境界付近には必ず「超対角成分の $1$」が少なくとも1つ紛れ込んでしまう。したがって、右上ブロックを完全に零行列 $0$ にすることは物理的に不可能である。
これは、「基底の取り方（表現）を工夫しない限り、単独方程式の自然な連立化は、自動的には小さな系に分解してくれない」という単独系の強固な構造を視覚的に表している。（※ただし、元の単独微分方程式がより低階の微分作用素に因数分解できるような特殊なケースに限り、適切なゲージ変換 $T(x)$ を挟むことで、系全体を可約な形に変換することは可能である。）
\end{Q}
%#endregion
%#region
\begin{definition}[単独高階化]
    連立微分方程式系を、特定の成分（例えば $y_1$）について $x$ で繰り返し微分し、代入と消去を繰り返すことで、元の成分を未知関数とする1つの $n$ 階単独微分方程式に帰着させる操作を \textbf{単独高階化} という。（※代数的には「巡回ベクトル (cyclic vector)」をとる操作に相当する。）
\end{definition}

既約（簡約不能）な $n$ 次連立微分方程式系に対して単独高階化を行うと、同値な $n$ 階単独微分方程式が得られ、後者に対してこれまでの単独方程式の議論を適用することができる。
また、元の連立系がフックス型（すべての特異点が確定特異点）であれば、得られた単独方程式もフックス型になる。なぜなら、「解が高々多項式オーダーの増大度しか持たない」という確定特異点の解析的な性質は、成分の足し算や微分といった線形代数的な操作を行っても破壊されないからである。

例えば、特異点 $\Xi=\{\xi_{1},\cdots,\xi_{m},\infty\}$ を持つシュレージンガー型方程式 \eqref{eq:schlesinger} について単独高階化を行ったとする。元の本来の確定特異点の位置は当然変わらない。しかしここで、方程式に含まれる「パラメータの数」について重大なパラドックスが生じる。

前述の通り、本来の特異点のみを持つ単独系のパラメータ数は
\[
E(n,m) = \frac{1}{2}n \big( m(n+1) - (n-1) \big)
\]
であった。一方で、シュレージンガー型微分方程式（連立系）のパラメータ数は、モノドロミー表現の次元と等しい
\[
M(n,m) = (m-1)n^{2} + 1
\]
であった。両者は全く等価な微分方程式であるはずなのに、パラメータの数が一致しない。

\begin{Q}[【核心】単独高階化を行った結果、$K(n,m) = M(n,m) - E(n,m)$ 個分のパラメータが空中に消え去ってしまったように見える。この「消えたパラメータ」はどこへ行ったのか？]
\textbf{A.} 消えたわけではない。連立系から単独系へ「翻訳」する過程で生じる歪みによって、ちょうど $K(n,m)$ 個の\textbf{「みかけの特異点」として微分方程式上に突如として出現した}のである。

連立系からある1つの変数 $y_1$ についての単独方程式を導出する際、$y_1$ の1階微分、2階微分、$\dots$ と計算を進め、それらを独立な基底として元のベクトル $\vec{y}$ を逆算して消去する作業を行う。
この「逆算」は、代数的にはある変換行列 $P(x)$ の逆行列 $P(x)^{-1}$ を掛ける操作（あるいはロンスキアンのような量での割り算）に他ならない。この変換行列 $P(x)$ は元の連立系の係数（有理関数）から作られるため、$P(x)$ の行列式が $0$ になる点 $x=\lambda_j$ では割り算ができず、得られた単独方程式の係数に「新たな極」が発生してしまう。

しかし、この極は「方程式の表示を変えたこと」に起因するだけの人工的な特異点であるため、解のモノドロミー（分岐）には一切影響を与えない。すなわち、これこそが「みかけの特異点」の正体である。
一般の位置にある巡回ベクトルを選んで単独高階化を行った場合、この割り算によって生じるみかけの特異点の個数は、驚くべきことに次元の不足分 $K(n,m)$ 個と完全に一致する。
連立系における「係数行列の成分の自由度」が、単独系に変換された瞬間に「みかけの特異点の位置 $\lambda_1, \dots, \lambda_K$ という空間座標」に姿を変えて、モノドロミーの自由度を裏から支えているのである。
\end{Q}

以下、$n=2$（2階）の微分方程式系について詳しく見てみる。係数行列と未知関数を
\[
A(x)=
\begin{pmatrix}
    a_{11}(x) & a_{12}(x)\\
    a_{21}(x) & a_{22}(x)\\
\end{pmatrix},
\quad 
\vec{y}=
\begin{pmatrix}
    y_{1}\\
    y_{2}\\
\end{pmatrix}
\]
とすると、成分ごとの方程式は
\[
\begin{cases}
\frac{dy_{1}}{dx}=a_{11}y_{1}+a_{12}y_{2}\\
\frac{dy_{2}}{dx}=a_{21}y_{1}+a_{22}y_{2}
\end{cases}
\]
となる。第1式から $y_2$ を $y_1$ で表し（$y_2 = \frac{1}{a_{12}}(\frac{dy_{1}}{dx}-a_{11}y_{1})$）、これを $x$ で微分して第2式に代入することで、単独高階化を行う。
\begin{align*}
\frac{d^{2}y_{1}}{dx^{2}} &= a_{11}\frac{dy_{1}}{dx} + \frac{da_{11}}{dx}y_{1} + a_{12}\frac{dy_{2}}{dx} + \frac{da_{12}}{dx}y_{2}\\
&= a_{11}\frac{dy_{1}}{dx} + \frac{da_{11}}{dx}y_{1} + a_{12}(a_{21}y_{1}+a_{22}y_{2}) + \frac{da_{12}}{dx}\frac{1}{a_{12}}\left(\frac{dy_{1}}{dx}-a_{11}y_{1}\right)\\
&= a_{11}\frac{dy_{1}}{dx} + \frac{da_{11}}{dx}y_{1} + a_{12}a_{21}y_{1} + a_{22}\left(\frac{dy_{1}}{dx}-a_{11}y_{1}\right) + \frac{da_{12}}{dx}\frac{1}{a_{12}}\left(\frac{dy_{1}}{dx}-a_{11}y_{1}\right)\\
&= \left(a_{11} + a_{22} + \frac{1}{a_{12}}\frac{da_{12}}{dx}\right)\frac{dy_{1}}{dx} + \left(a_{12}a_{21} - a_{11}a_{22} + \frac{da_{11}}{dx} - \frac{a_{11}}{a_{12}}\frac{da_{12}}{dx}\right)y_{1}\\
&= \left(\text{Tr} A(x) + \frac{d}{dx}\log a_{12}\right)\frac{dy_{1}}{dx} + \left(-\det A(x) + \frac{da_{11}}{dx} - a_{11}\frac{d}{dx}\log a_{12}\right)y_{1}
\end{align*}
これを標準的な単独2階線形微分方程式の形 $\frac{d^{2}y_{1}}{dx^{2}} + p(x)\frac{dy_{1}}{dx} + q(x)y_{1} = 0$ に合わせるため、以下のように置く。
\begin{align*}
p(x) &= -\left(\text{Tr} A(x) + \frac{d}{dx}\log a_{12}\right)\\
q(x) &= \det A(x) - \frac{da_{11}}{dx} + a_{11}\frac{d}{dx}\log a_{12}
\end{align*}

ここで、特異点の数が $m=3$ （無限遠点を含めて4つ）であるシュレージンガー型方程式
\begin{equation}\label{eq:schlesinger_m3}
\frac{d\vec{y}}{dx} = A(x)\vec{y}, \quad A(x) = \frac{A_{0}}{x} + \frac{A_{1}}{x-1} + \frac{A_{t}}{x-t}
\end{equation}
を考察する。本来の確定特異点は一次分数変換の自由度を用いて $\Xi=\{0, 1, t, \infty\}$ と固定してよい。（係数行列 $A_0, A_1, A_t$ が零行列に退化しない限り、これらは特異点として残る。）
さらに、無限遠点における留数行列 $A_{\infty} = -(A_{0} + A_{1} + A_{t})$ は、事前のゲージ変換によって対角化され、
\[
A_{\infty} =
\begin{pmatrix}
s_{\infty} & 0\\
0 & s_{\infty}+\kappa_{\infty}
\end{pmatrix}
\]
となっているものと仮定する。

単独方程式の係数 $p(x), q(x)$ には対数微分 $\frac{d}{dx}\log a_{12} = \frac{a_{12}'}{a_{12}}$ が含まれているため、元の連立系の特異点 $\{0, 1, t, \infty\}$ に加えて、\textbf{「$a_{12}(x)$ の零点」が単独方程式の新たな特異点として出現する。}
$A(x)$ を通分すると $A(x) = \frac{\tilde{A}(x)}{x(x-1)(x-t)}$ となり、分子 $\tilde{A}(x)$ は
\begin{align*}
\tilde{A}(x) &= A_{0}(x-1)(x-t) + A_{1}x(x-t) + A_{t}x(x-1)\\
&= (A_{0}+A_{1}+A_{t})x^{2} - \big((t+1)A_{0}+tA_{1}+A_{t}\big)x + tA_{0}
\end{align*}
となる。(1,2)成分に注目すると、$A_{\infty}$ が対角行列であるため、$(A_{0}+A_{1}+A_{t})_{12} = -(A_{\infty})_{12} = 0$ となる。これにより $x^{2}$ の項が消滅し、分子は $x$ の1次式になる。定数 $X$ と $\lambda$ を用いて因数分解すると、
\[
a_{12}(x) = X\frac{x-\lambda}{x(x-1)(x-t)}
\]
となる。よって、$x=\lambda$ が新しく現れた特異点（$a_{12}$ の1位の零点）である。

\begin{Q}[【核心】新しく現れた特異点 $x=\lambda$ は、本当に「みかけの特異点」の厳しい条件（特性指数が整数差であり、かつ対数項を持たない）をクリアしているのか？]
\textbf{A.} 完全にクリアしている。

まず特性指数を調べる。$a_{12}(x)$ は $x=\lambda$ で1位の零点を持つため、対数微分 $\frac{d}{dx}\log a_{12}$ は $x=\lambda$ に留数 $1$ の極を持つ。したがって、
\begin{align*}
(x-\lambda)p(x)\big|_{x=\lambda} &= -1\\
(x-\lambda)^{2}q(x)\big|_{x=\lambda} &= 0 \quad (\text{$q(x)$ の極は高々1位であるため})
\end{align*}
となる。決定方程式（indicial equation）は $s(s-1) + s \cdot (-1) + 0 = s(s-2) = 0$ となり、特性指数は $s=0, 2$ となる。すべて整数であり、その差は2である。

次に非対数的であること（対数項の非存在）の証明だが、これは計算するまでもなく明らかである。元の連立方程式 \eqref{eq:schlesinger_m3} の特異点は $\{0, 1, t, \infty\}$ のみであり、$x=\lambda$ において係数行列 $A(x)$ は完全に正則である。微分方程式の基本定理より、正則点での解ベクトル $\vec{y}(x)$ は当然正則である。
単独高階化は単なる式変形であるため、得られた単独方程式の解 $y_1(x)$ は元の解の第1成分に他ならず、$x=\lambda$ でテイラー展開可能（正則）でなければならない。正則関数に対数項 $\log(x-\lambda)$ が入り込む余地はないため、$x=\lambda$ は決定方程式の特性指数の差にかかわらず、確実に非対数的特異点となる。
\end{Q}

\begin{proposition}
    シュレージンガー型方程式 \eqref{eq:schlesinger_m3}（$m=3, n=2$）の単独高階化は、本来の特異点 $\Xi=\{0,1,t,\infty\}$ の他に、ただ1つの点 $x=\lambda$ をみかけの特異点として持つ。またその特性指数は $0, 2$ である。
\end{proposition}

この結果をパラメータ（モジュライ空間の次元）の観点から振り返る。
\begin{align*}
E(2,3) &= \frac{1}{2} \cdot 2 \cdot (3(2+1)-(2-1)) = 8\\
M(2,3) &= (3-1) \cdot 2^{2} + 1 = 9\\
K(2,3) &= M(2,3) - E(2,3) = 9 - 8 = 1
\end{align*}
不足していたパラメータの数 $g = K(2,3) = 1$ が、まさに単独高階化によって生じた「みかけの特異点の位置 $\lambda$」の数 1 と完全に符号している。
%#endregion


\subsection{リーマンの問題}

%#region
これまで、フックス型微分方程式に対して、確定特異点の集合 $\Xi$ とそのモノドロミー表現 $\hat{\rho}$ をまとめた組 $(\mathbb{P}^{1}(\mathbb{C}), \Xi, \hat{\rho})$ をリーマンデータと呼んできた。ここでは、これとは「逆」の対応を問う以下の問題を考える。

\begin{itemize}
    \item[] \textbf{問題} \\
    任意の3つの組 $(\mathbb{P}^{1}(\mathbb{C}), \Xi, \hat{\rho})$ が与えられたとき、これをリーマンデータとして持つようなフックス型微分方程式は存在するだろうか？
\end{itemize}

\begin{definition}[リーマン・ヒルベルトの問題]
    上の「与えられたモノドロミー表現を実現する微分方程式の構成問題」を、\textbf{リーマンの問題} または \textbf{リーマン・ヒルベルトの問題 (Riemann-Hilbert problem)} という。ダフィット・ヒルベルトが1900年に提唱した「ヒルベルトの23の問題」における第21問題としても知られる。
\end{definition}

この問題を幾何学的に捉え直すために、パラメータ空間（モジュライ空間）を考える。
（特異点の位置 $\Xi$ を固定した）フックス型微分方程式の係数がなす空間を $\mathcal{E}$、モノドロミー表現の同値類がなす空間を $\mathcal{M}$ とする。
微分方程式が一つ与えられれば、そこから（基底の取り替えの自由度を除いて）モノドロミー表現が一意に定まるため、空間 $\mathcal{E}$ から空間 $\mathcal{M}$ への自然な写像 $\gamma$ を定義できる。

\begin{definition}[モノドロミー写像]
    この、微分方程式からモノドロミー表現を対応させる写像 $\gamma : \mathcal{E} \to \mathcal{M}$ を、\textbf{モノドロミー写像 (monodromy map)} という。
\end{definition}

\begin{Q}[【核心】リーマン・ヒルベルトの問題を、このモノドロミー写像 $\gamma$ の言葉で言い換えるとどうなるか？また、これまでのパラメータ数 $E(n,m)$ や $M(n,m)$ の議論は、この問題に対してどのような結論をもたらすか？]
\textbf{A.} リーマン・ヒルベルトの問題は、\textbf{「モノドロミー写像 $\gamma : \mathcal{E} \to \mathcal{M}$ は全射（onto）であるか？」}という写像の性質を問う問題に完全に翻訳される。

もし $\gamma$ が全射であれば、任意のモノドロミー表現（$\mathcal{M}$ の元）に対して、それを写像先とする微分方程式（$\mathcal{E}$ の元）が必ず存在することになる。
しかし、これまでの議論で見たように、両空間の次元（パラメータの数）を比較すると、一般に
\[
\dim \mathcal{E} = E(n,m) < M(n,m) = \dim \mathcal{M}
\]
となる。定義域 $\mathcal{E}$ の次元が値域 $\mathcal{M}$ の次元よりも真に小さいため、写像 $\gamma$ が全射になることはあり得ない。

つまり、「本来の確定特異点 $\Xi$ のみを持つ微分方程式の空間 $\mathcal{E}$ だけでは、リーマン・ヒルベルトの問題は一般には（否定的に）解けない」という強烈な事実が、単純な次元の比較からエレガントに証明されるのである。そして、この「次元の不足分」を補い、写像を全射にして問題を肯定的に解決するための代償として導入されたのが、「みかけの特異点」に他ならない。
\end{Q}

ここで、シュレージンガー方程式系の空間 $\mathcal{E}_{Sch}$ を考えれば、その次元は
\[
\dim \mathcal{E}_{Sch} = \dim \mathcal{M} = (m-1)n^{2}+1
\]
となり、モノドロミー表現のモジュライ空間の次元と完全に一致するため、リーマン・ヒルベルトの問題を考えるに足る十分なパラメータを持っているように見える。
しかし、シュレージンガー方程式系の空間からモノドロミー表現の空間へのモノドロミー写像 $\gamma : \mathcal{E}_{Sch} \to \mathcal{M}$ は、一般には全射（surjective）にならないことが知られている（ボリブルフの反例）。

この「次元が一致しているにもかかわらず、連立系では実現できないモノドロミーが存在する」という困難を解決するため、再び単独高階微分方程式の視点に立ち戻り、十分な数のみかけの特異点
\[
g \ge K(n,m) = M(n,m) - E(n,m)
\]
をあえて追加した微分方程式の空間 $\mathcal{E}_g$ を定義し、この空間上での構成問題を考える。

\begin{Q}[【核心】シュレージンガー型（連立系）で次元が一致しても失敗したのに、なぜ「みかけの特異点を追加した単独方程式」に望みを託すのか？]
\textbf{A.} みかけの特異点の位置 $\{\lambda_j\}$ が、連立系では到達できなかったモノドロミー表現の空間の「欠落した領域」を補完する新たな幾何学的自由度として機能するからである。

ボリブルフの反例は、特定のモノドロミー表現に対して、それを実現する「1位の極のみを持つ連立系（シュレージンガー型）」が存在しないことを示した。
しかし、高階単独微分方程式においてみかけの特異点 $g$ 個を許容すれば、方程式の自由度は $E(n,m) + g$ となり、少なくとも $g = K(n,m)$ と取れば次元の不足は解消される。単独高階化の議論で見た通り、連立系の「係数パラメータ」を単独系の「特異点の座標パラメータ」へ変換して扱うこの枠組みは、モノドロミー表現を網羅する（全射性を確保する）上で、連立系よりも柔軟な構造を提供のである。
\end{Q}

\begin{Q}[【核心】$g$ を $K(n,m)$ よりも大きく取る（$g > K$）ことには、どのような数学的メリットがあるのか？]
\textbf{A.} リーマン・ヒルベルト問題を「確実に」肯定的に解決するためである。
$g = K(n,m)$ はあくまで次元を一致させるための「必要最低限」の数であるが、モノドロミー写像がその次元において常に全射である保証はない。しかし、みかけの特異点の数 $g$ を十分に大きく取れば、モノドロミー表現を構成する際の代数的な制約をより多くの自由度で吸収できるようになり、任意のリーマンデータに対してそれを実現する微分方程式の存在を保証しやすくなる。
\end{Q}
%#endregion
%#region
微分方程式の具体的な形からは一旦離れ、純粋に幾何学的・代数的なデータのみを抽出する。
$\Xi=\{\xi_{1},\cdots,\xi_{m},\xi_{m+1}\}$ とし、$\mathbb{X}=\mathbb{P}^{1}(\mathbb{C})\setminus \Xi$ とする。基本群 $\pi_1(\mathbb{X})$ から一般線形群への準同型（表現）
\[
\rho:\pi_{1}(\mathbb{X})\to GL(n,\mathbb{C})
\]
が与えられたとき、適当な正則行列 $P$ による相似変換 $\rho'(\gamma) = P^{-1}\rho(\gamma)P$ で移り合う表現の同値類 $\hat{\rho}$ が定義できる。

\begin{definition}[リーマンデータの純粋化]
    微分方程式の存在とは独立に、幾何学的に定まる組 $(\mathbb{P}^{1}(\mathbb{C}),\Xi,\hat{\rho})$ のことを \textbf{リーマンデータ} と呼ぶ。
\end{definition}

さて、この表現 $\rho$ は、ベクトル空間 $\mathbb{C}^n$ に対する群 $\pi_1(\mathbb{X})$ の作用とみなすことができる。ここで、微分方程式の解の解析接続の性質から、作用の「向き」について極めて重要な注意が必要となる。

\begin{Q}[【核心】表現 $\rho$ による作用を考えるとき、$\mathbb{C}^n$ は $\pi_1(\mathbb{X})$ の「左加群」と「右加群」のどちらとして扱うのが自然か？]
\textbf{A.} 微分方程式のモノドロミーの自然な構造に従うならば、\textbf{「右 $\pi_1(\mathbb{X})$-加群」}として扱うのが最も理にかなっている。

この理由は、「経路の積の順序」と「行列の積の順序」の衝突にある。
基本群における経路の積 $\gamma_1 \circ \gamma_2$ を「$\gamma_1$ を辿った後に $\gamma_2$ を辿る」と定義する。
微分方程式の解の基底を並べた基本行列 $Y(x)$ を解析接続すると、モノドロミー行列 $M_\gamma$ は $Y(x)$ の「右から」掛かる。
したがって、$\gamma_1$ の後に $\gamma_2$ を辿った場合、基本行列の変換は
\[
Y \xrightarrow{\gamma_1} Y M_{\gamma_1} \xrightarrow{\gamma_2} (Y M_{\gamma_2}) M_{\gamma_1} = Y (M_{\gamma_2} M_{\gamma_1})
\]
となり、表現行列の積の順序が逆転し、$M_{\gamma_1 \circ \gamma_2} = M_{\gamma_2} M_{\gamma_1}$ となる（すなわち表現が反準同型になる）。

この順序のねじれを解消するためには、ベクトルを列ベクトルではなく「行ベクトル $\vec{v}^T$」とし、行列を右から作用させる
\[
\vec{v}^T \cdot \gamma := \vec{v}^T \rho(\gamma)
\]
という定義を採用するのが自然である。このとき、$(\vec{v}^T \cdot \gamma_1) \cdot \gamma_2 = (\vec{v}^T \rho(\gamma_1)) \rho(\gamma_2) = \vec{v}^T (\rho(\gamma_1)\rho(\gamma_2))$ となり、群の積と作用の順序が完璧に整合する。ゆえに、$\mathbb{C}^n$ は「非可換な群環 $\mathbb{C}[\pi_1(\mathbb{X})]$ 上の右加群（ベクトル空間の非可換版）」とみなすことができるのである。

（※注：文献によっては $\mathbb{C}^n$ を無理やり「左加群」として扱うものもあるが、その場合は「経路の積 $\gamma_1 \circ \gamma_2$ を関数の合成のように右から順に辿るように定義し直す」か、あるいは「表現 $\rho(\gamma)$ をモノドロミー行列の逆行列 $M_\gamma^{-1}$ で定義する」といった、規約（コンベンション）の意図的な書き換えが裏で行われていることに注意が必要である。）
\end{Q}

\begin{definition}[表現の既約性]
    この右 $\pi_{1}(\mathbb{X})$-加群 $\mathbb{C}^{n}$ において、任意の $\gamma \in \pi_{1}(\mathbb{X})$ の作用に対して不変となる（すなわち $W \rho(\gamma) \subseteq W$ となる）真の部分空間 $W \subsetneq \mathbb{C}^{n}$ が $\{0\}$ に限られるとき、表現 $\rho$ は \textbf{既約 (irreducible)} であるという。
    このとき、その同値類 $\hat{\rho}$ も \textbf{既約} であるという。既約でない表現（$\{0\}$ 以外の不変部分空間を持つ表現）を一般に \textbf{可約 (reducible)} という。
\end{definition}

\begin{Q}[【核心】表現 $\rho$ が可約であるとき、その具体的な行列表現はどうなるか？]
\textbf{A.} 適切な基底を取ることで、すべての $\rho(\gamma)$ が共通の\textbf{ブロック行列（行ベクトルへの右作用の場合はブロック下三角行列）}になる。

表現 $\rho$ が可約であるとする。定義より、$0$ ではない $n'$ 次元（$0 < n' < n$）の不変部分空間 $W$ が存在する。$W$ の基底を最初に取り、それにベクトルを付け加えて $\mathbb{C}^{n}$ 全体の基底とする。
行ベクトルに対する右作用 $\vec{v}^T \mapsto \vec{v}^T \rho(\gamma)$ において $W$ が閉じているためには、行列 $\rho(\gamma)$ は
\[
\rho(\gamma_{l})=
\begin{pmatrix}
A_{l} & 0\\
B_{l} & C_{l}
\end{pmatrix}
\]
というブロック下三角の形にならなければならない。（ここで、$A_{l}$ は $n'$ 次正方行列、$C_{l}$ は $(n-n')$ 次正方行列）。
このとき、左上のブロック $\rho'(\gamma_{l}):=A_{l}$ を取り出すと、これも準同型の性質を満たすため、$\pi_{1}(\mathbb{X})$ の $\mathbb{C}^{n'}$ 上の新たな表現（部分表現）となる。
\end{Q}

%#endregion
%#region
これまでの議論を踏まえ、リーマンデータ $(\mathbb{P}^{1}(\mathbb{C}),\Xi,\hat{\rho})$ が与えられているとき、ここに新たな特異点の集合 $\Lambda=\{\lambda_{1},\cdots,\lambda_{g}\}$ を追加した新しいリーマンデータ
\[
(\mathbb{P}^{1}(\mathbb{C}),\Xi\cup\Lambda,\hat{\rho'})
\]
について考えたい。
ここで、新しい表現類 $\hat{\rho'}$ は以下の条件を満たす代表元 $\rho'$ を持つとする。
元の確定特異点 $\xi_{1},\cdots,\xi_{m+1}$ を反時計回りに1周する道を $\gamma_{1},\cdots,\gamma_{m+1}$、新しく追加した特異点 $\lambda_{1},\cdots,\lambda_{g}$ を1周する道を $\gamma'_{1},\cdots,\gamma'_{g}$ とする。このとき、元の表現 $\rho \in \hat{\rho}$ が存在して、任意の $i, j$ に対して
\begin{align*}
\rho'(\gamma_{i}) &= \rho(\gamma_{i}) \\
\rho'(\gamma'_{j}) &= I_{n} \quad (\text{単位行列})
\end{align*}
を満たすものとする。

\begin{Q}[【核心】新しく追加した点に対する表現を $\rho'(\gamma'_{j}) = I_{n}$ と設定することは、代数的にどのような意味を持つか？また、これは元の表現と矛盾しないか？]
\textbf{A.} これは、解析的な「みかけの特異点」を、代数的な「表現の核（自明な作用）」として完璧に翻訳した設定であり、数学的に一切の矛盾を生じない。

解析において、みかけの特異点とは「その周りを回っても解が分岐しない（元の解にそのまま戻る）」点であった。表現論の言葉では、ベクトルに右から作用させたとき $\vec{v}^T \rho'(\gamma'_j) = \vec{v}^T$ となること、すなわち作用する行列が単位行列 $I_n$ であることと完全に同値である。

また、穴の開いた球面（$\mathbb{P}^{1}(\mathbb{C})\setminus (\Xi\cup\Lambda)$）の基本群は、すべての特異点を1周するループの積が自明になるという関係式
\[
\gamma_{1} \cdots \gamma_{m+1} \gamma'_{1} \cdots \gamma'_{g} = 1
\]
を持つ。これに新しい表現 $\rho'$ を適用すると、
\[
\rho'(\gamma_{1}) \cdots \rho'(\gamma_{m+1}) \rho'(\gamma'_{1}) \cdots \rho'(\gamma'_{g}) = \rho(\gamma_{1}) \cdots \rho(\gamma_{m+1}) I_{n} \cdots I_{n} = I_{n}
\]
となり、元の表現 $\rho$ が持っていた関係式 $\prod \rho(\gamma_i) = I_n$ に完全に帰着する。したがって、$\rho'(\gamma'_j) = I_n$ と置くことは、元のモノドロミー表現の代数的な構造を一切破壊することなく、空間の穴（特異点）だけを増やす極めて自然な操作なのである。
\end{Q}

\begin{definition}[リーマンデータの拡大]
    上記の条件を満たすような新しいデータ $(\mathbb{P}^{1}(\mathbb{C}),\Xi\cup\Lambda,\hat{\rho'})$ が存在するとき、これを元のリーマンデータ $(\mathbb{P}^{1}(\mathbb{C}),\Xi,\hat{\rho})$ の \textbf{拡大 (extension)} という。
\end{definition}

\begin{definition}[本質的なリーマンデータ]
    与えられたリーマンデータ $(\mathbb{P}^{1}(\mathbb{C}),\Xi,\hat{\rho})$ の特異点集合 $\Xi$ が「みかけの特異点（モノドロミーが $I_n$ となる点）」を一つも含まないとき、すなわち他のいかなるデータの真の拡大（$\Lambda \neq \emptyset$ による拡大）にもなっていないとき、これを \textbf{本質的なリーマンデータ (essential Riemann data)} という。
\end{definition}

今後のリーマン・ヒルベルト問題の構成においては、まず余分な贅肉を一切持たない「本質的なリーマンデータ」を出発点として固定し、そこから次元を補うために人為的に拡大を行う、というステップを踏むことになる。
%#endregion
%#region
そこで、リーマンの問題を以下のように考え直す。
\begin{itemize}
    \item []\textbf{問題}\\
    $(\mathbb{P}^{1}(\mathbb{C}),\Xi,\hat{\rho})$が本質的なリーマンデータだとする。このとき、フックス型微分方程式でその拡大$(\mathbb{P}^{1}(\mathbb{C}),\Xi\cup\Lambda,\hat{\rho'})$をリーマンデータに持つものはあるだろうか?
\end{itemize}
また、この問題に対して以下の事実が知られている。
\begin{proposition}
    リーマンの問題は常に解をもつ。
\end{proposition}
%#endregion

\subsection{フックスの問題}

\begin{definition}[モノドロミー保存変形]\label{mono_deform}
%#region
%FIXME(書き直す。)
%FIXME(確定特異点を理解する。) 
    \begin{equation}\label{eq:before}
    \frac{d^{n}y}{dx^{n}}+\sum_{j=1}^{n} p_{j}(x)\frac{d^{n-1}y}{dx^{n-1}} = 0   
    \end{equation}
    という以下の(P1)-(P3)を満たすフックス型微分方程式を考える。
    \begin{itemize}
        \item[(P1)] $x=\xi_{1},...,\xi_{n},\xi_{n+1}=\infty$を確定特異点、$x=\lambda_{1},...,\lambda_{g}$をみかけの特異点としてもつ。
        \item[(P2)] 各$x=\xi_{i}$における特性指数のどの2つも、差が整数ではない。
        \item[(P3)] モノドロミー$\hat{\rho}$は既約である。
    \end{itemize}
    $p_{j}(x)$が$t$に解析的に依存しているとする。
    \begin{equation}\label{eq:after}
    \frac{d^{n}y}{dx^{n}}+\sum_{j=1}^{n} p_{j}(x,t)\frac{d^{n-1}y}{dx^{n-1}} = 0   
    \end{equation}
    が以下の条件(H)を満たす時、\eqref{eq:after}を\eqref{eq:before}の\textbf{モノドロミー保存変形}または\textbf{ホロノミック変形}という。
    \begin{itemize}
        \item[\textbf{Type:}] \eqref{eq:before}は$\mathbb{P}^{1}(\mathbb{C})$上定義されている。$U\subset\mathbb{C}^{L}$は領域。$t=(t_{1},...,t_{L})\in U$
        \item[\textbf{Insight:}] 方程式のパラメータ $t$ を動かしても、基本群の表現（モノドロミー）が変化しないような変形。Schlesinger系やPainlevé方程式の導出における出発点となる。
        \item[\textbf{Link:}] 岡本2.8節の定義2.14/p.105を参照。
    \end{itemize}
%#endregion
\end{definition}